\documentclass[a4paper,10pt]{article}
\usepackage{amsmath,amssymb}
\usepackage{fontspec}
\setmainfont{Liberation Serif}
\usepackage{graphicx}
\usepackage{eso-pic}
\usepackage{tikz}
\usepackage[most]{tcolorbox}
\usepackage{geometry}
\usepackage{fancyhdr}
\usepackage{lastpage}
\usepackage{booktabs}
\usepackage{listings}
\usepackage{enumitem}
\usepackage{microtype}

\geometry{
    a4paper,
    top=1.5cm,
    bottom=1.8cm,
    left=1.2cm,
    right=1.2cm
}

% Persistent Unremovable Canvas Watermark on Every Page
\newsavebox{\watermarkbox}
\sbox{\watermarkbox}{%
    \begin{tabular}{c}
        \includegraphics[width=8cm]{/home/johan/GATE/GATE_CS_PYQ_DB/assets/studiva-butterfly-icon.jpg} \\[-0.5em]
        {\fontsize{18}{22}\selectfont\bfseries\scshape www.studiva.co.in}
    \end{tabular}%
}
\AddToShipoutPictureBG{%
    \begin{tikzpicture}[remember picture, overlay]
        \node[opacity=0.12] at (current page.center) {\usebox{\watermarkbox}};
    \end{tikzpicture}%
}

% Header and Footer
\pagestyle{fancy}
\fancyhf{}
\fancyhead[L]{\small\scshape GATE 2018--2025 CS/IT $\cdot$ \textbf{Easy Questions}}
\fancyhead[R]{\small\scshape Practice Question Bank with Solutions}
\fancyfoot[L]{\scriptsize\textsc{Copyright \copyright\ 2026} Studiva\textsuperscript{TM}}
\fancyfoot[C]{\small Page \thepage\ of \pageref{LastPage}}
\fancyfoot[R]{\scriptsize\texttt{www.studiva.co.in}}
\renewcommand{\headrulewidth}{0.4pt}
\renewcommand{\footrulewidth}{0.4pt}

% Code Listing Style (Monochrome)
\lstset{
    basicstyle=\footnotesize\ttfamily,
    numbers=none,
    frame=none,
    tabsize=2,
    breaklines=true,
    showstringspaces=false
}

\begin{document}

% ==================== PAGE 1: IMPRINT & TITLE PAGE ====================
\thispagestyle{empty}
\vspace*{0.5cm}
\begin{center}
    \includegraphics[width=2.5cm]{/home/johan/GATE/GATE_CS_PYQ_DB/assets/studiva-butterfly-icon.jpg}\\[0.6em]
    {\LARGE \textbf{GATE 2018--2025 COMPUTER SCIENCE \& IT}}\\[0.3em]
    {\Large \textbf{PREVIOUS YEAR QUESTION BANK WITH SOLUTIONS}}\\[0.4em]
    {\large \textsc{\textbf{Difficulty Wise: Easy Questions}}}\\[0.2em]
    {\normalsize \textsc{Combined 2018--2025 Examination Papers $\cdot$ Comprehensive Edition ($\mathbf{166}$ Questions)}}\\[1.2em]
    \rule{0.6\textwidth}{0.6pt}
\end{center}

\vspace{1em}
\noindent\begin{tabular}{ll}
\textbf{Author / Compiler:} & Rupam Ghosh \\
\textbf{Publisher \& Platform:} & Studiva\textsuperscript{TM} (\texttt{www.studiva.co.in}) \\
\textbf{Official Inquiries:} & \texttt{contact@studiva.co.in} \\
\textbf{Edition:} & First Standard Edition (2026) \\
\textbf{Coverage:} & Complete GATE CS/IT Papers (2018, 2022, 2023, 2024 CS1, 2024 CS2, 2025 CS1, 2025 CS2) \\
\textbf{Syllabus Reference:} & Official GATE 2026 Syllabus (Organizing Institute: IIT Guwahati) \\
\end{tabular}

\vspace{1.5em}
\begin{tcolorbox}[enhanced, standard jigsaw, opacityback=0, colframe=black!40, arc=1mm, boxrule=0.5pt, left=10pt, right=10pt, top=8pt, bottom=8pt]
\textbf{\large COPYRIGHT \& INTELLECTUAL PROPERTY NOTICE}\\[0.6em]
Copyright \copyright\ 2026 Rupam Ghosh \& Studiva\textsuperscript{TM}. All Rights Reserved.\\[0.5em]
This educational compilation, problem organization, taxonomy classification, LaTeX typesetting, hints, and explanatory solutions are the proprietary intellectual property of \textbf{Rupam Ghosh} and \textbf{Studiva\textsuperscript{TM}}.\\[0.6em]
\textbf{Terms of Use \& Attribution Rules:}
\begin{itemize}[leftmargin=*, nosep]
    \item \textbf{Personal Educational Use:} Permitted freely for individual students and educators preparing for competitive examinations.
    \item \textbf{Mandatory Attribution:} Any reproduction, republication, hosting, or sharing (digital or physical) must prominently credit \textbf{Rupam Ghosh \& Studiva\textsuperscript{TM}} and include an active, direct link to \texttt{https://www.studiva.co.in}.
    \item \textbf{Commercial Restrictions:} Commercial resale, sublicensing, course bundling, or inclusion in paid test series without prior written permission is strictly prohibited.
    \item \textbf{Anti-Tampering Notice:} Removing, altering, or obscuring copyright notices, page watermarks, or website links from this document is strictly prohibited.
\end{itemize}
\vspace{0.6em}
\textbf{Fair Use Disclaimer:}\\
The underlying examination questions originate from the Graduate Aptitude Test in Engineering (GATE), conducted by the Indian Institutes of Technology (IITs) and Indian Institute of Science (IISc). Question texts are reproduced for non-commercial educational study under Fair Use doctrines. No copyright is claimed over the original GATE examination question texts.
\end{tcolorbox}

\vspace{1.2em}
\begin{center}
\small\itshape
Practice Recommendation: Use an opaque sheet of paper or ruler to cover the right-hand column during mock attempts. Tick your chosen answers directly in the check boxes on the left, then slide the cover to check the solution and hint.
\end{center}

\newpage
% ==================== PAGE 2 ONWARDS: QUESTION BANK ====================

\vspace{8pt}
\subsection*{General Aptitude}
\vspace{-4pt}
\noindent
\begin{tcolorbox}[
    blanker,
    sidebyside,
    sidebyside align=top,
    lefthand width=0.55\textwidth,
    top=4pt, bottom=4pt, left=0pt, right=0pt,
    middle=16pt,
    underlay={
        \draw[dashed, black!80, line width=0.8pt, dash pattern=on 4pt off 3pt]
            ([xshift=0.55\textwidth+8pt]frame.north west) --
            ([xshift=0.55\textwidth+8pt]frame.south west);
    }
]
\textbf{Q.1} \quad {\footnotesize\textsc{[MCQ $\cdot$ 1 Mark(s) $\cdot$ Neg: 0.33]}} \hfill {\footnotesize\textbf{GATE 2018 CS} \quad \tiny\texttt{2018\_CS\_GA\_Q01}}
\vspace{2pt}\hrule height 0.3pt\vspace{4pt}

``From where are they bringing their books? \_\_\_\_\_\_\_\_ bringing \_\_\_\_\_\_\_ books from \_\_\_\_\_.''
\par\medskip
The words that best fill the blanks in the above sentence are

\par\vspace{6pt}\noindent
\begin{tabular}{@{}p{0.48\linewidth}@{\hspace{0.04\linewidth}}p{0.48\linewidth}@{}}
$\square$\ \textbf{(A)} Their, they're, there & $\square$\ \textbf{(B)} They're, their, there \\
\addlinespace[3pt]
$\square$\ \textbf{(C)} There, their, they're & $\square$\ \textbf{(D)} They're, there, there
\end{tabular}

\tcblower

\textbf{Official Key:} \textbf{\large B} \hfill \textit{\small [Easy]}
\vspace{2pt}\hrule height 0.3pt\vspace{3pt}
\textbf{Topic:} Verbal Aptitude \\
\textbf{Subtopic:} \textit{Vocabulary: words, idioms, and phrases in context}
\vspace{2pt}\hrule height 0.3pt\vspace{3pt}
\textbf{Guide / Hint:} \\
{\small "They're" is the contraction for "they are", "their" is the possessive pronoun, and "there" refers to a place. Thus, the complete sentence is: "They're bringing their books from there."}

\end{tcolorbox}
\vspace{6pt}\hrule height 0.35pt\vspace{6pt}
\noindent
\begin{tcolorbox}[
    blanker,
    sidebyside,
    sidebyside align=top,
    lefthand width=0.55\textwidth,
    top=4pt, bottom=4pt, left=0pt, right=0pt,
    middle=16pt,
    underlay={
        \draw[dashed, black!80, line width=0.8pt, dash pattern=on 4pt off 3pt]
            ([xshift=0.55\textwidth+8pt]frame.north west) --
            ([xshift=0.55\textwidth+8pt]frame.south west);
    }
]
\textbf{Q.2} \quad {\footnotesize\textsc{[MCQ $\cdot$ 1 Mark(s) $\cdot$ Neg: 0.33]}} \hfill {\footnotesize\textbf{GATE 2018 CS} \quad \tiny\texttt{2018\_CS\_GA\_Q02}}
\vspace{2pt}\hrule height 0.3pt\vspace{4pt}

``A \_\_\_\_\_\_\_\_\_ investigation can sometimes yield new facts, but typically organized ones are more successful.''
\par\medskip
The word that best fills the blank in the above sentence is

\par\vspace{6pt}\noindent
\begin{tabular}{@{}p{0.48\linewidth}@{\hspace{0.04\linewidth}}p{0.48\linewidth}@{}}
$\square$\ \textbf{(A)} meandering & $\square$\ \textbf{(B)} timely \\
\addlinespace[3pt]
$\square$\ \textbf{(C)} consistent & $\square$\ \textbf{(D)} systematic
\end{tabular}

\tcblower

\textbf{Official Key:} \textbf{\large A} \hfill \textit{\small [Easy]}
\vspace{2pt}\hrule height 0.3pt\vspace{3pt}
\textbf{Topic:} Verbal Aptitude \\
\textbf{Subtopic:} \textit{Vocabulary: words, idioms, and phrases in context}
\vspace{2pt}\hrule height 0.3pt\vspace{3pt}
\textbf{Guide / Hint:} \\
{\small The sentence contrasts an erratic or wandering investigation with an "organized" one. "Meandering" means proceeding in a winding or circuitous course without a clear direction, which fits the context precisely.}

\end{tcolorbox}
\vspace{6pt}\hrule height 0.35pt\vspace{6pt}
\noindent
\begin{tcolorbox}[
    blanker,
    sidebyside,
    sidebyside align=top,
    lefthand width=0.55\textwidth,
    top=4pt, bottom=4pt, left=0pt, right=0pt,
    middle=16pt,
    underlay={
        \draw[dashed, black!80, line width=0.8pt, dash pattern=on 4pt off 3pt]
            ([xshift=0.55\textwidth+8pt]frame.north west) --
            ([xshift=0.55\textwidth+8pt]frame.south west);
    }
]
\textbf{Q.3} \quad {\footnotesize\textsc{[MCQ $\cdot$ 1 Mark(s) $\cdot$ Neg: 0.33]}} \hfill {\footnotesize\textbf{GATE 2018 CS} \quad \tiny\texttt{2018\_CS\_GA\_Q03}}
\vspace{2pt}\hrule height 0.3pt\vspace{4pt}

The area of a square is $d$. What is the area of the circle which has the diagonal of the square as its diameter?

\par\vspace{6pt}\noindent
\begin{tabular}{@{}p{0.48\linewidth}@{\hspace{0.04\linewidth}}p{0.48\linewidth}@{}}
$\square$\ \textbf{(A)} $\pi d$ & $\square$\ \textbf{(B)} $\pi d^2$ \\
\addlinespace[3pt]
$\square$\ \textbf{(C)} $\frac{1}{4}\pi d^2$ & $\square$\ \textbf{(D)} $\frac{1}{2}\pi d$
\end{tabular}

\tcblower

\textbf{Official Key:} \textbf{\large D} \hfill \textit{\small [Easy]}
\vspace{2pt}\hrule height 0.3pt\vspace{3pt}
\textbf{Topic:} Quantitative Aptitude \\
\textbf{Subtopic:} \textit{Mensuration and geometry}
\vspace{2pt}\hrule height 0.3pt\vspace{3pt}
\textbf{Guide / Hint:} \\
{\small For a square of area $d$, the side length is $s = \sqrt{d}$ and its diagonal is $s\sqrt{2} = \sqrt{2d}$. Since this diagonal is the diameter of the circle, the radius is $r = \frac{\sqrt{2d}}{2} = \sqrt{\frac{d}{2}}$. Hence, the area of the circle is $\pi r^2 = \pi \left(\frac{d}{2}\right) = \frac{1}{2}\pi d$.}

\end{tcolorbox}
\vspace{6pt}\hrule height 0.35pt\vspace{6pt}
\noindent
\begin{tcolorbox}[
    blanker,
    sidebyside,
    sidebyside align=top,
    lefthand width=0.55\textwidth,
    top=4pt, bottom=4pt, left=0pt, right=0pt,
    middle=16pt,
    underlay={
        \draw[dashed, black!80, line width=0.8pt, dash pattern=on 4pt off 3pt]
            ([xshift=0.55\textwidth+8pt]frame.north west) --
            ([xshift=0.55\textwidth+8pt]frame.south west);
    }
]
\textbf{Q.4} \quad {\footnotesize\textsc{[MCQ $\cdot$ 1 Mark(s) $\cdot$ Neg: 0.33]}} \hfill {\footnotesize\textbf{GATE 2018 CS} \quad \tiny\texttt{2018\_CS\_GA\_Q04}}
\vspace{2pt}\hrule height 0.3pt\vspace{4pt}

What would be the smallest natural number which when divided either by 20 or by 42 or by 76 leaves a remainder of 7 in each case?

\par\vspace{6pt}\noindent
\begin{tabular}{@{}p{0.48\linewidth}@{\hspace{0.04\linewidth}}p{0.48\linewidth}@{}}
$\square$\ \textbf{(A)} 3047 & $\square$\ \textbf{(B)} 6047 \\
\addlinespace[3pt]
$\square$\ \textbf{(C)} 7987 & $\square$\ \textbf{(D)} 63847
\end{tabular}

\tcblower

\textbf{Official Key:} \textbf{\large C} \hfill \textit{\small [Easy]}
\vspace{2pt}\hrule height 0.3pt\vspace{3pt}
\textbf{Topic:} Quantitative Aptitude \\
\textbf{Subtopic:} \textit{Numerical computation and estimation: ratios, percentages, powers, exponents and logarithms, permutations and combinations, series}
\vspace{2pt}\hrule height 0.3pt\vspace{3pt}
\textbf{Guide / Hint:} \\
{\small The required number is $\text{LCM}(20, 42, 76) + 7$. Calculating prime factorizations: $20 = 2^2 \times 5$, $42 = 2 \times 3 \times 7$, $76 = 2^2 \times 19$. Thus $\text{LCM} = 2^2 \times 3 \times 5 \times 7 \times 19 = 7980$. Adding the remainder 7 gives $7980 + 7 = 7987$.}

\end{tcolorbox}
\vspace{6pt}\hrule height 0.35pt\vspace{6pt}
\noindent
\begin{tcolorbox}[
    blanker,
    sidebyside,
    sidebyside align=top,
    lefthand width=0.55\textwidth,
    top=4pt, bottom=4pt, left=0pt, right=0pt,
    middle=16pt,
    underlay={
        \draw[dashed, black!80, line width=0.8pt, dash pattern=on 4pt off 3pt]
            ([xshift=0.55\textwidth+8pt]frame.north west) --
            ([xshift=0.55\textwidth+8pt]frame.south west);
    }
]
\textbf{Q.5} \quad {\footnotesize\textsc{[MCQ $\cdot$ 1 Mark(s) $\cdot$ Neg: 0.33]}} \hfill {\footnotesize\textbf{GATE 2018 CS} \quad \tiny\texttt{2018\_CS\_GA\_Q05}}
\vspace{2pt}\hrule height 0.3pt\vspace{4pt}

What is the missing number in the following sequence?
\[ 2,\; 12,\; 60,\; 240,\; 720,\; 1440,\; \_\_\_\_\_,\; 0 \]

\par\vspace{6pt}\noindent
\begin{tabular}{@{}p{0.48\linewidth}@{\hspace{0.04\linewidth}}p{0.48\linewidth}@{}}
$\square$\ \textbf{(A)} 2880 & $\square$\ \textbf{(B)} 1440 \\
\addlinespace[3pt]
$\square$\ \textbf{(C)} 720 & $\square$\ \textbf{(D)} 0
\end{tabular}

\tcblower

\textbf{Official Key:} \textbf{\large B} \hfill \textit{\small [Easy]}
\vspace{2pt}\hrule height 0.3pt\vspace{3pt}
\textbf{Topic:} Quantitative Aptitude \\
\textbf{Subtopic:} \textit{Numerical computation and estimation: ratios, percentages, powers, exponents and logarithms, permutations and combinations, series}
\vspace{2pt}\hrule height 0.3pt\vspace{3pt}
\textbf{Guide / Hint:} \\
{\small Examining consecutive multipliers: $2 \times 6 = 12$, $12 \times 5 = 60$, $60 \times 4 = 240$, $240 \times 3 = 720$, $720 \times 2 = 1440$. Following this decreasing multiplier pattern, the next term is $1440 \times 1 = 1440$, and the term after is $1440 \times 0 = 0$.}

\end{tcolorbox}
\vspace{6pt}\hrule height 0.35pt\vspace{6pt}
\noindent
\begin{tcolorbox}[
    blanker,
    sidebyside,
    sidebyside align=top,
    lefthand width=0.55\textwidth,
    top=4pt, bottom=4pt, left=0pt, right=0pt,
    middle=16pt,
    underlay={
        \draw[dashed, black!80, line width=0.8pt, dash pattern=on 4pt off 3pt]
            ([xshift=0.55\textwidth+8pt]frame.north west) --
            ([xshift=0.55\textwidth+8pt]frame.south west);
    }
]
\textbf{Q.6} \quad {\footnotesize\textsc{[MCQ $\cdot$ 2 Mark(s) $\cdot$ Neg: 0.66]}} \hfill {\footnotesize\textbf{GATE 2018 CS} \quad \tiny\texttt{2018\_CS\_GA\_Q07}}
\vspace{2pt}\hrule height 0.3pt\vspace{4pt}

If $p^{-x} = \frac{1}{q}$, $q^{-y} = \frac{1}{r}$, and $r^{-z} = \frac{1}{p}$, what is the value of $xyz$?

\par\vspace{6pt}\noindent
\begin{tabular}{@{}p{0.48\linewidth}@{\hspace{0.04\linewidth}}p{0.48\linewidth}@{}}
$\square$\ \textbf{(A)} $-1$ & $\square$\ \textbf{(B)} 0 \\
\addlinespace[3pt]
$\square$\ \textbf{(C)} 1 & $\square$\ \textbf{(D)} $pqr$
\end{tabular}

\tcblower

\textbf{Official Key:} \textbf{\large C} \hfill \textit{\small [Easy]}
\vspace{2pt}\hrule height 0.3pt\vspace{3pt}
\textbf{Topic:} Quantitative Aptitude \\
\textbf{Subtopic:} \textit{Numerical computation and estimation: ratios, percentages, powers, exponents and logarithms, permutations and combinations, series}
\vspace{2pt}\hrule height 0.3pt\vspace{3pt}
\textbf{Guide / Hint:} \\
{\small From $p^{-x} = \frac{1}{q}$, we have $q = p^x$. Substituting into $q^{-y} = \frac{1}{r}$ gives $r = q^y = (p^x)^y = p^{xy}$. Then $r^{-z} = \frac{1}{p}$ gives $p = r^z = (p^{xy})^z = p^{xyz}$. Since $p > 0$ and $p \neq 1$, it follows that $xyz = 1$.}

\end{tcolorbox}
\vspace{6pt}\hrule height 0.35pt\vspace{6pt}
\noindent
\begin{tcolorbox}[
    blanker,
    sidebyside,
    sidebyside align=top,
    lefthand width=0.55\textwidth,
    top=4pt, bottom=4pt, left=0pt, right=0pt,
    middle=16pt,
    underlay={
        \draw[dashed, black!80, line width=0.8pt, dash pattern=on 4pt off 3pt]
            ([xshift=0.55\textwidth+8pt]frame.north west) --
            ([xshift=0.55\textwidth+8pt]frame.south west);
    }
]
\textbf{Q.7} \quad {\footnotesize\textsc{[MCQ $\cdot$ 2 Mark(s) $\cdot$ Neg: 0.66]}} \hfill {\footnotesize\textbf{GATE 2018 CS} \quad \tiny\texttt{2018\_CS\_GA\_Q08}}
\vspace{2pt}\hrule height 0.3pt\vspace{4pt}

In a party, 60\% of the invited guests are male and 40\% are female. 80\% of the invited guests attend the party. If all the invited female guests attend the party, what is the ratio of male guests who attended to the female guests who attended?

\par\vspace{6pt}\noindent
\begin{tabular}{@{}p{0.48\linewidth}@{\hspace{0.04\linewidth}}p{0.48\linewidth}@{}}
$\square$\ \textbf{(A)} 2 : 3 & $\square$\ \textbf{(B)} 1 : 1 \\
\addlinespace[3pt]
$\square$\ \textbf{(C)} 3 : 2 & $\square$\ \textbf{(D)} 2 : 1
\end{tabular}

\tcblower

\textbf{Official Key:} \textbf{\large B} \hfill \textit{\small [Easy]}
\vspace{2pt}\hrule height 0.3pt\vspace{3pt}
\textbf{Topic:} Quantitative Aptitude \\
\textbf{Subtopic:} \textit{Numerical computation and estimation: ratios, percentages, powers, exponents and logarithms, permutations and combinations, series}
\vspace{2pt}\hrule height 0.3pt\vspace{3pt}
\textbf{Guide / Hint:} \\
{\small Out of 100 invited guests, 60 are male and 40 are female. Since 80 guests attend and all 40 invited females attend, the remaining $80 - 40 = 40$ attendees must be male. Thus, the ratio of attending males to attending females is $40 : 40 = 1 : 1$.}

\end{tcolorbox}
\vspace{6pt}\hrule height 0.35pt\vspace{6pt}
\noindent
\begin{tcolorbox}[
    blanker,
    sidebyside,
    sidebyside align=top,
    lefthand width=0.55\textwidth,
    top=4pt, bottom=4pt, left=0pt, right=0pt,
    middle=16pt,
    underlay={
        \draw[dashed, black!80, line width=0.8pt, dash pattern=on 4pt off 3pt]
            ([xshift=0.55\textwidth+8pt]frame.north west) --
            ([xshift=0.55\textwidth+8pt]frame.south west);
    }
]
\textbf{Q.8} \quad {\footnotesize\textsc{[MCQ $\cdot$ 1 Mark(s) $\cdot$ Neg: 0.33]}} \hfill {\footnotesize\textbf{GATE 2022 CS} \quad \tiny\texttt{2022\_CS\_GA\_Q01}}
\vspace{2pt}\hrule height 0.3pt\vspace{4pt}

The \_\_\_\_\_\_\_\_\_ is too high for it to be considered \_\_\_\_\_\_\_\_\_.

\par\vspace{6pt}\noindent
\begin{tabular}{@{}p{0.48\linewidth}@{\hspace{0.04\linewidth}}p{0.48\linewidth}@{}}
$\square$\ \textbf{(A)} fair / fare & $\square$\ \textbf{(B)} faer / fair \\
\addlinespace[3pt]
$\square$\ \textbf{(C)} fare / fare & $\square$\ \textbf{(D)} fare / fair
\end{tabular}

\tcblower

\textbf{Official Key:} \textbf{\large D} \hfill \textit{\small [Easy]}
\vspace{2pt}\hrule height 0.3pt\vspace{3pt}
\textbf{Topic:} Verbal Aptitude \\
\textbf{Subtopic:} \textit{Vocabulary: words, idioms, and phrases in context}
\vspace{2pt}\hrule height 0.3pt\vspace{3pt}
\textbf{Guide / Hint:} \\
{\small 'Fare' denotes the price charged for a journey or conveyance, whereas 'fair' means reasonable or just. Therefore, 'The fare is too high for it to be considered fair.'}

\end{tcolorbox}
\vspace{6pt}\hrule height 0.35pt\vspace{6pt}
\noindent
\begin{tcolorbox}[
    blanker,
    sidebyside,
    sidebyside align=top,
    lefthand width=0.55\textwidth,
    top=4pt, bottom=4pt, left=0pt, right=0pt,
    middle=16pt,
    underlay={
        \draw[dashed, black!80, line width=0.8pt, dash pattern=on 4pt off 3pt]
            ([xshift=0.55\textwidth+8pt]frame.north west) --
            ([xshift=0.55\textwidth+8pt]frame.south west);
    }
]
\textbf{Q.9} \quad {\footnotesize\textsc{[MCQ $\cdot$ 1 Mark(s) $\cdot$ Neg: 0.33]}} \hfill {\footnotesize\textbf{GATE 2022 CS} \quad \tiny\texttt{2022\_CS\_GA\_Q04}}
\vspace{2pt}\hrule height 0.3pt\vspace{4pt}

Given below are four statements:
\begin{itemize}[itemsep=0pt,topsep=2pt]
\item Statement 1: All students are inquisitive.
\item Statement 2: Some students are inquisitive.
\item Statement 3: No student is inquisitive.
\item Statement 4: Some students are not inquisitive.
\end{itemize}
From the given four statements, find the two statements that CANNOT BE TRUE simultaneously, assuming that there is at least one student in the class.

\par\vspace{6pt}\noindent
\begin{tabular}{@{}p{0.48\linewidth}@{\hspace{0.04\linewidth}}p{0.48\linewidth}@{}}
$\square$\ \textbf{(A)} Statement 1 and Statement 3 & $\square$\ \textbf{(B)} Statement 1 and Statement 2 \\
\addlinespace[3pt]
$\square$\ \textbf{(C)} Statement 2 and Statement 4 & $\square$\ \textbf{(D)} Statement 3 and Statement 4
\end{tabular}

\tcblower

\textbf{Official Key:} \textbf{\large A} \hfill \textit{\small [Easy]}
\vspace{2pt}\hrule height 0.3pt\vspace{3pt}
\textbf{Topic:} Analytical Aptitude \\
\textbf{Subtopic:} \textit{Logic: deduction and induction}
\vspace{2pt}\hrule height 0.3pt\vspace{3pt}
\textbf{Guide / Hint:} \\
{\small Statement 1 ("All students are inquisitive") and Statement 3 ("No student is inquisitive") are universal affirmative and universal negative statements (contraries). In classical categorical logic with a non-empty class, both statements cannot be simultaneously true.}

\end{tcolorbox}
\vspace{6pt}\hrule height 0.35pt\vspace{6pt}
\noindent
\begin{tcolorbox}[
    blanker,
    sidebyside,
    sidebyside align=top,
    lefthand width=0.55\textwidth,
    top=4pt, bottom=4pt, left=0pt, right=0pt,
    middle=16pt,
    underlay={
        \draw[dashed, black!80, line width=0.8pt, dash pattern=on 4pt off 3pt]
            ([xshift=0.55\textwidth+8pt]frame.north west) --
            ([xshift=0.55\textwidth+8pt]frame.south west);
    }
]
\textbf{Q.10} \quad {\footnotesize\textsc{[MCQ $\cdot$ 2 Mark(s) $\cdot$ Neg: 0.66]}} \hfill {\footnotesize\textbf{GATE 2022 CS} \quad \tiny\texttt{2022\_CS\_GA\_Q06}}
\vspace{2pt}\hrule height 0.3pt\vspace{4pt}

Some people believe that ``what gets measured, improves''. Some others believe that ``what gets measured, gets gamed''. One possible reason for the difference in the beliefs is the work culture in organizations. In organizations with good work culture, metrics help improve outcomes. However, the same metrics are counterproductive in organizations with poor work culture.
Which one of the following is the CORRECT logical inference based on the information in the above passage?

\par\vspace{6pt}\noindent
\begin{tabular}{@{}p{0.98\linewidth}@{}}
$\square$\ \textbf{(A)} Metrics are useful in organizations with poor work culture \\
\addlinespace[3pt]
$\square$\ \textbf{(B)} Metrics are useful in organizations with good work culture \\
\addlinespace[3pt]
$\square$\ \textbf{(C)} Metrics are always counterproductive in organizations with good work culture \\
\addlinespace[3pt]
$\square$\ \textbf{(D)} Metrics are never useful in organizations with good work culture
\end{tabular}

\tcblower

\textbf{Official Key:} \textbf{\large B} \hfill \textit{\small [Easy]}
\vspace{2pt}\hrule height 0.3pt\vspace{3pt}
\textbf{Topic:} Verbal Aptitude \\
\textbf{Subtopic:} \textit{Reading and comprehension}
\vspace{2pt}\hrule height 0.3pt\vspace{3pt}
\textbf{Guide / Hint:} \\
{\small The passage explicitly notes that "In organizations with good work culture, metrics help improve outcomes." This directly implies that metrics are useful in organizations with good work culture.}

\end{tcolorbox}
\vspace{6pt}\hrule height 0.35pt\vspace{6pt}
\noindent
\begin{tcolorbox}[
    blanker,
    sidebyside,
    sidebyside align=top,
    lefthand width=0.55\textwidth,
    top=4pt, bottom=4pt, left=0pt, right=0pt,
    middle=16pt,
    underlay={
        \draw[dashed, black!80, line width=0.8pt, dash pattern=on 4pt off 3pt]
            ([xshift=0.55\textwidth+8pt]frame.north west) --
            ([xshift=0.55\textwidth+8pt]frame.south west);
    }
]
\textbf{Q.11} \quad {\footnotesize\textsc{[MCQ $\cdot$ 2 Mark(s) $\cdot$ Neg: 0.66]}} \hfill {\footnotesize\textbf{GATE 2022 CS} \quad \tiny\texttt{2022\_CS\_GA\_Q07}}
\vspace{2pt}\hrule height 0.3pt\vspace{4pt}

In a recently conducted national entrance test, boys constituted 65\% of those who appeared for the test. Girls constituted the remaining candidates and they accounted for 60\% of the qualified candidates.
Which one of the following is the correct logical inference based on the information provided in the above passage?

\par\vspace{6pt}\noindent
\begin{tabular}{@{}p{0.98\linewidth}@{}}
$\square$\ \textbf{(A)} Equal number of boys and girls qualified \\
\addlinespace[3pt]
$\square$\ \textbf{(B)} Equal number of boys and girls appeared for the test \\
\addlinespace[3pt]
$\square$\ \textbf{(C)} The number of boys who appeared for the test is less than the number of girls who appeared \\
\addlinespace[3pt]
$\square$\ \textbf{(D)} The number of boys who qualified the test is less than the number of girls who qualified
\end{tabular}

\tcblower

\textbf{Official Key:} \textbf{\large D} \hfill \textit{\small [Easy]}
\vspace{2pt}\hrule height 0.3pt\vspace{3pt}
\textbf{Topic:} Quantitative Aptitude \\
\textbf{Subtopic:} \textit{Numerical computation and estimation: ratios, percentages, powers, exponents and logarithms, permutations and combinations, series}
\vspace{2pt}\hrule height 0.3pt\vspace{3pt}
\textbf{Guide / Hint:} \\
{\small Among all candidates who qualified, girls accounted for 60\%, meaning boys accounted for the remaining $100\% - 60\% = 40\%$. Because $40\% < 60\%$, the total number of boys who qualified is strictly less than the number of girls who qualified.}

\end{tcolorbox}
\vspace{6pt}\hrule height 0.35pt\vspace{6pt}
\noindent
\begin{tcolorbox}[
    blanker,
    sidebyside,
    sidebyside align=top,
    lefthand width=0.55\textwidth,
    top=4pt, bottom=4pt, left=0pt, right=0pt,
    middle=16pt,
    underlay={
        \draw[dashed, black!80, line width=0.8pt, dash pattern=on 4pt off 3pt]
            ([xshift=0.55\textwidth+8pt]frame.north west) --
            ([xshift=0.55\textwidth+8pt]frame.south west);
    }
]
\textbf{Q.12} \quad {\footnotesize\textsc{[MCQ $\cdot$ 1 Mark(s) $\cdot$ Neg: 0.33]}} \hfill {\footnotesize\textbf{GATE 2023 CS} \quad \tiny\texttt{2023\_CS\_GA\_Q01}}
\vspace{2pt}\hrule height 0.3pt\vspace{4pt}

We reached the station late, and \_\_\_\_\_\_\_ missed the train.

\par\vspace{6pt}\noindent
\begin{tabular}{@{}p{0.48\linewidth}@{\hspace{0.04\linewidth}}p{0.48\linewidth}@{}}
$\square$\ \textbf{(A)} near & $\square$\ \textbf{(B)} nearly \\
\addlinespace[3pt]
$\square$\ \textbf{(C)} utterly & $\square$\ \textbf{(D)} mostly
\end{tabular}

\tcblower

\textbf{Official Key:} \textbf{\large B} \hfill \textit{\small [Easy]}
\vspace{2pt}\hrule height 0.3pt\vspace{3pt}
\textbf{Topic:} Verbal Aptitude \\
\textbf{Subtopic:} \textit{Basic English grammar: tenses, articles, adjectives, prepositions, conjunctions, verb-noun agreement}
\vspace{2pt}\hrule height 0.3pt\vspace{3pt}
\textbf{Guide / Hint:} \\
{\small The adverb 'nearly' means 'almost' and correctly modifies the past-tense verb 'missed'. 'Near' is an adjective/preposition, while 'utterly' and 'mostly' do not fit the intended semantic meaning of narrowly missing a train.}

\end{tcolorbox}
\vspace{6pt}\hrule height 0.35pt\vspace{6pt}
\noindent
\begin{tcolorbox}[
    blanker,
    sidebyside,
    sidebyside align=top,
    lefthand width=0.55\textwidth,
    top=4pt, bottom=4pt, left=0pt, right=0pt,
    middle=16pt,
    underlay={
        \draw[dashed, black!80, line width=0.8pt, dash pattern=on 4pt off 3pt]
            ([xshift=0.55\textwidth+8pt]frame.north west) --
            ([xshift=0.55\textwidth+8pt]frame.south west);
    }
]
\textbf{Q.13} \quad {\footnotesize\textsc{[MCQ $\cdot$ 1 Mark(s) $\cdot$ Neg: 0.33]}} \hfill {\footnotesize\textbf{GATE 2023 CS} \quad \tiny\texttt{2023\_CS\_GA\_Q02}}
\vspace{2pt}\hrule height 0.3pt\vspace{4pt}

Kind : \_\_\_\_\_\_\_ : : Often : Frequently
(By word meaning)

\par\vspace{6pt}\noindent
\begin{tabular}{@{}p{0.48\linewidth}@{\hspace{0.04\linewidth}}p{0.48\linewidth}@{}}
$\square$\ \textbf{(A)} Mean & $\square$\ \textbf{(B)} Type \\
\addlinespace[3pt]
$\square$\ \textbf{(C)} Cruel & $\square$\ \textbf{(D)} Kindly
\end{tabular}

\tcblower

\textbf{Official Key:} \textbf{\large MTA} \hfill \textit{\small [Easy]}
\vspace{2pt}\hrule height 0.3pt\vspace{3pt}
\textbf{Topic:} Verbal Aptitude \\
\textbf{Subtopic:} \textit{Vocabulary: words, idioms, and phrases in context}
\vspace{2pt}\hrule height 0.3pt\vspace{3pt}
\textbf{Guide / Hint:} \\
{\small 'Often' and 'Frequently' are direct synonyms. 'Kind' can be a noun meaning 'Type' (category) or an adjective meaning benevolent/caring. Because of multiple ambiguous interpretations of 'Kind', GATE awarded Marks to All (MTA) for this question.}

\end{tcolorbox}
\vspace{6pt}\hrule height 0.35pt\vspace{6pt}
\noindent
\begin{tcolorbox}[
    blanker,
    sidebyside,
    sidebyside align=top,
    lefthand width=0.55\textwidth,
    top=4pt, bottom=4pt, left=0pt, right=0pt,
    middle=16pt,
    underlay={
        \draw[dashed, black!80, line width=0.8pt, dash pattern=on 4pt off 3pt]
            ([xshift=0.55\textwidth+8pt]frame.north west) --
            ([xshift=0.55\textwidth+8pt]frame.south west);
    }
]
\textbf{Q.14} \quad {\footnotesize\textsc{[MCQ $\cdot$ 1 Mark(s) $\cdot$ Neg: 0.33]}} \hfill {\footnotesize\textbf{GATE 2023 CS} \quad \tiny\texttt{2023\_CS\_GA\_Q03}}
\vspace{2pt}\hrule height 0.3pt\vspace{4pt}

A series of natural numbers $F_1, F_2, F_3, F_4, F_5, F_6, F_7, \dots$ obeys $F_{n+1} = F_n + F_{n-1}$ for all integers $n \ge 2$.
If $F_6 = 37$, and $F_7 = 60$, then what is $F_1$?

\par\vspace{6pt}\noindent
\begin{tabular}{@{}p{0.48\linewidth}@{\hspace{0.04\linewidth}}p{0.48\linewidth}@{}}
$\square$\ \textbf{(A)} 4 & $\square$\ \textbf{(B)} 5 \\
\addlinespace[3pt]
$\square$\ \textbf{(C)} 8 & $\square$\ \textbf{(D)} 9
\end{tabular}

\tcblower

\textbf{Official Key:} \textbf{\large A} \hfill \textit{\small [Easy]}
\vspace{2pt}\hrule height 0.3pt\vspace{3pt}
\textbf{Topic:} Quantitative Aptitude \\
\textbf{Subtopic:} \textit{Numerical computation and estimation: ratios, percentages, powers, exponents and logarithms, permutations and combinations, series}
\vspace{2pt}\hrule height 0.3pt\vspace{3pt}
\textbf{Guide / Hint:} \\
{\small The recurrence relation is $F_{n+1} = F_n + F_{n-1} \implies F_{n-1} = F_{n+1} - F_n$. Working backward from $F_7 = 60$ and $F_6 = 37$: $F_5 = 60 - 37 = 23$; $F_4 = 37 - 23 = 14$; $F_3 = 23 - 14 = 9$; $F_2 = 14 - 9 = 5$; and $F_1 = 9 - 5 = 4$.}

\end{tcolorbox}
\vspace{6pt}\hrule height 0.35pt\vspace{6pt}
\noindent
\begin{tcolorbox}[
    blanker,
    sidebyside,
    sidebyside align=top,
    lefthand width=0.55\textwidth,
    top=4pt, bottom=4pt, left=0pt, right=0pt,
    middle=16pt,
    underlay={
        \draw[dashed, black!80, line width=0.8pt, dash pattern=on 4pt off 3pt]
            ([xshift=0.55\textwidth+8pt]frame.north west) --
            ([xshift=0.55\textwidth+8pt]frame.south west);
    }
]
\textbf{Q.15} \quad {\footnotesize\textsc{[MCQ $\cdot$ 1 Mark(s) $\cdot$ Neg: 0.33]}} \hfill {\footnotesize\textbf{GATE 2023 CS} \quad \tiny\texttt{2023\_CS\_GA\_Q04}}
\vspace{2pt}\hrule height 0.3pt\vspace{4pt}

A survey for a certain year found that 90\% of pregnant women received medical care at least once before giving birth. Of these women, 60\% received medical care from doctors, while 40\% received medical care from other healthcare providers.
Given this information, which one of the following statements can be inferred with certainty?

\par\vspace{6pt}\noindent
\begin{tabular}{@{}p{0.98\linewidth}@{}}
$\square$\ \textbf{(A)} More than half of the pregnant women received medical care at least once from a doctor. \\
\addlinespace[3pt]
$\square$\ \textbf{(B)} Less than half of the pregnant women received medical care at least once from a doctor. \\
\addlinespace[3pt]
$\square$\ \textbf{(C)} More than half of the pregnant women received medical care at most once from a doctor. \\
\addlinespace[3pt]
$\square$\ \textbf{(D)} Less than half of the pregnant women received medical care at most once from a doctor.
\end{tabular}

\tcblower

\textbf{Official Key:} \textbf{\large A} \hfill \textit{\small [Easy]}
\vspace{2pt}\hrule height 0.3pt\vspace{3pt}
\textbf{Topic:} Quantitative Aptitude \\
\textbf{Subtopic:} \textit{Numerical computation and estimation: ratios, percentages, powers, exponents and logarithms, permutations and combinations, series}
\vspace{2pt}\hrule height 0.3pt\vspace{3pt}
\textbf{Guide / Hint:} \\
{\small Out of all pregnant women, 90\% received medical care at least once. Among those who received medical care, 60\% received care from doctors. The percentage of all pregnant women who received care at least once from a doctor is $90\% \times 60\% = 54\%$, which is strictly greater than half ($> 50\%$).}

\end{tcolorbox}
\vspace{6pt}\hrule height 0.35pt\vspace{6pt}
\noindent
\begin{tcolorbox}[
    blanker,
    sidebyside,
    sidebyside align=top,
    lefthand width=0.55\textwidth,
    top=4pt, bottom=4pt, left=0pt, right=0pt,
    middle=16pt,
    underlay={
        \draw[dashed, black!80, line width=0.8pt, dash pattern=on 4pt off 3pt]
            ([xshift=0.55\textwidth+8pt]frame.north west) --
            ([xshift=0.55\textwidth+8pt]frame.south west);
    }
]
\textbf{Q.16} \quad {\footnotesize\textsc{[MCQ $\cdot$ 1 Mark(s) $\cdot$ Neg: 0.33]}} \hfill {\footnotesize\textbf{GATE 2023 CS} \quad \tiny\texttt{2023\_CS\_GA\_Q05}}
\vspace{2pt}\hrule height 0.3pt\vspace{4pt}

Looking at the surface of a smooth 3-dimensional object from the outside, which one of the following options is TRUE?

\par\vspace{6pt}\noindent
\begin{tabular}{@{}p{0.98\linewidth}@{}}
$\square$\ \textbf{(A)} The surface of the object must be concave everywhere. \\
\addlinespace[3pt]
$\square$\ \textbf{(B)} The surface of the object must be convex everywhere. \\
\addlinespace[3pt]
$\square$\ \textbf{(C)} The surface of the object may be concave in some places and convex in other places. \\
\addlinespace[3pt]
$\square$\ \textbf{(D)} The object can have edges, but no corners.
\end{tabular}

\tcblower

\textbf{Official Key:} \textbf{\large C} \hfill \textit{\small [Easy]}
\vspace{2pt}\hrule height 0.3pt\vspace{3pt}
\textbf{Topic:} Spatial Aptitude \\
\textbf{Subtopic:} \textit{Transformation of shapes: translation, rotation, scaling, mirroring, assembling, and grouping}
\vspace{2pt}\hrule height 0.3pt\vspace{3pt}
\textbf{Guide / Hint:} \\
{\small A smooth 3-dimensional object (such as a torus or a dumbbell shape) can possess regions of positive Gaussian curvature where it is locally convex, as well as regions of negative curvature (saddle-like regions) where it is locally concave. Hence, its surface may be concave in some places and convex in others.}

\end{tcolorbox}
\vspace{6pt}\hrule height 0.35pt\vspace{6pt}
\noindent
\begin{tcolorbox}[
    blanker,
    sidebyside,
    sidebyside align=top,
    lefthand width=0.55\textwidth,
    top=4pt, bottom=4pt, left=0pt, right=0pt,
    middle=16pt,
    underlay={
        \draw[dashed, black!80, line width=0.8pt, dash pattern=on 4pt off 3pt]
            ([xshift=0.55\textwidth+8pt]frame.north west) --
            ([xshift=0.55\textwidth+8pt]frame.south west);
    }
]
\textbf{Q.17} \quad {\footnotesize\textsc{[MCQ $\cdot$ 2 Mark(s) $\cdot$ Neg: 0.66]}} \hfill {\footnotesize\textbf{GATE 2023 CS} \quad \tiny\texttt{2023\_CS\_GA\_Q06}}
\vspace{2pt}\hrule height 0.3pt\vspace{4pt}

The country of Zombieland is in distress since more than 75\% of its working population is suffering from serious health issues. Studies conducted by competent health experts concluded that a complete lack of physical exercise among its working population was one of the leading causes of their health issues. As one of the measures to address the problem, the Government of Zombieland has decided to provide monetary incentives to those who ride bicycles to work.
\par\medskip
Based only on the information provided above, which one of the following statements can be logically inferred with certainty?

\par\vspace{6pt}\noindent
\begin{tabular}{@{}p{0.98\linewidth}@{}}
$\square$\ \textbf{(A)} All the working population of Zombieland will henceforth ride bicycles to work. \\
\addlinespace[3pt]
$\square$\ \textbf{(B)} Riding bicycles will ensure that all of the working population of Zombieland is free of health issues. \\
\addlinespace[3pt]
$\square$\ \textbf{(C)} The health experts suggested to the Government of Zombieland to declare riding bicycles as mandatory. \\
\addlinespace[3pt]
$\square$\ \textbf{(D)} The Government of Zombieland believes that riding bicycles is a form of physical exercise.
\end{tabular}

\tcblower

\textbf{Official Key:} \textbf{\large D} \hfill \textit{\small [Easy]}
\vspace{2pt}\hrule height 0.3pt\vspace{3pt}
\textbf{Topic:} Verbal Aptitude \\
\textbf{Subtopic:} \textit{Reading and comprehension}
\vspace{2pt}\hrule height 0.3pt\vspace{3pt}
\textbf{Guide / Hint:} \\
{\small The passage establishes that a complete lack of physical exercise is a primary cause of health issues, and in response, the government provides financial incentives to ride bicycles to work. This direct causal measure logically implies that the government views bicycle riding as a form of physical exercise.}

\end{tcolorbox}
\vspace{6pt}\hrule height 0.35pt\vspace{6pt}
\noindent
\begin{tcolorbox}[
    blanker,
    sidebyside,
    sidebyside align=top,
    lefthand width=0.55\textwidth,
    top=4pt, bottom=4pt, left=0pt, right=0pt,
    middle=16pt,
    underlay={
        \draw[dashed, black!80, line width=0.8pt, dash pattern=on 4pt off 3pt]
            ([xshift=0.55\textwidth+8pt]frame.north west) --
            ([xshift=0.55\textwidth+8pt]frame.south west);
    }
]
\textbf{Q.18} \quad {\footnotesize\textsc{[MCQ $\cdot$ 2 Mark(s) $\cdot$ Neg: 0.66]}} \hfill {\footnotesize\textbf{GATE 2023 CS} \quad \tiny\texttt{2023\_CS\_GA\_Q07}}
\vspace{2pt}\hrule height 0.3pt\vspace{4pt}

Consider two functions of time ($t$),
\[ f(t) = 0.01 t^2, \quad g(t) = 4 t \]
where $0 < t < \infty$.
Now consider the following two statements:
\begin{enumerate}[label=(\roman*),itemsep=0pt,topsep=2pt]
\item For some $t > 0$, $g(t) > f(t)$.
\item There exists a $T$, such that $f(t) > g(t)$ for all $t > T$.
\end{enumerate}
Which one of the following options is TRUE?

\par\vspace{6pt}\noindent
\begin{tabular}{@{}p{0.48\linewidth}@{\hspace{0.04\linewidth}}p{0.48\linewidth}@{}}
$\square$\ \textbf{(A)} only (i) is correct & $\square$\ \textbf{(B)} only (ii) is correct \\
\addlinespace[3pt]
$\square$\ \textbf{(C)} both (i) and (ii) are correct & $\square$\ \textbf{(D)} neither (i) nor (ii) is correct
\end{tabular}

\tcblower

\textbf{Official Key:} \textbf{\large C} \hfill \textit{\small [Easy]}
\vspace{2pt}\hrule height 0.3pt\vspace{3pt}
\textbf{Topic:} Quantitative Aptitude \\
\textbf{Subtopic:} \textit{Numerical computation and estimation: ratios, percentages, powers, exponents and logarithms, permutations and combinations, series}
\vspace{2pt}\hrule height 0.3pt\vspace{3pt}
\textbf{Guide / Hint:} \\
{\small Setting $f(t) = g(t) \implies 0.01 t^2 = 4t \implies t = 400$. For any $0 < t < 400$ (e.g., $t = 10$), $g(10) = 40 > f(10) = 1$, making statement (i) correct. For all $t > 400$, the quadratic growth of $f(t)$ strictly dominates the linear growth of $g(t)$, so choosing $T = 400$ satisfies statement (ii). Hence, both statements are correct.}

\end{tcolorbox}
\vspace{6pt}\hrule height 0.35pt\vspace{6pt}
\noindent
\begin{tcolorbox}[
    blanker,
    sidebyside,
    sidebyside align=top,
    lefthand width=0.55\textwidth,
    top=4pt, bottom=4pt, left=0pt, right=0pt,
    middle=16pt,
    underlay={
        \draw[dashed, black!80, line width=0.8pt, dash pattern=on 4pt off 3pt]
            ([xshift=0.55\textwidth+8pt]frame.north west) --
            ([xshift=0.55\textwidth+8pt]frame.south west);
    }
]
\textbf{Q.19} \quad {\footnotesize\textsc{[MCQ $\cdot$ 2 Mark(s) $\cdot$ Neg: 0.66]}} \hfill {\footnotesize\textbf{GATE 2023 CS} \quad \tiny\texttt{2023\_CS\_GA\_Q08}}
\vspace{2pt}\hrule height 0.3pt\vspace{4pt}

Which one of the following sentence sequences creates a coherent narrative?
\begin{enumerate}[label=(\roman*),itemsep=0pt,topsep=2pt]
\item Once on the terrace, on her way to her small room in the corner, she notices the man right away.
\item She begins to pant by the time she has climbed all the stairs.
\item Mina has bought vegetables and rice at the market, so her bags are heavy.
\item He was leaning against the parapet, watching the traffic below.
\end{enumerate}

\par\vspace{6pt}\noindent
\begin{tabular}{@{}p{0.48\linewidth}@{\hspace{0.04\linewidth}}p{0.48\linewidth}@{}}
$\square$\ \textbf{(A)} (i), (ii), (iv), (iii) & $\square$\ \textbf{(B)} (ii), (iii), (i), (iv) \\
\addlinespace[3pt]
$\square$\ \textbf{(C)} (iv), (ii), (i), (iii) & $\square$\ \textbf{(D)} (iii), (ii), (i), (iv)
\end{tabular}

\tcblower

\textbf{Official Key:} \textbf{\large D} \hfill \textit{\small [Easy]}
\vspace{2pt}\hrule height 0.3pt\vspace{3pt}
\textbf{Topic:} Verbal Aptitude \\
\textbf{Subtopic:} \textit{Narrative sequencing}
\vspace{2pt}\hrule height 0.3pt\vspace{3pt}
\textbf{Guide / Hint:} \\
{\small Sentence (iii) introduces the protagonist Mina carrying heavy market purchases. Sentence (ii) describes the resulting exertion as she climbs the stairs. Sentence (i) brings her onto the terrace where she spots a man. Sentence (iv) details the man's posture against the parapet. Thus, the coherent sequence is (iii), (ii), (i), (iv).}

\end{tcolorbox}
\vspace{6pt}\hrule height 0.35pt\vspace{6pt}
\noindent
\begin{tcolorbox}[
    blanker,
    sidebyside,
    sidebyside align=top,
    lefthand width=0.55\textwidth,
    top=4pt, bottom=4pt, left=0pt, right=0pt,
    middle=16pt,
    underlay={
        \draw[dashed, black!80, line width=0.8pt, dash pattern=on 4pt off 3pt]
            ([xshift=0.55\textwidth+8pt]frame.north west) --
            ([xshift=0.55\textwidth+8pt]frame.south west);
    }
]
\textbf{Q.20} \quad {\footnotesize\textsc{[MCQ $\cdot$ 2 Mark(s) $\cdot$ Neg: 0.66]}} \hfill {\footnotesize\textbf{GATE 2023 CS} \quad \tiny\texttt{2023\_CS\_GA\_Q09}}
\vspace{2pt}\hrule height 0.3pt\vspace{4pt}

$f(x)$ and $g(y)$ are functions of $x$ and $y$, respectively, and $f(x) = g(y)$ for all real values of $x$ and $y$. Which one of the following options is necessarily TRUE for all $x$ and $y$?

\par\vspace{6pt}\noindent
\begin{tabular}{@{}p{0.98\linewidth}@{}}
$\square$\ \textbf{(A)} $f(x) = 0$ and $g(y) = 0$ \\
\addlinespace[3pt]
$\square$\ \textbf{(B)} $f(x) = g(y) = \text{constant}$ \\
\addlinespace[3pt]
$\square$\ \textbf{(C)} $f(x) \neq \text{constant}$ and $g(y) \neq \text{constant}$ \\
\addlinespace[3pt]
$\square$\ \textbf{(D)} $f(x) + g(y) = f(x) - g(y)$
\end{tabular}

\tcblower

\textbf{Official Key:} \textbf{\large B} \hfill \textit{\small [Easy]}
\vspace{2pt}\hrule height 0.3pt\vspace{3pt}
\textbf{Topic:} Quantitative Aptitude \\
\textbf{Subtopic:} \textit{Numerical computation and estimation: ratios, percentages, powers, exponents and logarithms, permutations and combinations, series}
\vspace{2pt}\hrule height 0.3pt\vspace{3pt}
\textbf{Guide / Hint:} \\
{\small We are given $f(x) = g(y)$ for all independent real variables $x$ and $y$. If $x$ varies while $y$ remains constant, $f(x)$ cannot change because $g(y)$ is fixed, meaning $f(x)$ must be a constant. Symmetrically, $g(y)$ must be identical to that same constant. Therefore, $f(x) = g(y) = \text{constant}$.}

\end{tcolorbox}
\vspace{6pt}\hrule height 0.35pt\vspace{6pt}
\noindent
\begin{tcolorbox}[
    blanker,
    sidebyside,
    sidebyside align=top,
    lefthand width=0.55\textwidth,
    top=4pt, bottom=4pt, left=0pt, right=0pt,
    middle=16pt,
    underlay={
        \draw[dashed, black!80, line width=0.8pt, dash pattern=on 4pt off 3pt]
            ([xshift=0.55\textwidth+8pt]frame.north west) --
            ([xshift=0.55\textwidth+8pt]frame.south west);
    }
]
\textbf{Q.21} \quad {\footnotesize\textsc{[MCQ $\cdot$ 1 Mark(s) $\cdot$ Neg: 0.33]}} \hfill {\footnotesize\textbf{GATE 2024 CS1} \quad \tiny\texttt{2024\_CS1\_GA\_Q01}}
\vspace{2pt}\hrule height 0.3pt\vspace{4pt}

If `$\to$' denotes increasing order of intensity, then the meaning of the words $[\text{dry} \to \text{arid} \to \text{parched}]$ is analogous to $[\text{diet} \to \text{fast} \to \text{\_\_\_\_\_\_\_\_}]$.
Which one of the given options is appropriate to fill the blank?

\par\vspace{6pt}\noindent
\begin{tabular}{@{}p{0.48\linewidth}@{\hspace{0.04\linewidth}}p{0.48\linewidth}@{}}
$\square$\ \textbf{(A)} starve & $\square$\ \textbf{(B)} reject \\
\addlinespace[3pt]
$\square$\ \textbf{(C)} feast & $\square$\ \textbf{(D)} deny
\end{tabular}

\tcblower

\textbf{Official Key:} \textbf{\large A} \hfill \textit{\small [Easy]}
\vspace{2pt}\hrule height 0.3pt\vspace{3pt}
\textbf{Topic:} Verbal Aptitude \\
\textbf{Subtopic:} \textit{Vocabulary: words, idioms, and phrases in context}
\vspace{2pt}\hrule height 0.3pt\vspace{3pt}
\textbf{Guide / Hint:} \\
{\small 'dry $\to$ arid $\to$ parched' depicts an increasing scale of lack of moisture. Similarly, 'diet $\to$ fast $\to$ starve' depicts an increasing scale of food restriction/deprivation, where starving represents extreme lack of nourishment.}

\end{tcolorbox}
\vspace{6pt}\hrule height 0.35pt\vspace{6pt}
\noindent
\begin{tcolorbox}[
    blanker,
    sidebyside,
    sidebyside align=top,
    lefthand width=0.55\textwidth,
    top=4pt, bottom=4pt, left=0pt, right=0pt,
    middle=16pt,
    underlay={
        \draw[dashed, black!80, line width=0.8pt, dash pattern=on 4pt off 3pt]
            ([xshift=0.55\textwidth+8pt]frame.north west) --
            ([xshift=0.55\textwidth+8pt]frame.south west);
    }
]
\textbf{Q.22} \quad {\footnotesize\textsc{[MCQ $\cdot$ 1 Mark(s) $\cdot$ Neg: 0.33]}} \hfill {\footnotesize\textbf{GATE 2024 CS1} \quad \tiny\texttt{2024\_CS1\_GA\_Q02}}
\vspace{2pt}\hrule height 0.3pt\vspace{4pt}

If two distinct non-zero real variables $x$ and $y$ are such that $(x+y)$ is proportional to $(x-y)$ then the value of $\frac{x}{y}$

\par\vspace{6pt}\noindent
\begin{tabular}{@{}p{0.98\linewidth}@{}}
$\square$\ \textbf{(A)} depends on $xy$ \\
\addlinespace[3pt]
$\square$\ \textbf{(B)} depends only on $x$ and not on $y$ \\
\addlinespace[3pt]
$\square$\ \textbf{(C)} depends only on $y$ and not on $x$ \\
\addlinespace[3pt]
$\square$\ \textbf{(D)} is a constant
\end{tabular}

\tcblower

\textbf{Official Key:} \textbf{\large D} \hfill \textit{\small [Easy]}
\vspace{2pt}\hrule height 0.3pt\vspace{3pt}
\textbf{Topic:} Quantitative Aptitude \\
\textbf{Subtopic:} \textit{Numerical computation and estimation: ratios, percentages, powers, exponents and logarithms, permutations and combinations, series}
\vspace{2pt}\hrule height 0.3pt\vspace{3pt}
\textbf{Guide / Hint:} \\
{\small If $(x+y) \propto (x-y)$, then $x+y = k(x-y)$ for some non-zero constant $k \neq 1$. Rearranging terms gives $x(1-k) = -y(1+k) \implies \frac{x}{y} = \frac{k+1}{k-1}$, which is a constant independent of $x$ and $y$.}

\end{tcolorbox}
\vspace{6pt}\hrule height 0.35pt\vspace{6pt}
\noindent
\begin{tcolorbox}[
    blanker,
    sidebyside,
    sidebyside align=top,
    lefthand width=0.55\textwidth,
    top=4pt, bottom=4pt, left=0pt, right=0pt,
    middle=16pt,
    underlay={
        \draw[dashed, black!80, line width=0.8pt, dash pattern=on 4pt off 3pt]
            ([xshift=0.55\textwidth+8pt]frame.north west) --
            ([xshift=0.55\textwidth+8pt]frame.south west);
    }
]
\textbf{Q.23} \quad {\footnotesize\textsc{[MCQ $\cdot$ 1 Mark(s) $\cdot$ Neg: 0.33]}} \hfill {\footnotesize\textbf{GATE 2024 CS1} \quad \tiny\texttt{2024\_CS1\_GA\_Q03}}
\vspace{2pt}\hrule height 0.3pt\vspace{4pt}

Consider the following sample of numbers:
\[ 9, 18, 11, 14, 15, 17, 10, 69, 11, 13 \]
The median of the sample is

\par\vspace{6pt}\noindent
\begin{tabular}{@{}p{0.48\linewidth}@{\hspace{0.04\linewidth}}p{0.48\linewidth}@{}}
$\square$\ \textbf{(A)} 13.5 & $\square$\ \textbf{(B)} 14 \\
\addlinespace[3pt]
$\square$\ \textbf{(C)} 11 & $\square$\ \textbf{(D)} 18.7
\end{tabular}

\tcblower

\textbf{Official Key:} \textbf{\large A} \hfill \textit{\small [Easy]}
\vspace{2pt}\hrule height 0.3pt\vspace{3pt}
\textbf{Topic:} Quantitative Aptitude \\
\textbf{Subtopic:} \textit{Elementary statistics and probability}
\vspace{2pt}\hrule height 0.3pt\vspace{3pt}
\textbf{Guide / Hint:} \\
{\small Arranging the 10 sample values in ascending order: $9, 10, 11, 11, 13, 14, 15, 17, 18, 69$. Since the sample size $n = 10$ is even, the median is the arithmetic mean of the 5th and 6th values: $\frac{13 + 14}{2} = 13.5$.}

\end{tcolorbox}
\vspace{6pt}\hrule height 0.35pt\vspace{6pt}
\noindent
\begin{tcolorbox}[
    blanker,
    sidebyside,
    sidebyside align=top,
    lefthand width=0.55\textwidth,
    top=4pt, bottom=4pt, left=0pt, right=0pt,
    middle=16pt,
    underlay={
        \draw[dashed, black!80, line width=0.8pt, dash pattern=on 4pt off 3pt]
            ([xshift=0.55\textwidth+8pt]frame.north west) --
            ([xshift=0.55\textwidth+8pt]frame.south west);
    }
]
\textbf{Q.24} \quad {\footnotesize\textsc{[MCQ $\cdot$ 1 Mark(s) $\cdot$ Neg: 0.33]}} \hfill {\footnotesize\textbf{GATE 2024 CS1} \quad \tiny\texttt{2024\_CS1\_GA\_Q04}}
\vspace{2pt}\hrule height 0.3pt\vspace{4pt}

The number of coins of Rs.~1, Rs.~5, and Rs.~10 denominations that a person has are in the ratio $5:3:13$. Of the total amount, the percentage of money in Rs.~5 coins is

\par\vspace{6pt}\noindent
\begin{tabular}{@{}p{0.48\linewidth}@{\hspace{0.04\linewidth}}p{0.48\linewidth}@{}}
$\square$\ \textbf{(A)} 21\% & $\square$\ \textbf{(B)} $14\frac{2}{7}\%$ \\
\addlinespace[3pt]
$\square$\ \textbf{(C)} 10\% & $\square$\ \textbf{(D)} 30\%
\end{tabular}

\tcblower

\textbf{Official Key:} \textbf{\large C} \hfill \textit{\small [Easy]}
\vspace{2pt}\hrule height 0.3pt\vspace{3pt}
\textbf{Topic:} Quantitative Aptitude \\
\textbf{Subtopic:} \textit{Numerical computation and estimation: ratios, percentages, powers, exponents and logarithms, permutations and combinations, series}
\vspace{2pt}\hrule height 0.3pt\vspace{3pt}
\textbf{Guide / Hint:} \\
{\small Let the number of Rs.~1, Rs.~5, and Rs.~10 coins be $5k$, $3k$, and $13k$ respectively. The total value is $1(5k) + 5(3k) + 10(13k) = 5k + 15k + 130k = 150k$. The percentage of money in Rs.~5 coins is $\frac{15k}{150k} \times 100\% = 10\%$.}

\end{tcolorbox}
\vspace{6pt}\hrule height 0.35pt\vspace{6pt}
\noindent
\begin{tcolorbox}[
    blanker,
    sidebyside,
    sidebyside align=top,
    lefthand width=0.55\textwidth,
    top=4pt, bottom=4pt, left=0pt, right=0pt,
    middle=16pt,
    underlay={
        \draw[dashed, black!80, line width=0.8pt, dash pattern=on 4pt off 3pt]
            ([xshift=0.55\textwidth+8pt]frame.north west) --
            ([xshift=0.55\textwidth+8pt]frame.south west);
    }
]
\textbf{Q.25} \quad {\footnotesize\textsc{[MCQ $\cdot$ 2 Mark(s) $\cdot$ Neg: 0.66]}} \hfill {\footnotesize\textbf{GATE 2024 CS1} \quad \tiny\texttt{2024\_CS1\_GA\_Q06}}
\vspace{2pt}\hrule height 0.3pt\vspace{4pt}

In the given text, the blanks are numbered (i)--(iv). Select the best match for all the blanks.
\par\medskip
Steve was advised to keep his head (i)\_\_\_\_ before heading (ii)\_\_\_\_ to bat; for, while he had a head (iii)\_\_\_\_ batting, he could only do so with a cool head (iv)\_\_\_\_ his shoulders.

\par\vspace{6pt}\noindent
\begin{tabular}{@{}p{0.98\linewidth}@{}}
$\square$\ \textbf{(A)} (i) down (ii) down (iii) on (iv) for \\
\addlinespace[3pt]
$\square$\ \textbf{(B)} (i) on (ii) down (iii) for (iv) on \\
\addlinespace[3pt]
$\square$\ \textbf{(C)} (i) down (ii) out (iii) for (iv) on \\
\addlinespace[3pt]
$\square$\ \textbf{(D)} (i) on (ii) out (iii) on (iv) for
\end{tabular}

\tcblower

\textbf{Official Key:} \textbf{\large C} \hfill \textit{\small [Easy]}
\vspace{2pt}\hrule height 0.3pt\vspace{3pt}
\textbf{Topic:} Verbal Aptitude \\
\textbf{Subtopic:} \textit{Vocabulary: words, idioms, and phrases in context}
\vspace{2pt}\hrule height 0.3pt\vspace{3pt}
\textbf{Guide / Hint:} \\
{\small 'Keep one's head down' means to stay focused and avoid trouble. 'Heading out' means going forth/departing. Having a 'head for' an activity signifies natural aptitude for it, and having a cool 'head on' one's shoulders means remaining calm and composed.}

\end{tcolorbox}
\vspace{6pt}\hrule height 0.35pt\vspace{6pt}
\noindent
\begin{tcolorbox}[
    blanker,
    sidebyside,
    sidebyside align=top,
    lefthand width=0.55\textwidth,
    top=4pt, bottom=4pt, left=0pt, right=0pt,
    middle=16pt,
    underlay={
        \draw[dashed, black!80, line width=0.8pt, dash pattern=on 4pt off 3pt]
            ([xshift=0.55\textwidth+8pt]frame.north west) --
            ([xshift=0.55\textwidth+8pt]frame.south west);
    }
]
\textbf{Q.26} \quad {\footnotesize\textsc{[MCQ $\cdot$ 1 Mark(s) $\cdot$ Neg: 0.33]}} \hfill {\footnotesize\textbf{GATE 2024 CS2} \quad \tiny\texttt{2024\_CS2\_GA\_Q01}}
\vspace{2pt}\hrule height 0.3pt\vspace{4pt}

If `$\to$' denotes increasing order of intensity, then the meaning of the words $[\text{walk} \to \text{jog} \to \text{sprint}]$ is analogous to $[\text{bothered} \to \text{\_\_\_\_\_\_\_\_} \to \text{daunted}]$.
Which one of the given options is appropriate to fill the blank?

\par\vspace{6pt}\noindent
\begin{tabular}{@{}p{0.48\linewidth}@{\hspace{0.04\linewidth}}p{0.48\linewidth}@{}}
$\square$\ \textbf{(A)} phased & $\square$\ \textbf{(B)} phrased \\
\addlinespace[3pt]
$\square$\ \textbf{(C)} fazed & $\square$\ \textbf{(D)} fused
\end{tabular}

\tcblower

\textbf{Official Key:} \textbf{\large C} \hfill \textit{\small [Easy]}
\vspace{2pt}\hrule height 0.3pt\vspace{3pt}
\textbf{Topic:} Verbal Aptitude \\
\textbf{Subtopic:} \textit{Vocabulary: words, idioms, and phrases in context}
\vspace{2pt}\hrule height 0.3pt\vspace{3pt}
\textbf{Guide / Hint:} \\
{\small 'Walk -> jog -> sprint' depicts an increasing scale of physical speed/exertion. Similarly, 'bothered -> fazed -> daunted' depicts an increasing degree of perturbation/discouragement. 'Fazed' means disconcerted or perturbed, fitting between bothered and daunted.}

\end{tcolorbox}
\vspace{6pt}\hrule height 0.35pt\vspace{6pt}
\noindent
\begin{tcolorbox}[
    blanker,
    sidebyside,
    sidebyside align=top,
    lefthand width=0.55\textwidth,
    top=4pt, bottom=4pt, left=0pt, right=0pt,
    middle=16pt,
    underlay={
        \draw[dashed, black!80, line width=0.8pt, dash pattern=on 4pt off 3pt]
            ([xshift=0.55\textwidth+8pt]frame.north west) --
            ([xshift=0.55\textwidth+8pt]frame.south west);
    }
]
\textbf{Q.27} \quad {\footnotesize\textsc{[MCQ $\cdot$ 1 Mark(s) $\cdot$ Neg: 0.33]}} \hfill {\footnotesize\textbf{GATE 2024 CS2} \quad \tiny\texttt{2024\_CS2\_GA\_Q02}}
\vspace{2pt}\hrule height 0.3pt\vspace{4pt}

Two wizards try to create a spell using all the four elements, water, air, fire, and earth. For this, they decide to mix all these elements in all possible orders. They also decide to work independently. After trying all possible combination of elements, they conclude that the spell does not work.
How many attempts does each wizard make before coming to this conclusion, independently?

\par\vspace{6pt}\noindent
\begin{tabular}{@{}p{0.48\linewidth}@{\hspace{0.04\linewidth}}p{0.48\linewidth}@{}}
$\square$\ \textbf{(A)} 24 & $\square$\ \textbf{(B)} 48 \\
\addlinespace[3pt]
$\square$\ \textbf{(C)} 16 & $\square$\ \textbf{(D)} 12
\end{tabular}

\tcblower

\textbf{Official Key:} \textbf{\large A} \hfill \textit{\small [Easy]}
\vspace{2pt}\hrule height 0.3pt\vspace{3pt}
\textbf{Topic:} Quantitative Aptitude \\
\textbf{Subtopic:} \textit{Numerical computation and estimation: ratios, percentages, powers, exponents and logarithms, permutations and combinations, series}
\vspace{2pt}\hrule height 0.3pt\vspace{3pt}
\textbf{Guide / Hint:} \\
{\small There are 4 distinct elements (water, air, fire, earth). The number of distinct permutations/orderings of these 4 elements is $4! = 4 \times 3 \times 2 \times 1 = 24$. Since each wizard tries all possible orders independently, each wizard makes 24 attempts.}

\end{tcolorbox}
\vspace{6pt}\hrule height 0.35pt\vspace{6pt}
\noindent
\begin{tcolorbox}[
    blanker,
    sidebyside,
    sidebyside align=top,
    lefthand width=0.55\textwidth,
    top=4pt, bottom=4pt, left=0pt, right=0pt,
    middle=16pt,
    underlay={
        \draw[dashed, black!80, line width=0.8pt, dash pattern=on 4pt off 3pt]
            ([xshift=0.55\textwidth+8pt]frame.north west) --
            ([xshift=0.55\textwidth+8pt]frame.south west);
    }
]
\textbf{Q.28} \quad {\footnotesize\textsc{[MCQ $\cdot$ 1 Mark(s) $\cdot$ Neg: 0.33]}} \hfill {\footnotesize\textbf{GATE 2024 CS2} \quad \tiny\texttt{2024\_CS2\_GA\_Q04}}
\vspace{2pt}\hrule height 0.3pt\vspace{4pt}

For positive non-zero real variables $x$ and $y$, if
\[ \ln\left(\frac{x+y}{2}\right) = \frac{1}{2} [\ln(x) + \ln(y)] \]
then, the value of $\frac{x}{y} + \frac{y}{x}$ is:

\par\vspace{6pt}\noindent
\begin{tabular}{@{}p{0.48\linewidth}@{\hspace{0.04\linewidth}}p{0.48\linewidth}@{}}
$\square$\ \textbf{(A)} 1 & $\square$\ \textbf{(B)} 1/2 \\
\addlinespace[3pt]
$\square$\ \textbf{(C)} 2 & $\square$\ \textbf{(D)} 4
\end{tabular}

\tcblower

\textbf{Official Key:} \textbf{\large C} \hfill \textit{\small [Easy]}
\vspace{2pt}\hrule height 0.3pt\vspace{3pt}
\textbf{Topic:} Quantitative Aptitude \\
\textbf{Subtopic:} \textit{Numerical computation and estimation: ratios, percentages, powers, exponents and logarithms, permutations and combinations, series}
\vspace{2pt}\hrule height 0.3pt\vspace{3pt}
\textbf{Guide / Hint:} \\
{\small From $\ln\left(\frac{x+y}{2}\right) = \frac{1}{2}[\ln(x) + \ln(y)] = \ln(\sqrt{xy})$, we obtain $\frac{x+y}{2} = \sqrt{xy} \implies x + y = 2\sqrt{xy} \implies (\sqrt{x} - \sqrt{y})^2 = 0 \implies x = y$. Substituting $x = y$ gives $\frac{x}{y} + \frac{y}{x} = 1 + 1 = 2$.}

\end{tcolorbox}
\vspace{6pt}\hrule height 0.35pt\vspace{6pt}
\noindent
\begin{tcolorbox}[
    blanker,
    sidebyside,
    sidebyside align=top,
    lefthand width=0.55\textwidth,
    top=4pt, bottom=4pt, left=0pt, right=0pt,
    middle=16pt,
    underlay={
        \draw[dashed, black!80, line width=0.8pt, dash pattern=on 4pt off 3pt]
            ([xshift=0.55\textwidth+8pt]frame.north west) --
            ([xshift=0.55\textwidth+8pt]frame.south west);
    }
]
\textbf{Q.29} \quad {\footnotesize\textsc{[MCQ $\cdot$ 1 Mark(s) $\cdot$ Neg: 0.33]}} \hfill {\footnotesize\textbf{GATE 2024 CS2} \quad \tiny\texttt{2024\_CS2\_GA\_Q05}}
\vspace{2pt}\hrule height 0.3pt\vspace{4pt}

In the sequence 6, 9, 14, $x$, 30, 41, a possible value of $x$ is:

\par\vspace{6pt}\noindent
\begin{tabular}{@{}p{0.48\linewidth}@{\hspace{0.04\linewidth}}p{0.48\linewidth}@{}}
$\square$\ \textbf{(A)} 25 & $\square$\ \textbf{(B)} 21 \\
\addlinespace[3pt]
$\square$\ \textbf{(C)} 18 & $\square$\ \textbf{(D)} 20
\end{tabular}

\tcblower

\textbf{Official Key:} \textbf{\large B} \hfill \textit{\small [Easy]}
\vspace{2pt}\hrule height 0.3pt\vspace{3pt}
\textbf{Topic:} Quantitative Aptitude \\
\textbf{Subtopic:} \textit{Numerical computation and estimation: ratios, percentages, powers, exponents and logarithms, permutations and combinations, series}
\vspace{2pt}\hrule height 0.3pt\vspace{3pt}
\textbf{Guide / Hint:} \\
{\small Examine the consecutive differences: $9 - 6 = 3$, $14 - 9 = 5$, $x - 14 = d_3$, $30 - x = d_4$, $41 - 30 = 11$. The differences are consecutive odd numbers: $3, 5, 7, 9, 11$. For $d_3 = 7$, $x = 14 + 7 = 21$, and checking $d_4 = 30 - 21 = 9$ matches. Hence, $x = 21$.}

\end{tcolorbox}
\vspace{6pt}\hrule height 0.35pt\vspace{6pt}
\noindent
\begin{tcolorbox}[
    blanker,
    sidebyside,
    sidebyside align=top,
    lefthand width=0.55\textwidth,
    top=4pt, bottom=4pt, left=0pt, right=0pt,
    middle=16pt,
    underlay={
        \draw[dashed, black!80, line width=0.8pt, dash pattern=on 4pt off 3pt]
            ([xshift=0.55\textwidth+8pt]frame.north west) --
            ([xshift=0.55\textwidth+8pt]frame.south west);
    }
]
\textbf{Q.30} \quad {\footnotesize\textsc{[MCQ $\cdot$ 2 Mark(s) $\cdot$ Neg: 0.66]}} \hfill {\footnotesize\textbf{GATE 2024 CS2} \quad \tiny\texttt{2024\_CS2\_GA\_Q07}}
\vspace{2pt}\hrule height 0.3pt\vspace{4pt}

A person sold two different items at the same price. He made 10\% profit in one item, and 10\% loss in the other item. In selling these two items, the person made a total of:

\par\vspace{6pt}\noindent
\begin{tabular}{@{}p{0.48\linewidth}@{\hspace{0.04\linewidth}}p{0.48\linewidth}@{}}
$\square$\ \textbf{(A)} 1\% profit & $\square$\ \textbf{(B)} 2\% profit \\
\addlinespace[3pt]
$\square$\ \textbf{(C)} 1\% loss & $\square$\ \textbf{(D)} 2\% loss
\end{tabular}

\tcblower

\textbf{Official Key:} \textbf{\large C} \hfill \textit{\small [Easy]}
\vspace{2pt}\hrule height 0.3pt\vspace{3pt}
\textbf{Topic:} Quantitative Aptitude \\
\textbf{Subtopic:} \textit{Numerical computation and estimation: ratios, percentages, powers, exponents and logarithms, permutations and combinations, series}
\vspace{2pt}\hrule height 0.3pt\vspace{3pt}
\textbf{Guide / Hint:} \\
{\small When two items are sold at the same selling price $S$, one at a profit of $x\%$ and the other at a loss of $x\%$, there is always an overall net loss given by $\frac{x^2}{100}\%$. Here $x = 10$, so the net loss is $\frac{10^2}{100}\% = 1\%$ loss.}

\end{tcolorbox}
\vspace{6pt}\hrule height 0.35pt\vspace{6pt}
\noindent
\begin{tcolorbox}[
    blanker,
    sidebyside,
    sidebyside align=top,
    lefthand width=0.55\textwidth,
    top=4pt, bottom=4pt, left=0pt, right=0pt,
    middle=16pt,
    underlay={
        \draw[dashed, black!80, line width=0.8pt, dash pattern=on 4pt off 3pt]
            ([xshift=0.55\textwidth+8pt]frame.north west) --
            ([xshift=0.55\textwidth+8pt]frame.south west);
    }
]
\textbf{Q.31} \quad {\footnotesize\textsc{[MCQ $\cdot$ 2 Mark(s) $\cdot$ Neg: 0.66]}} \hfill {\footnotesize\textbf{GATE 2024 CS2} \quad \tiny\texttt{2024\_CS2\_GA\_Q09}}
\vspace{2pt}\hrule height 0.3pt\vspace{4pt}

A cube is to be cut into 8 pieces of equal size and shape. Here, each cut should be straight and it should not stop till it reaches the other end of the cube.
The minimum number of such cuts required is:

\par\vspace{6pt}\noindent
\begin{tabular}{@{}p{0.48\linewidth}@{\hspace{0.04\linewidth}}p{0.48\linewidth}@{}}
$\square$\ \textbf{(A)} 3 & $\square$\ \textbf{(B)} 4 \\
\addlinespace[3pt]
$\square$\ \textbf{(C)} 7 & $\square$\ \textbf{(D)} 8
\end{tabular}

\tcblower

\textbf{Official Key:} \textbf{\large A} \hfill \textit{\small [Easy]}
\vspace{2pt}\hrule height 0.3pt\vspace{3pt}
\textbf{Topic:} Spatial Aptitude \\
\textbf{Subtopic:} \textit{Paper folding, cutting, and patterns in 2 and 3 dimensions}
\vspace{2pt}\hrule height 0.3pt\vspace{3pt}
\textbf{Guide / Hint:} \\
{\small A single full planar straight cut can at most double the number of pieces. With $k$ cuts through 3 orthogonal dimensions (1 cut parallel to $xy$-plane, 1 parallel to $yz$-plane, 1 parallel to $zx$-plane), the cube is divided into $2 \times 2 \times 2 = 8$ equal pieces. Thus, a minimum of 3 cuts is required.}

\end{tcolorbox}
\vspace{6pt}\hrule height 0.35pt\vspace{6pt}
\noindent
\begin{tcolorbox}[
    blanker,
    sidebyside,
    sidebyside align=top,
    lefthand width=0.55\textwidth,
    top=4pt, bottom=4pt, left=0pt, right=0pt,
    middle=16pt,
    underlay={
        \draw[dashed, black!80, line width=0.8pt, dash pattern=on 4pt off 3pt]
            ([xshift=0.55\textwidth+8pt]frame.north west) --
            ([xshift=0.55\textwidth+8pt]frame.south west);
    }
]
\textbf{Q.32} \quad {\footnotesize\textsc{[MCQ $\cdot$ 1 Mark(s) $\cdot$ Neg: 0.33]}} \hfill {\footnotesize\textbf{GATE 2025 CS1} \quad \tiny\texttt{2025\_CS1\_GA\_Q01}}
\vspace{2pt}\hrule height 0.3pt\vspace{4pt}

Ravi had \_\_\_\_\_\_ younger brother who taught at \_\_\_\_\_\_ university. He was widely regarded as \_\_\_\_\_\_ honorable man.
Select the option with the correct sequence of articles to fill in the blanks.

\par\vspace{6pt}\noindent
\begin{tabular}{@{}p{0.48\linewidth}@{\hspace{0.04\linewidth}}p{0.48\linewidth}@{}}
$\square$\ \textbf{(A)} a; a; an & $\square$\ \textbf{(B)} the; an; a \\
\addlinespace[3pt]
$\square$\ \textbf{(C)} a; an; a & $\square$\ \textbf{(D)} an; an; a
\end{tabular}

\tcblower

\textbf{Official Key:} \textbf{\large A} \hfill \textit{\small [Easy]}
\vspace{2pt}\hrule height 0.3pt\vspace{3pt}
\textbf{Topic:} Verbal Aptitude \\
\textbf{Subtopic:} \textit{Basic English grammar: tenses, articles, adjectives, prepositions, conjunctions, verb-noun agreement}
\vspace{2pt}\hrule height 0.3pt\vspace{3pt}
\textbf{Guide / Hint:} \\
{\small 'Younger' begins with a consonant sound (/j/), requiring the indefinite article 'a'. 'University' also begins with a consonant sound (/juː/), taking 'a'. 'Honorable' has a silent 'h' and begins with a vowel sound (/ɒ/), requiring 'an'. Hence, the sequence is a; a; an.}

\end{tcolorbox}
\vspace{6pt}\hrule height 0.35pt\vspace{6pt}
\noindent
\begin{tcolorbox}[
    blanker,
    sidebyside,
    sidebyside align=top,
    lefthand width=0.55\textwidth,
    top=4pt, bottom=4pt, left=0pt, right=0pt,
    middle=16pt,
    underlay={
        \draw[dashed, black!80, line width=0.8pt, dash pattern=on 4pt off 3pt]
            ([xshift=0.55\textwidth+8pt]frame.north west) --
            ([xshift=0.55\textwidth+8pt]frame.south west);
    }
]
\textbf{Q.33} \quad {\footnotesize\textsc{[MCQ $\cdot$ 1 Mark(s) $\cdot$ Neg: 0.33]}} \hfill {\footnotesize\textbf{GATE 2025 CS1} \quad \tiny\texttt{2025\_CS1\_GA\_Q02}}
\vspace{2pt}\hrule height 0.3pt\vspace{4pt}

The CEO’s decision to downsize the workforce was considered myopic because it sacrificed long-term stability to accommodate short-term gains.
Select the most appropriate option that can replace the word ``myopic'' without changing the meaning of the sentence.

\par\vspace{6pt}\noindent
\begin{tabular}{@{}p{0.48\linewidth}@{\hspace{0.04\linewidth}}p{0.48\linewidth}@{}}
$\square$\ \textbf{(A)} visionary & $\square$\ \textbf{(B)} shortsighted \\
\addlinespace[3pt]
$\square$\ \textbf{(C)} progressive & $\square$\ \textbf{(D)} innovative
\end{tabular}

\tcblower

\textbf{Official Key:} \textbf{\large B} \hfill \textit{\small [Easy]}
\vspace{2pt}\hrule height 0.3pt\vspace{3pt}
\textbf{Topic:} Verbal Aptitude \\
\textbf{Subtopic:} \textit{Vocabulary: words, idioms, and phrases in context}
\vspace{2pt}\hrule height 0.3pt\vspace{3pt}
\textbf{Guide / Hint:} \\
{\small 'Myopic' literally means nearsighted, and metaphorically refers to lacking foresight or being shortsighted by focusing only on immediate considerations rather than long-term outcomes.}

\end{tcolorbox}
\vspace{6pt}\hrule height 0.35pt\vspace{6pt}
\noindent
\begin{tcolorbox}[
    blanker,
    sidebyside,
    sidebyside align=top,
    lefthand width=0.55\textwidth,
    top=4pt, bottom=4pt, left=0pt, right=0pt,
    middle=16pt,
    underlay={
        \draw[dashed, black!80, line width=0.8pt, dash pattern=on 4pt off 3pt]
            ([xshift=0.55\textwidth+8pt]frame.north west) --
            ([xshift=0.55\textwidth+8pt]frame.south west);
    }
]
\textbf{Q.34} \quad {\footnotesize\textsc{[MCQ $\cdot$ 1 Mark(s) $\cdot$ Neg: 0.33]}} \hfill {\footnotesize\textbf{GATE 2025 CS1} \quad \tiny\texttt{2025\_CS1\_GA\_Q03}}
\vspace{2pt}\hrule height 0.3pt\vspace{4pt}

The average marks obtained by a class in an examination were calculated as 30.8. However, while checking the marks entered, the teacher found that the marks of one student were entered incorrectly as 24 instead of 42. After correcting the marks, the average becomes 31.4. How many students does the class have?

\par\vspace{6pt}\noindent
\begin{tabular}{@{}p{0.48\linewidth}@{\hspace{0.04\linewidth}}p{0.48\linewidth}@{}}
$\square$\ \textbf{(A)} 25 & $\square$\ \textbf{(B)} 28 \\
\addlinespace[3pt]
$\square$\ \textbf{(C)} 30 & $\square$\ \textbf{(D)} 32
\end{tabular}

\tcblower

\textbf{Official Key:} \textbf{\large C} \hfill \textit{\small [Easy]}
\vspace{2pt}\hrule height 0.3pt\vspace{3pt}
\textbf{Topic:} Quantitative Aptitude \\
\textbf{Subtopic:} \textit{Elementary statistics and probability}
\vspace{2pt}\hrule height 0.3pt\vspace{3pt}
\textbf{Guide / Hint:} \\
{\small Let $n$ be the number of students. The total increase in marks after correction is $42 - 24 = 18$. The increase in average is $31.4 - 30.8 = 0.6$. Therefore, $0.6 \times n = 18$, giving $n = \frac{18}{0.6} = 30$ students.}

\end{tcolorbox}
\vspace{6pt}\hrule height 0.35pt\vspace{6pt}
\noindent
\begin{tcolorbox}[
    blanker,
    sidebyside,
    sidebyside align=top,
    lefthand width=0.55\textwidth,
    top=4pt, bottom=4pt, left=0pt, right=0pt,
    middle=16pt,
    underlay={
        \draw[dashed, black!80, line width=0.8pt, dash pattern=on 4pt off 3pt]
            ([xshift=0.55\textwidth+8pt]frame.north west) --
            ([xshift=0.55\textwidth+8pt]frame.south west);
    }
]
\textbf{Q.35} \quad {\footnotesize\textsc{[MCQ $\cdot$ 1 Mark(s) $\cdot$ Neg: 0.33]}} \hfill {\footnotesize\textbf{GATE 2025 CS1} \quad \tiny\texttt{2025\_CS1\_GA\_Q04}}
\vspace{2pt}\hrule height 0.3pt\vspace{4pt}

Consider the relationships among P, Q, R, S, and T:
\begin{itemize}
\item P is the brother of Q.
\item S is the daughter of Q.
\item T is the sister of S.
\item R is the mother of Q.
\end{itemize}
The following statements are made based on the relationships given above:
(1) R is the grandmother of S.
(2) P is the uncle of S and T.
(3) R has only one son.
(4) Q has only one daughter.
Which one of the following options is correct?

\par\vspace{6pt}\noindent
\begin{tabular}{@{}p{0.48\linewidth}@{\hspace{0.04\linewidth}}p{0.48\linewidth}@{}}
$\square$\ \textbf{(A)} Both (1) and (2) are true. & $\square$\ \textbf{(B)} Both (1) and (3) are true. \\
\addlinespace[3pt]
$\square$\ \textbf{(C)} Only (3) is true. & $\square$\ \textbf{(D)} Only (4) is true.
\end{tabular}

\tcblower

\textbf{Official Key:} \textbf{\large A} \hfill \textit{\small [Easy]}
\vspace{2pt}\hrule height 0.3pt\vspace{3pt}
\textbf{Topic:} Analytical Aptitude \\
\textbf{Subtopic:} \textit{Logic: deduction and induction}
\vspace{2pt}\hrule height 0.3pt\vspace{3pt}
\textbf{Guide / Hint:} \\
{\small Since R is the mother of Q and S is the daughter of Q, R is definitely the grandmother of S (statement 1 is true). P is the brother of Q, so P is the maternal/paternal uncle of Q's children S and T (statement 2 is true). The genders and number of other children of R and Q are not specified, so (3) and (4) cannot be established.}

\end{tcolorbox}
\vspace{6pt}\hrule height 0.35pt\vspace{6pt}
\noindent
\begin{tcolorbox}[
    blanker,
    sidebyside,
    sidebyside align=top,
    lefthand width=0.55\textwidth,
    top=4pt, bottom=4pt, left=0pt, right=0pt,
    middle=16pt,
    underlay={
        \draw[dashed, black!80, line width=0.8pt, dash pattern=on 4pt off 3pt]
            ([xshift=0.55\textwidth+8pt]frame.north west) --
            ([xshift=0.55\textwidth+8pt]frame.south west);
    }
]
\textbf{Q.36} \quad {\footnotesize\textsc{[MCQ $\cdot$ 2 Mark(s) $\cdot$ Neg: 0.66]}} \hfill {\footnotesize\textbf{GATE 2025 CS1} \quad \tiny\texttt{2025\_CS1\_GA\_Q06}}
\vspace{2pt}\hrule height 0.3pt\vspace{4pt}

``I put the brown paper in my pocket along with the chalks, and possibly other things. I suppose every one must have reflected how primeval and how poetical are the things that one carries in one's pocket: the pocket-knife, for instance the type of all human tools, the infant of the sword. Once I planned to write a book of poems entirely about the things in my pocket. But I found it would be too long: and the age of the great epics is past.''
(From G.K. Chesterton's ``A Piece of Chalk'')
\par\medskip
Based only on the information provided in the above passage, which one of the following statements is true?

\par\vspace{6pt}\noindent
\begin{tabular}{@{}p{0.98\linewidth}@{}}
$\square$\ \textbf{(A)} The author of the passage carries a mirror in his pocket to reflect upon things. \\
\addlinespace[3pt]
$\square$\ \textbf{(B)} The author of the passage had decided to write a poem on epics. \\
\addlinespace[3pt]
$\square$\ \textbf{(C)} The pocket-knife is described as the infant of the sword. \\
\addlinespace[3pt]
$\square$\ \textbf{(D)} Epics are described as too inconvenient to write.
\end{tabular}

\tcblower

\textbf{Official Key:} \textbf{\large C} \hfill \textit{\small [Easy]}
\vspace{2pt}\hrule height 0.3pt\vspace{3pt}
\textbf{Topic:} Verbal Aptitude \\
\textbf{Subtopic:} \textit{Reading and comprehension}
\vspace{2pt}\hrule height 0.3pt\vspace{3pt}
\textbf{Guide / Hint:} \\
{\small The passage explicitly states: 'the pocket-knife, for instance the type of all human tools, the infant of the sword.' Statement C directly reflects this verbatim description.}

\end{tcolorbox}
\vspace{6pt}\hrule height 0.35pt\vspace{6pt}
\noindent
\begin{tcolorbox}[
    blanker,
    sidebyside,
    sidebyside align=top,
    lefthand width=0.55\textwidth,
    top=4pt, bottom=4pt, left=0pt, right=0pt,
    middle=16pt,
    underlay={
        \draw[dashed, black!80, line width=0.8pt, dash pattern=on 4pt off 3pt]
            ([xshift=0.55\textwidth+8pt]frame.north west) --
            ([xshift=0.55\textwidth+8pt]frame.south west);
    }
]
\textbf{Q.37} \quad {\footnotesize\textsc{[MCQ $\cdot$ 2 Mark(s) $\cdot$ Neg: 0.66]}} \hfill {\footnotesize\textbf{GATE 2025 CS1} \quad \tiny\texttt{2025\_CS1\_GA\_Q10}}
\vspace{2pt}\hrule height 0.3pt\vspace{4pt}

A shop has 4 distinct flavors of ice-cream. One can purchase any number of scoops of any flavor. The order in which the scoops are purchased is inconsequential. If one wants to purchase 3 scoops of ice-cream, in how many ways can one make that purchase?

\par\vspace{6pt}\noindent
\begin{tabular}{@{}p{0.48\linewidth}@{\hspace{0.04\linewidth}}p{0.48\linewidth}@{}}
$\square$\ \textbf{(A)} 4 & $\square$\ \textbf{(B)} 20 \\
\addlinespace[3pt]
$\square$\ \textbf{(C)} 24 & $\square$\ \textbf{(D)} 48
\end{tabular}

\tcblower

\textbf{Official Key:} \textbf{\large B} \hfill \textit{\small [Easy]}
\vspace{2pt}\hrule height 0.3pt\vspace{3pt}
\textbf{Topic:} Quantitative Aptitude \\
\textbf{Subtopic:} \textit{Numerical computation and estimation: ratios, percentages, powers, exponents and logarithms, permutations and combinations, series}
\vspace{2pt}\hrule height 0.3pt\vspace{3pt}
\textbf{Guide / Hint:} \\
{\small This is a multichoose problem (combination with repetition), where $r = 3$ scoops are chosen from $n = 4$ flavors. The number of ways is given by the stars and bars formula $\binom{n + r - 1}{r} = \binom{4 + 3 - 1}{3} = \binom{6}{3} = \frac{6 \times 5 \times 4}{3 \times 2 \times 1} = 20$.}

\end{tcolorbox}
\vspace{6pt}\hrule height 0.35pt\vspace{6pt}
\noindent
\begin{tcolorbox}[
    blanker,
    sidebyside,
    sidebyside align=top,
    lefthand width=0.55\textwidth,
    top=4pt, bottom=4pt, left=0pt, right=0pt,
    middle=16pt,
    underlay={
        \draw[dashed, black!80, line width=0.8pt, dash pattern=on 4pt off 3pt]
            ([xshift=0.55\textwidth+8pt]frame.north west) --
            ([xshift=0.55\textwidth+8pt]frame.south west);
    }
]
\textbf{Q.38} \quad {\footnotesize\textsc{[MCQ $\cdot$ 1 Mark(s) $\cdot$ Neg: 0.33]}} \hfill {\footnotesize\textbf{GATE 2025 CS2} \quad \tiny\texttt{2025\_CS2\_GA\_Q01}}
\vspace{2pt}\hrule height 0.3pt\vspace{4pt}

Despite his initial hesitation, Rehman’s \_\_\_\_\_\_\_\_\_ to contribute to the success of the project never wavered.
Select the most appropriate option to complete the above sentence.

\par\vspace{6pt}\noindent
\begin{tabular}{@{}p{0.48\linewidth}@{\hspace{0.04\linewidth}}p{0.48\linewidth}@{}}
$\square$\ \textbf{(A)} ambivalence & $\square$\ \textbf{(B)} satisfaction \\
\addlinespace[3pt]
$\square$\ \textbf{(C)} resolve & $\square$\ \textbf{(D)} revolve
\end{tabular}

\tcblower

\textbf{Official Key:} \textbf{\large C} \hfill \textit{\small [Easy]}
\vspace{2pt}\hrule height 0.3pt\vspace{3pt}
\textbf{Topic:} Verbal Aptitude \\
\textbf{Subtopic:} \textit{Vocabulary: words, idioms, and phrases in context}
\vspace{2pt}\hrule height 0.3pt\vspace{3pt}
\textbf{Guide / Hint:} \\
{\small The sentence sets up a contrast between Rehman's initial hesitation and his continuing commitment. 'Resolve' denotes firm determination and perseverance, which fits the context accurately.}

\end{tcolorbox}
\vspace{6pt}\hrule height 0.35pt\vspace{6pt}
\noindent
\begin{tcolorbox}[
    blanker,
    sidebyside,
    sidebyside align=top,
    lefthand width=0.55\textwidth,
    top=4pt, bottom=4pt, left=0pt, right=0pt,
    middle=16pt,
    underlay={
        \draw[dashed, black!80, line width=0.8pt, dash pattern=on 4pt off 3pt]
            ([xshift=0.55\textwidth+8pt]frame.north west) --
            ([xshift=0.55\textwidth+8pt]frame.south west);
    }
]
\textbf{Q.39} \quad {\footnotesize\textsc{[MCQ $\cdot$ 1 Mark(s) $\cdot$ Neg: 0.33]}} \hfill {\footnotesize\textbf{GATE 2025 CS2} \quad \tiny\texttt{2025\_CS2\_GA\_Q02}}
\vspace{2pt}\hrule height 0.3pt\vspace{4pt}

$\text{Bird} : \text{Nest} :: \text{Bee} : \text{\_\_\_\_\_\_\_}$
Select the correct option to complete the analogy.

\par\vspace{6pt}\noindent
\begin{tabular}{@{}p{0.48\linewidth}@{\hspace{0.04\linewidth}}p{0.48\linewidth}@{}}
$\square$\ \textbf{(A)} Kennel & $\square$\ \textbf{(B)} Hammock \\
\addlinespace[3pt]
$\square$\ \textbf{(C)} Hive & $\square$\ \textbf{(D)} Lair
\end{tabular}

\tcblower

\textbf{Official Key:} \textbf{\large C} \hfill \textit{\small [Easy]}
\vspace{2pt}\hrule height 0.3pt\vspace{3pt}
\textbf{Topic:} Verbal Aptitude \\
\textbf{Subtopic:} \textit{Vocabulary: word analogy and contextual relations}
\vspace{2pt}\hrule height 0.3pt\vspace{3pt}
\textbf{Guide / Hint:} \\
{\small This is a habitat/dwelling analogy: a bird naturally lives in a nest, and similarly, a bee lives in a hive.}

\end{tcolorbox}
\vspace{6pt}\hrule height 0.35pt\vspace{6pt}
\noindent
\begin{tcolorbox}[
    blanker,
    sidebyside,
    sidebyside align=top,
    lefthand width=0.55\textwidth,
    top=4pt, bottom=4pt, left=0pt, right=0pt,
    middle=16pt,
    underlay={
        \draw[dashed, black!80, line width=0.8pt, dash pattern=on 4pt off 3pt]
            ([xshift=0.55\textwidth+8pt]frame.north west) --
            ([xshift=0.55\textwidth+8pt]frame.south west);
    }
]
\textbf{Q.40} \quad {\footnotesize\textsc{[MCQ $\cdot$ 1 Mark(s) $\cdot$ Neg: 0.33]}} \hfill {\footnotesize\textbf{GATE 2025 CS2} \quad \tiny\texttt{2025\_CS2\_GA\_Q05}}
\vspace{2pt}\hrule height 0.3pt\vspace{4pt}

Let $p_1$ and $p_2$ denote two arbitrary prime numbers. Which one of the following statements is correct for all values of $p_1$ and $p_2$?

\par\vspace{6pt}\noindent
\begin{tabular}{@{}p{0.98\linewidth}@{}}
$\square$\ \textbf{(A)} $p_1 + p_2$ is not a prime number. \\
\addlinespace[3pt]
$\square$\ \textbf{(B)} $p_1 p_2$ is not a prime number. \\
\addlinespace[3pt]
$\square$\ \textbf{(C)} $p_1 + p_2 + 1$ is a prime number. \\
\addlinespace[3pt]
$\square$\ \textbf{(D)} $p_1 p_2 + 1$ is a prime number.
\end{tabular}

\tcblower

\textbf{Official Key:} \textbf{\large B} \hfill \textit{\small [Easy]}
\vspace{2pt}\hrule height 0.3pt\vspace{3pt}
\textbf{Topic:} Quantitative Aptitude \\
\textbf{Subtopic:} \textit{Elementary number theory and prime numbers}
\vspace{2pt}\hrule height 0.3pt\vspace{3pt}
\textbf{Guide / Hint:} \\
{\small For any two prime numbers $p_1$ and $p_2$, the product $p_1 p_2$ has at least three distinct positive divisors ($1, p_1, p_1 p_2$ when $p_1=p_2$, or $1, p_1, p_2, p_1 p_2$ when $p_1 \neq p_2$). Hence, $p_1 p_2$ can never be prime.}

\end{tcolorbox}
\vspace{6pt}\hrule height 0.35pt\vspace{6pt}
\noindent
\begin{tcolorbox}[
    blanker,
    sidebyside,
    sidebyside align=top,
    lefthand width=0.55\textwidth,
    top=4pt, bottom=4pt, left=0pt, right=0pt,
    middle=16pt,
    underlay={
        \draw[dashed, black!80, line width=0.8pt, dash pattern=on 4pt off 3pt]
            ([xshift=0.55\textwidth+8pt]frame.north west) --
            ([xshift=0.55\textwidth+8pt]frame.south west);
    }
]
\textbf{Q.41} \quad {\footnotesize\textsc{[MCQ $\cdot$ 2 Mark(s) $\cdot$ Neg: 0.66]}} \hfill {\footnotesize\textbf{GATE 2025 CS2} \quad \tiny\texttt{2025\_CS2\_GA\_Q06}}
\vspace{2pt}\hrule height 0.3pt\vspace{4pt}

Based only on the conversation below, identify the logically correct inference:
``Even if I had known that you were in the hospital, I would not have gone there to see you'', Ramya told Josephine.

\par\vspace{6pt}\noindent
\begin{tabular}{@{}p{0.98\linewidth}@{}}
$\square$\ \textbf{(A)} Ramya knew that Josephine was in the hospital. \\
\addlinespace[3pt]
$\square$\ \textbf{(B)} Ramya did not know that Josephine was in the hospital. \\
\addlinespace[3pt]
$\square$\ \textbf{(C)} Ramya and Josephine were once close friends; but now, they are not. \\
\addlinespace[3pt]
$\square$\ \textbf{(D)} Josephine was in the hospital due to an injury to her leg.
\end{tabular}

\tcblower

\textbf{Official Key:} \textbf{\large B} \hfill \textit{\small [Easy]}
\vspace{2pt}\hrule height 0.3pt\vspace{3pt}
\textbf{Topic:} Analytical Aptitude \\
\textbf{Subtopic:} \textit{Logic: deduction and conditional statements}
\vspace{2pt}\hrule height 0.3pt\vspace{3pt}
\textbf{Guide / Hint:} \\
{\small The subjunctive statement 'Even if I had known that you were in the hospital...' logically presupposes that the speaker did not know about Josephine's hospitalization at that time.}

\end{tcolorbox}
\vspace{6pt}\hrule height 0.35pt\vspace{6pt}
\vspace{8pt}
\subsection*{Engineering Mathematics}
\vspace{-4pt}
\noindent
\begin{tcolorbox}[
    blanker,
    sidebyside,
    sidebyside align=top,
    lefthand width=0.55\textwidth,
    top=4pt, bottom=4pt, left=0pt, right=0pt,
    middle=16pt,
    underlay={
        \draw[dashed, black!80, line width=0.8pt, dash pattern=on 4pt off 3pt]
            ([xshift=0.55\textwidth+8pt]frame.north west) --
            ([xshift=0.55\textwidth+8pt]frame.south west);
    }
]
\textbf{Q.42} \quad {\footnotesize\textsc{[NAT $\cdot$ 1 Mark(s) $\cdot$ No Neg]}} \hfill {\footnotesize\textbf{GATE 2018 CS} \quad \tiny\texttt{2018\_CS\_GA\_Q16}}
\vspace{2pt}\hrule height 0.3pt\vspace{4pt}

The value of $\int_0^{\pi/4} x \cos(x^2) \, dx$ correct to three decimal places (assuming that $\pi = 3.14$) is \_\_\_\_\_\_\_\_.

\par\vspace{6pt}\noindent\textbf{Your Answer:} \fbox{\rule{0pt}{1.3ex}\hspace{3cm}}

\tcblower

\textbf{Official Key:} \textbf{\large 0.27 to 0.30} \hfill \textit{\small [Easy]}
\vspace{2pt}\hrule height 0.3pt\vspace{3pt}
\textbf{Topic:} Calculus \\
\textbf{Subtopic:} \textit{Integration}
\vspace{2pt}\hrule height 0.3pt\vspace{3pt}
\textbf{Guide / Hint:} \\
{\small Substitute $u = x^2$, so $du = 2x\,dx$. The limits transform from $x = 0 \implies u = 0$ to $x = \frac{\pi}{4} \implies u = \frac{\pi^2}{16} \approx \frac{9.8596}{16} \approx 0.61685$ radians. The integral is $\frac{1}{2}\int_0^{\pi^2/16} \cos u \, du = \frac{1}{2} \sin\left(\frac{\pi^2}{16}\right) \approx \frac{1}{2} \sin(0.61685) \approx \frac{1}{2}(0.5784) \approx 0.289$.}

\end{tcolorbox}
\vspace{6pt}\hrule height 0.35pt\vspace{6pt}
\noindent
\begin{tcolorbox}[
    blanker,
    sidebyside,
    sidebyside align=top,
    lefthand width=0.55\textwidth,
    top=4pt, bottom=4pt, left=0pt, right=0pt,
    middle=16pt,
    underlay={
        \draw[dashed, black!80, line width=0.8pt, dash pattern=on 4pt off 3pt]
            ([xshift=0.55\textwidth+8pt]frame.north west) --
            ([xshift=0.55\textwidth+8pt]frame.south west);
    }
]
\textbf{Q.43} \quad {\footnotesize\textsc{[NAT $\cdot$ 1 Mark(s) $\cdot$ No Neg]}} \hfill {\footnotesize\textbf{GATE 2018 CS} \quad \tiny\texttt{2018\_CS\_GA\_Q17}}
\vspace{2pt}\hrule height 0.3pt\vspace{4pt}

Consider a matrix $A = u v^T$ where $u = \begin{bmatrix} 1 \\ 2 \end{bmatrix}$ and $v = \begin{bmatrix} 1 \\ 1 \end{bmatrix}$. Note that $v^T$ denotes the transpose of $v$. The largest eigenvalue of $A$ is \_\_\_\_\_\_\_\_.

\par\vspace{6pt}\noindent\textbf{Your Answer:} \fbox{\rule{0pt}{1.3ex}\hspace{3cm}}

\tcblower

\textbf{Official Key:} \textbf{\large 3 to 3} \hfill \textit{\small [Easy]}
\vspace{2pt}\hrule height 0.3pt\vspace{3pt}
\textbf{Topic:} Linear Algebra \\
\textbf{Subtopic:} \textit{Eigenvalues and eigenvectors}
\vspace{2pt}\hrule height 0.3pt\vspace{3pt}
\textbf{Guide / Hint:} \\
{\small For an outer product matrix $A = u v^T$ where $u = \begin{bmatrix} 1 \\ 2 \end{bmatrix}$ and $v = \begin{bmatrix} 1 \\ 1 \end{bmatrix}$, the rank of $A$ is 1. Its non-zero eigenvalue equals the trace of $A$, which is given by the inner product $v^T u = 1(1) + 1(2) = 3$.}

\end{tcolorbox}
\vspace{6pt}\hrule height 0.35pt\vspace{6pt}
\noindent
\begin{tcolorbox}[
    blanker,
    sidebyside,
    sidebyside align=top,
    lefthand width=0.55\textwidth,
    top=4pt, bottom=4pt, left=0pt, right=0pt,
    middle=16pt,
    underlay={
        \draw[dashed, black!80, line width=0.8pt, dash pattern=on 4pt off 3pt]
            ([xshift=0.55\textwidth+8pt]frame.north west) --
            ([xshift=0.55\textwidth+8pt]frame.south west);
    }
]
\textbf{Q.44} \quad {\footnotesize\textsc{[MCQ $\cdot$ 1 Mark(s) $\cdot$ Neg: 0.33]}} \hfill {\footnotesize\textbf{GATE 2022 CS} \quad \tiny\texttt{2022\_CS\_GA\_Q02}}
\vspace{2pt}\hrule height 0.3pt\vspace{4pt}

A function $y(x)$ is defined in the interval $[0, 1]$ on the $x$-axis as
\[ y(x) = \begin{cases} 2 & \text{if } 0 \le x < \frac{1}{3} \\ 3 & \text{if } \frac{1}{3} \le x < \frac{3}{4} \\ 1 & \text{if } \frac{3}{4} \le x \le 1 \end{cases} \]
Which one of the following is the area under the curve for the interval $[0, 1]$ on the $x$-axis?

\par\vspace{6pt}\noindent
\begin{tabular}{@{}p{0.48\linewidth}@{\hspace{0.04\linewidth}}p{0.48\linewidth}@{}}
$\square$\ \textbf{(A)} $5/6$ & $\square$\ \textbf{(B)} $6/5$ \\
\addlinespace[3pt]
$\square$\ \textbf{(C)} $13/6$ & $\square$\ \textbf{(D)} $6/13$
\end{tabular}

\tcblower

\textbf{Official Key:} \textbf{\large C} \hfill \textit{\small [Easy]}
\vspace{2pt}\hrule height 0.3pt\vspace{3pt}
\textbf{Topic:} Calculus \\
\textbf{Subtopic:} \textit{Integration}
\vspace{2pt}\hrule height 0.3pt\vspace{3pt}
\textbf{Guide / Hint:} \\
{\small The area under the piecewise step function is the sum of the rectangular areas: $\int_0^1 y(x) \, dx = 2\left(\frac{1}{3} - 0\right) + 3\left(\frac{3}{4} - \frac{1}{3}\right) + 1\left(1 - \frac{3}{4}\right) = \frac{2}{3} + 3\left(\frac{5}{12}\right) + \frac{1}{4} = \frac{2}{3} + \frac{5}{4} + \frac{1}{4} = \frac{2}{3} + \frac{3}{2} = \frac{13}{6}$.}

\end{tcolorbox}
\vspace{6pt}\hrule height 0.35pt\vspace{6pt}
\noindent
\begin{tcolorbox}[
    blanker,
    sidebyside,
    sidebyside align=top,
    lefthand width=0.55\textwidth,
    top=4pt, bottom=4pt, left=0pt, right=0pt,
    middle=16pt,
    underlay={
        \draw[dashed, black!80, line width=0.8pt, dash pattern=on 4pt off 3pt]
            ([xshift=0.55\textwidth+8pt]frame.north west) --
            ([xshift=0.55\textwidth+8pt]frame.south west);
    }
]
\textbf{Q.45} \quad {\footnotesize\textsc{[MCQ $\cdot$ 1 Mark(s) $\cdot$ Neg: 0.33]}} \hfill {\footnotesize\textbf{GATE 2022 CS} \quad \tiny\texttt{2022\_CS\_CS\_Q20}}
\vspace{2pt}\hrule height 0.3pt\vspace{4pt}

Consider the following two statements with respect to the matrices $A_{m \times n}, B_{n \times m}, C_{n \times n}$ and $D_{n \times n}$:
\begin{itemize}[itemsep=0pt,topsep=2pt]
\item Statement 1: $\text{tr}(AB) = \text{tr}(BA)$
\item Statement 2: $\text{tr}(CD) = \text{tr}(DC)$
\end{itemize}
where $\text{tr}()$ represents the trace of a matrix. Which one of the following holds?

\par\vspace{6pt}\noindent
\begin{tabular}{@{}p{0.98\linewidth}@{}}
$\square$\ \textbf{(A)} Statement 1 is correct and Statement 2 is wrong. \\
\addlinespace[3pt]
$\square$\ \textbf{(B)} Statement 1 is wrong and Statement 2 is correct. \\
\addlinespace[3pt]
$\square$\ \textbf{(C)} Both Statement 1 and Statement 2 are correct. \\
\addlinespace[3pt]
$\square$\ \textbf{(D)} Both Statement 1 and Statement 2 are wrong.
\end{tabular}

\tcblower

\textbf{Official Key:} \textbf{\large C} \hfill \textit{\small [Easy]}
\vspace{2pt}\hrule height 0.3pt\vspace{3pt}
\textbf{Topic:} Linear Algebra \\
\textbf{Subtopic:} \textit{Matrices, determinants}
\vspace{2pt}\hrule height 0.3pt\vspace{3pt}
\textbf{Guide / Hint:} \\
{\small The trace is cyclic: $\text{tr}(AB) = \text{tr}(BA)$ holds whenever $A$ is $m \times n$ and $B$ is $n \times m$. Similarly, for $n \times n$ square matrices $C$ and $D$, $\text{tr}(CD) = \text{tr}(DC)$. Thus, both Statement 1 and Statement 2 are correct.}

\end{tcolorbox}
\vspace{6pt}\hrule height 0.35pt\vspace{6pt}
\noindent
\begin{tcolorbox}[
    blanker,
    sidebyside,
    sidebyside align=top,
    lefthand width=0.55\textwidth,
    top=4pt, bottom=4pt, left=0pt, right=0pt,
    middle=16pt,
    underlay={
        \draw[dashed, black!80, line width=0.8pt, dash pattern=on 4pt off 3pt]
            ([xshift=0.55\textwidth+8pt]frame.north west) --
            ([xshift=0.55\textwidth+8pt]frame.south west);
    }
]
\textbf{Q.46} \quad {\footnotesize\textsc{[NAT $\cdot$ 1 Mark(s) $\cdot$ No Neg]}} \hfill {\footnotesize\textbf{GATE 2022 CS} \quad \tiny\texttt{2022\_CS\_CS\_Q34}}
\vspace{2pt}\hrule height 0.3pt\vspace{4pt}

The value of the following limit is \_\_\_\_\_\_\_\_\_\_\_\_\. (Answer in numeric format)
\[ \lim\_{x \to 0^+} \frac{x}{1 - e^{2x}} \]

\par\vspace{6pt}\noindent\textbf{Your Answer:} \fbox{\rule{0pt}{1.3ex}\hspace{3cm}}

\tcblower

\textbf{Official Key:} \textbf{\large -0.5 to -0.5} \hfill \textit{\small [Easy]}
\vspace{2pt}\hrule height 0.3pt\vspace{3pt}
\textbf{Topic:} Calculus \\
\textbf{Subtopic:} \textit{Limits, continuity and differentiability}
\vspace{2pt}\hrule height 0.3pt\vspace{3pt}
\textbf{Guide / Hint:} \\
{\small The limit $\lim_{x \to 0} \frac{x}{1 - e^{2x}}$ has indeterminate form $\frac{0}{0}$. Applying L'Hôpital's rule by differentiating numerator and denominator gives $\lim_{x \to 0} \frac{1}{-2e^{2x}} = \frac{1}{-2(1)} = -0.5$.}

\end{tcolorbox}
\vspace{6pt}\hrule height 0.35pt\vspace{6pt}
\noindent
\begin{tcolorbox}[
    blanker,
    sidebyside,
    sidebyside align=top,
    lefthand width=0.55\textwidth,
    top=4pt, bottom=4pt, left=0pt, right=0pt,
    middle=16pt,
    underlay={
        \draw[dashed, black!80, line width=0.8pt, dash pattern=on 4pt off 3pt]
            ([xshift=0.55\textwidth+8pt]frame.north west) --
            ([xshift=0.55\textwidth+8pt]frame.south west);
    }
]
\textbf{Q.47} \quad {\footnotesize\textsc{[MCQ $\cdot$ 1 Mark(s) $\cdot$ Neg: 0.33]}} \hfill {\footnotesize\textbf{GATE 2023 CS} \quad \tiny\texttt{2023\_CS\_CS\_Q18}}
\vspace{2pt}\hrule height 0.3pt\vspace{4pt}

Let
\[ A = \begin{bmatrix} 1 & 2 & 3 & 4 \\ 4 & 1 & 2 & 3 \\ 3 & 4 & 1 & 2 \\ 2 & 3 & 4 & 1 \end{bmatrix} \quad \text{and} \quad B = \begin{bmatrix} 3 & 4 & 1 & 2 \\ 4 & 1 & 2 & 3 \\ 1 & 2 & 3 & 4 \\ 2 & 3 & 4 & 1 \end{bmatrix}. \]
Let $\det(A)$ and $\det(B)$ denote the determinants of the matrices $A$ and $B$, respectively.
Which one of the options given below is TRUE?

\par\vspace{6pt}\noindent
\begin{tabular}{@{}p{0.48\linewidth}@{\hspace{0.04\linewidth}}p{0.48\linewidth}@{}}
$\square$\ \textbf{(A)} $\det(A) = \det(B)$ & $\square$\ \textbf{(B)} $\det(B) = -\det(A)$ \\
\addlinespace[3pt]
$\square$\ \textbf{(C)} $\det(A) = 0$ & $\square$\ \textbf{(D)} $\det(AB) = \det(A) + \det(B)$
\end{tabular}

\tcblower

\textbf{Official Key:} \textbf{\large B} \hfill \textit{\small [Easy]}
\vspace{2pt}\hrule height 0.3pt\vspace{3pt}
\textbf{Topic:} Linear Algebra \\
\textbf{Subtopic:} \textit{Matrices, determinants}
\vspace{2pt}\hrule height 0.3pt\vspace{3pt}
\textbf{Guide / Hint:} \\
{\small Matrix $B$ is obtained from matrix $A$ by interchanging row 1 and row 3 ($R_1 \leftrightarrow R_3$), leaving all other rows unchanged. By the elementary row operation property of determinants, swapping any two rows of a square matrix negates its determinant: $\det(B) = -\det(A)$.}

\end{tcolorbox}
\vspace{6pt}\hrule height 0.35pt\vspace{6pt}
\noindent
\begin{tcolorbox}[
    blanker,
    sidebyside,
    sidebyside align=top,
    lefthand width=0.55\textwidth,
    top=4pt, bottom=4pt, left=0pt, right=0pt,
    middle=16pt,
    underlay={
        \draw[dashed, black!80, line width=0.8pt, dash pattern=on 4pt off 3pt]
            ([xshift=0.55\textwidth+8pt]frame.north west) --
            ([xshift=0.55\textwidth+8pt]frame.south west);
    }
]
\textbf{Q.48} \quad {\footnotesize\textsc{[MSQ $\cdot$ 1 Mark(s) $\cdot$ No Neg]}} \hfill {\footnotesize\textbf{GATE 2023 CS} \quad \tiny\texttt{2023\_CS\_CS\_Q28}}
\vspace{2pt}\hrule height 0.3pt\vspace{4pt}

Let
\[ f(x) = x^3 + 15x^2 - 33x - 36 \]
be a real-valued function.
Which of the following statements is/are TRUE?

\par\vspace{6pt}\noindent
\begin{tabular}{@{}p{0.98\linewidth}@{}}
$\square$\ \textbf{(A)} $f(x)$ does not have a local maximum. \\
\addlinespace[3pt]
$\square$\ \textbf{(B)} $f(x)$ has a local maximum. \\
\addlinespace[3pt]
$\square$\ \textbf{(C)} $f(x)$ does not have a local minimum. \\
\addlinespace[3pt]
$\square$\ \textbf{(D)} $f(x)$ has a local minimum.
\end{tabular}

\tcblower

\textbf{Official Key:} \textbf{\large B, D} \hfill \textit{\small [Easy]}
\vspace{2pt}\hrule height 0.3pt\vspace{3pt}
\textbf{Topic:} Calculus \\
\textbf{Subtopic:} \textit{Maxima and minima}
\vspace{2pt}\hrule height 0.3pt\vspace{3pt}
\textbf{Guide / Hint:} \\
{\small Setting $f'(x) = 3x^2 + 30x - 33 = 3(x+11)(x-1) = 0$ yields critical points $x = -11$ and $x = 1$. The second derivative is $f''(x) = 6x + 30$. Since $f''(-11) = -36 < 0$, $x = -11$ is a local maximum (B is TRUE). Since $f''(1) = 36 > 0$, $x = 1$ is a local minimum (D is TRUE).}

\end{tcolorbox}
\vspace{6pt}\hrule height 0.35pt\vspace{6pt}
\noindent
\begin{tcolorbox}[
    blanker,
    sidebyside,
    sidebyside align=top,
    lefthand width=0.55\textwidth,
    top=4pt, bottom=4pt, left=0pt, right=0pt,
    middle=16pt,
    underlay={
        \draw[dashed, black!80, line width=0.8pt, dash pattern=on 4pt off 3pt]
            ([xshift=0.55\textwidth+8pt]frame.north west) --
            ([xshift=0.55\textwidth+8pt]frame.south west);
    }
]
\textbf{Q.49} \quad {\footnotesize\textsc{[NAT $\cdot$ 1 Mark(s) $\cdot$ No Neg]}} \hfill {\footnotesize\textbf{GATE 2023 CS} \quad \tiny\texttt{2023\_CS\_CS\_Q31}}
\vspace{2pt}\hrule height 0.3pt\vspace{4pt}

The value of the definite integral
\[ \int\_{-3}^3 \int\_{-2}^2 \int\_{-1}^1 (4x^2 y - z^3) \, dz \, dy \, dx \]
is \_\_\_\_\_\_\_\_\. (Rounded off to the nearest integer)

\par\vspace{6pt}\noindent\textbf{Your Answer:} \fbox{\rule{0pt}{1.3ex}\hspace{3cm}}

\tcblower

\textbf{Official Key:} \textbf{\large 0 to 0} \hfill \textit{\small [Easy]}
\vspace{2pt}\hrule height 0.3pt\vspace{3pt}
\textbf{Topic:} Calculus \\
\textbf{Subtopic:} \textit{Integration}
\vspace{2pt}\hrule height 0.3pt\vspace{3pt}
\textbf{Guide / Hint:} \\
{\small The integral splits into $\int_{-3}^3 \int_{-2}^2 \int_{-1}^1 4x^2 y \, dz \, dy \, dx - \int_{-3}^3 \int_{-2}^2 \int_{-1}^1 z^3 \, dz \, dy \, dx$. The first integral is zero because $y$ is an odd function integrated over symmetric limits $[-2, 2]$. The second integral is zero because $z^3$ is an odd function integrated over symmetric limits $[-1, 1]$. Thus, the total value is $0 - 0 = 0$.}

\end{tcolorbox}
\vspace{6pt}\hrule height 0.35pt\vspace{6pt}
\noindent
\begin{tcolorbox}[
    blanker,
    sidebyside,
    sidebyside align=top,
    lefthand width=0.55\textwidth,
    top=4pt, bottom=4pt, left=0pt, right=0pt,
    middle=16pt,
    underlay={
        \draw[dashed, black!80, line width=0.8pt, dash pattern=on 4pt off 3pt]
            ([xshift=0.55\textwidth+8pt]frame.north west) --
            ([xshift=0.55\textwidth+8pt]frame.south west);
    }
]
\textbf{Q.50} \quad {\footnotesize\textsc{[MSQ $\cdot$ 2 Mark(s) $\cdot$ No Neg]}} \hfill {\footnotesize\textbf{GATE 2023 CS} \quad \tiny\texttt{2023\_CS\_CS\_Q53}}
\vspace{2pt}\hrule height 0.3pt\vspace{4pt}

Consider a random experiment where two fair coins are tossed. Let $A$ be the event that denotes HEAD on both the throws, $B$ be the event that denotes HEAD on the first throw, and $C$ be the event that denotes HEAD on the second throw.
Which of the following statements is/are TRUE?

\par\vspace{6pt}\noindent
\begin{tabular}{@{}p{0.98\linewidth}@{}}
$\square$\ \textbf{(A)} $A$ and $B$ are independent. \\
\addlinespace[3pt]
$\square$\ \textbf{(B)} $A$ and $C$ are independent. \\
\addlinespace[3pt]
$\square$\ \textbf{(C)} $B$ and $C$ are independent. \\
\addlinespace[3pt]
$\square$\ \textbf{(D)} $\text{Prob}(B|C) = \text{Prob}(B)$
\end{tabular}

\tcblower

\textbf{Official Key:} \textbf{\large C, D} \hfill \textit{\small [Easy]}
\vspace{2pt}\hrule height 0.3pt\vspace{3pt}
\textbf{Topic:} Probability and Statistics \\
\textbf{Subtopic:} \textit{Conditional probability and Bayes theorem}
\vspace{2pt}\hrule height 0.3pt\vspace{3pt}
\textbf{Guide / Hint:} \\
{\small Sample space $S = \{HH, HT, TH, TT\}$. Events are $A = \{HH\}$, $B = \{HH, HT\}$, $C = \{HH, TH\}$. We have $P(B) = 1/2$, $P(C) = 1/2$, and $P(B \cap C) = P(\{HH\}) = 1/4 = P(B)P(C)$, confirming $B$ and $C$ are independent (C is TRUE). By independence, $\text{Prob}(B|C) = \text{Prob}(B)$ (D is TRUE). $P(A \cap B) = 1/4 \neq P(A)P(B) = 1/8$, so $A$ and $B$ are dependent.}

\end{tcolorbox}
\vspace{6pt}\hrule height 0.35pt\vspace{6pt}
\noindent
\begin{tcolorbox}[
    blanker,
    sidebyside,
    sidebyside align=top,
    lefthand width=0.55\textwidth,
    top=4pt, bottom=4pt, left=0pt, right=0pt,
    middle=16pt,
    underlay={
        \draw[dashed, black!80, line width=0.8pt, dash pattern=on 4pt off 3pt]
            ([xshift=0.55\textwidth+8pt]frame.north west) --
            ([xshift=0.55\textwidth+8pt]frame.south west);
    }
]
\textbf{Q.51} \quad {\footnotesize\textsc{[MCQ $\cdot$ 1 Mark(s) $\cdot$ Neg: 0.33]}} \hfill {\footnotesize\textbf{GATE 2024 CS1} \quad \tiny\texttt{2024\_CS1\_CS\_Q12}}
\vspace{2pt}\hrule height 0.3pt\vspace{4pt}

The product of all eigenvalues of the matrix $\begin{bmatrix} 1 & 2 & 3 \\ 4 & 5 & 6 \\ 7 & 8 & 9 \end{bmatrix}$ is

\par\vspace{6pt}\noindent
\begin{tabular}{@{}p{0.48\linewidth}@{\hspace{0.04\linewidth}}p{0.48\linewidth}@{}}
$\square$\ \textbf{(A)} $-1$ & $\square$\ \textbf{(B)} 0 \\
\addlinespace[3pt]
$\square$\ \textbf{(C)} 1 & $\square$\ \textbf{(D)} 2
\end{tabular}

\tcblower

\textbf{Official Key:} \textbf{\large B} \hfill \textit{\small [Easy]}
\vspace{2pt}\hrule height 0.3pt\vspace{3pt}
\textbf{Topic:} Linear Algebra \\
\textbf{Subtopic:} \textit{Eigenvalues and eigenvectors}
\vspace{2pt}\hrule height 0.3pt\vspace{3pt}
\textbf{Guide / Hint:} \\
{\small The product of all eigenvalues of a square matrix equals its determinant $\det(A)$. In $A = \begin{bmatrix} 1 & 2 & 3 \\ 4 & 5 & 6 \\ 7 & 8 & 9 \end{bmatrix}$, the rows are in arithmetic progression with $R_3 - R_2 = R_2 - R_1 = [3, 3, 3]$, making the rows linearly dependent and $\det(A) = 0$. Hence, the product of the eigenvalues is 0.}

\end{tcolorbox}
\vspace{6pt}\hrule height 0.35pt\vspace{6pt}
\noindent
\begin{tcolorbox}[
    blanker,
    sidebyside,
    sidebyside align=top,
    lefthand width=0.55\textwidth,
    top=4pt, bottom=4pt, left=0pt, right=0pt,
    middle=16pt,
    underlay={
        \draw[dashed, black!80, line width=0.8pt, dash pattern=on 4pt off 3pt]
            ([xshift=0.55\textwidth+8pt]frame.north west) --
            ([xshift=0.55\textwidth+8pt]frame.south west);
    }
]
\textbf{Q.52} \quad {\footnotesize\textsc{[MSQ $\cdot$ 1 Mark(s) $\cdot$ No Neg]}} \hfill {\footnotesize\textbf{GATE 2024 CS1} \quad \tiny\texttt{2024\_CS1\_CS\_Q27}}
\vspace{2pt}\hrule height 0.3pt\vspace{4pt}

Let $A$ and $B$ be two events in a probability space with $P(A) = 0.3$, $P(B) = 0.5$, and $P(A \cap B) = 0.1$. Which of the following statements is/are TRUE?

\par\vspace{6pt}\noindent
\begin{tabular}{@{}p{0.98\linewidth}@{}}
$\square$\ \textbf{(A)} The two events $A$ and $B$ are independent \\
\addlinespace[3pt]
$\square$\ \textbf{(B)} $P(A \cup B) = 0.7$ \\
\addlinespace[3pt]
$\square$\ \textbf{(C)} $P(A \cap B^c) = 0.2$, where $B^c$ is the complement of the event $B$ \\
\addlinespace[3pt]
$\square$\ \textbf{(D)} $P(A^c \cap B^c) = 0.4$, where $A^c$ and $B^c$ are the complements of the events $A$ and $B$, respectively
\end{tabular}

\tcblower

\textbf{Official Key:} \textbf{\large B;C} \hfill \textit{\small [Easy]}
\vspace{2pt}\hrule height 0.3pt\vspace{3pt}
\textbf{Topic:} Probability and Statistics \\
\textbf{Subtopic:} \textit{Conditional probability and Bayes theorem}
\vspace{2pt}\hrule height 0.3pt\vspace{3pt}
\textbf{Guide / Hint:} \\
{\small We are given $P(A) = 0.3$, $P(B) = 0.5$, and $P(A \cap B) = 0.1$. First, $P(A \cup B) = P(A) + P(B) - P(A \cap B) = 0.3 + 0.5 - 0.1 = 0.7$ (B is TRUE). Second, $P(A \cap B^c) = P(A) - P(A \cap B) = 0.3 - 0.1 = 0.2$ (C is TRUE). Since $P(A)P(B) = 0.15 \neq 0.1$, the events are not independent.}

\end{tcolorbox}
\vspace{6pt}\hrule height 0.35pt\vspace{6pt}
\noindent
\begin{tcolorbox}[
    blanker,
    sidebyside,
    sidebyside align=top,
    lefthand width=0.55\textwidth,
    top=4pt, bottom=4pt, left=0pt, right=0pt,
    middle=16pt,
    underlay={
        \draw[dashed, black!80, line width=0.8pt, dash pattern=on 4pt off 3pt]
            ([xshift=0.55\textwidth+8pt]frame.north west) --
            ([xshift=0.55\textwidth+8pt]frame.south west);
    }
]
\textbf{Q.53} \quad {\footnotesize\textsc{[NAT $\cdot$ 2 Mark(s) $\cdot$ No Neg]}} \hfill {\footnotesize\textbf{GATE 2024 CS1} \quad \tiny\texttt{2024\_CS1\_CS\_Q63}}
\vspace{2pt}\hrule height 0.3pt\vspace{4pt}

A bag contains 10 red balls and 15 blue balls. Two balls are drawn randomly without replacement. Given that the first ball drawn is red, the probability (rounded off to 3 decimal places) that both balls drawn are red is \_\_\_\_\_\_\_\_\_\. (Answer in numeric format)

\par\vspace{6pt}\noindent\textbf{Your Answer:} \fbox{\rule{0pt}{1.3ex}\hspace{3cm}}

\tcblower

\textbf{Official Key:} \textbf{\large 0.370 to 0.380} \hfill \textit{\small [Easy]}
\vspace{2pt}\hrule height 0.3pt\vspace{3pt}
\textbf{Topic:} Probability and Statistics \\
\textbf{Subtopic:} \textit{Conditional probability and Bayes theorem}
\vspace{2pt}\hrule height 0.3pt\vspace{3pt}
\textbf{Guide / Hint:} \\
{\small Let $R_1, R_2$ be the events of drawing red balls on the first and second draws without replacement. The conditional probability that both are red given the first is red is $P(R_1 \cap R_2 \mid R_1) = P(R_2 \mid R_1)$. Given that 1 red ball is already drawn, 9 red balls remain out of $9 + 15 = 24$ total balls, so $P(R_2 \mid R_1) = \frac{9}{24} = \frac{3}{8} = 0.375$.}

\end{tcolorbox}
\vspace{6pt}\hrule height 0.35pt\vspace{6pt}
\noindent
\begin{tcolorbox}[
    blanker,
    sidebyside,
    sidebyside align=top,
    lefthand width=0.55\textwidth,
    top=4pt, bottom=4pt, left=0pt, right=0pt,
    middle=16pt,
    underlay={
        \draw[dashed, black!80, line width=0.8pt, dash pattern=on 4pt off 3pt]
            ([xshift=0.55\textwidth+8pt]frame.north west) --
            ([xshift=0.55\textwidth+8pt]frame.south west);
    }
]
\textbf{Q.54} \quad {\footnotesize\textsc{[MCQ $\cdot$ 1 Mark(s) $\cdot$ Neg: 0.33]}} \hfill {\footnotesize\textbf{GATE 2024 CS2} \quad \tiny\texttt{2024\_CS2\_CS\_Q18}}
\vspace{2pt}\hrule height 0.3pt\vspace{4pt}

When six unbiased dice are rolled simultaneously, the probability of getting all distinct numbers (i.e., 1, 2, 3, 4, 5, and 6) is:

\par\vspace{6pt}\noindent
\begin{tabular}{@{}p{0.48\linewidth}@{\hspace{0.04\linewidth}}p{0.48\linewidth}@{}}
$\square$\ \textbf{(A)} $\frac{1}{324}$ & $\square$\ \textbf{(B)} $\frac{5}{324}$ \\
\addlinespace[3pt]
$\square$\ \textbf{(C)} $\frac{7}{324}$ & $\square$\ \textbf{(D)} $\frac{11}{324}$
\end{tabular}

\tcblower

\textbf{Official Key:} \textbf{\large B} \hfill \textit{\small [Easy]}
\vspace{2pt}\hrule height 0.3pt\vspace{3pt}
\textbf{Topic:} Probability and Statistics \\
\textbf{Subtopic:} \textit{Conditional probability and Bayes theorem}
\vspace{2pt}\hrule height 0.3pt\vspace{3pt}
\textbf{Guide / Hint:} \\
{\small When six unbiased 6-faced dice are rolled, total outcomes are $6^6$. The number of outcomes where all 6 faces are distinct is $6! = 720$. The probability is $\frac{6!}{6^6} = \frac{720}{46656} = \frac{5}{324}$.}

\end{tcolorbox}
\vspace{6pt}\hrule height 0.35pt\vspace{6pt}
\noindent
\begin{tcolorbox}[
    blanker,
    sidebyside,
    sidebyside align=top,
    lefthand width=0.55\textwidth,
    top=4pt, bottom=4pt, left=0pt, right=0pt,
    middle=16pt,
    underlay={
        \draw[dashed, black!80, line width=0.8pt, dash pattern=on 4pt off 3pt]
            ([xshift=0.55\textwidth+8pt]frame.north west) --
            ([xshift=0.55\textwidth+8pt]frame.south west);
    }
]
\textbf{Q.55} \quad {\footnotesize\textsc{[MSQ $\cdot$ 2 Mark(s) $\cdot$ No Neg]}} \hfill {\footnotesize\textbf{GATE 2024 CS2} \quad \tiny\texttt{2024\_CS2\_CS\_Q47}}
\vspace{2pt}\hrule height 0.3pt\vspace{4pt}

Let $A$ be an $n \times n$ matrix over the set of all real numbers $\mathbb{R}$. Let $B$ be a matrix obtained from $A$ by swapping two rows. Which of the following statements is/are TRUE?

\par\vspace{6pt}\noindent
\begin{tabular}{@{}p{0.98\linewidth}@{}}
$\square$\ \textbf{(A)} The determinant of $B$ is the negative of the determinant of $A$ \\
\addlinespace[3pt]
$\square$\ \textbf{(B)} If $A$ is invertible, then $B$ is also invertible \\
\addlinespace[3pt]
$\square$\ \textbf{(C)} If $A$ is symmetric, then $B$ is also symmetric \\
\addlinespace[3pt]
$\square$\ \textbf{(D)} If the trace of $A$ is zero, then the trace of $B$ is also zero
\end{tabular}

\tcblower

\textbf{Official Key:} \textbf{\large A;B} \hfill \textit{\small [Easy]}
\vspace{2pt}\hrule height 0.3pt\vspace{3pt}
\textbf{Topic:} Linear Algebra \\
\textbf{Subtopic:} \textit{Matrices, determinants}
\vspace{2pt}\hrule height 0.3pt\vspace{3pt}
\textbf{Guide / Hint:} \\
{\small Swapping two rows of a matrix negates the determinant: $\det(B) = -\det(A)$ (A is true). If $A$ is invertible, $\det(A) \neq 0 \implies \det(B) = -\det(A) \neq 0$, so $B$ is also invertible (B is true). Row swapping does not preserve symmetry or trace.}

\end{tcolorbox}
\vspace{6pt}\hrule height 0.35pt\vspace{6pt}
\noindent
\begin{tcolorbox}[
    blanker,
    sidebyside,
    sidebyside align=top,
    lefthand width=0.55\textwidth,
    top=4pt, bottom=4pt, left=0pt, right=0pt,
    middle=16pt,
    underlay={
        \draw[dashed, black!80, line width=0.8pt, dash pattern=on 4pt off 3pt]
            ([xshift=0.55\textwidth+8pt]frame.north west) --
            ([xshift=0.55\textwidth+8pt]frame.south west);
    }
]
\textbf{Q.56} \quad {\footnotesize\textsc{[NAT $\cdot$ 1 Mark(s) $\cdot$ No Neg]}} \hfill {\footnotesize\textbf{GATE 2025 CS1} \quad \tiny\texttt{2025\_CS1\_CS\_Q31}}
\vspace{2pt}\hrule height 0.3pt\vspace{4pt}

Consider the given function $f(x)$:
\[ f(x) = \begin{cases} ax + b & \text{for } x < 1 \\ x^3 + x^2 + 1 & \text{for } x \ge 1 \end{cases} \]
If the function is differentiable everywhere, the value of $b$ must be \_\_\_\_\_\_\. (rounded off to one decimal place)

\par\vspace{6pt}\noindent\textbf{Your Answer:} \fbox{\rule{0pt}{1.3ex}\hspace{3cm}}

\tcblower

\textbf{Official Key:} \textbf{\large -2.1 to -1.9} \hfill \textit{\small [Easy]}
\vspace{2pt}\hrule height 0.3pt\vspace{3pt}
\textbf{Topic:} Calculus \\
\textbf{Subtopic:} \textit{Limits, continuity and differentiability}
\vspace{2pt}\hrule height 0.3pt\vspace{3pt}
\textbf{Guide / Hint:} \\
{\small For differentiability at $x = 1$, left and right derivatives must be equal: $\lim_{x \to 1^-} f'(x) = a$ and $\lim_{x \to 1^+} f'(x) = 3(1)^2 + 2(1) = 5$, so $a = 5$. Continuity at $x = 1$ requires $a(1) + b = 1^3 + 1^2 + 1 = 3$. Substituting $a = 5$ gives $5 + b = 3 \implies b = -2.0$.}

\end{tcolorbox}
\vspace{6pt}\hrule height 0.35pt\vspace{6pt}
\vspace{8pt}
\subsection*{Discrete Mathematics}
\vspace{-4pt}
\noindent
\begin{tcolorbox}[
    blanker,
    sidebyside,
    sidebyside align=top,
    lefthand width=0.55\textwidth,
    top=4pt, bottom=4pt, left=0pt, right=0pt,
    middle=16pt,
    underlay={
        \draw[dashed, black!80, line width=0.8pt, dash pattern=on 4pt off 3pt]
            ([xshift=0.55\textwidth+8pt]frame.north west) --
            ([xshift=0.55\textwidth+8pt]frame.south west);
    }
]
\textbf{Q.57} \quad {\footnotesize\textsc{[MCQ $\cdot$ 1 Mark(s) $\cdot$ Neg: 0.33]}} \hfill {\footnotesize\textbf{GATE 2018 CS} \quad \tiny\texttt{2018\_CS\_GA\_Q01\_v2}}
\vspace{2pt}\hrule height 0.3pt\vspace{4pt}

Which one of the following is a closed form expression for the generating function of the sequence $\{a_n\}$, where $a_n = 2n + 3$ for all $n = 0, 1, 2, \dots$?

\par\vspace{6pt}\noindent
\begin{tabular}{@{}p{0.48\linewidth}@{\hspace{0.04\linewidth}}p{0.48\linewidth}@{}}
$\square$\ \textbf{(A)} $\frac{3}{(1-x)^2}$ & $\square$\ \textbf{(B)} $\frac{3x}{(1-x)^2}$ \\
\addlinespace[3pt]
$\square$\ \textbf{(C)} $\frac{2-x}{(1-x)^2}$ & $\square$\ \textbf{(D)} $\frac{3-x}{(1-x)^2}$
\end{tabular}

\tcblower

\textbf{Official Key:} \textbf{\large D} \hfill \textit{\small [Easy]}
\vspace{2pt}\hrule height 0.3pt\vspace{3pt}
\textbf{Topic:} Discrete Mathematics \\
\textbf{Subtopic:} \textit{Combinatorics: counting, recurrence relations, generating functions}
\vspace{2pt}\hrule height 0.3pt\vspace{3pt}
\textbf{Guide / Hint:} \\
{\small The generating function is $G(x) = \sum_{n=0}^\infty (2n + 3)x^n = 2\sum_{n=0}^\infty n x^n + 3\sum_{n=0}^\infty x^n$. Using the standard power series $\sum_{n=0}^\infty x^n = \frac{1}{1-x}$ and $\sum_{n=0}^\infty n x^n = \frac{x}{(1-x)^2}$, we obtain $G(x) = \frac{2x}{(1-x)^2} + \frac{3(1-x)}{(1-x)^2} = \frac{3 - x}{(1-x)^2}$.}

\end{tcolorbox}
\vspace{6pt}\hrule height 0.35pt\vspace{6pt}
\noindent
\begin{tcolorbox}[
    blanker,
    sidebyside,
    sidebyside align=top,
    lefthand width=0.55\textwidth,
    top=4pt, bottom=4pt, left=0pt, right=0pt,
    middle=16pt,
    underlay={
        \draw[dashed, black!80, line width=0.8pt, dash pattern=on 4pt off 3pt]
            ([xshift=0.55\textwidth+8pt]frame.north west) --
            ([xshift=0.55\textwidth+8pt]frame.south west);
    }
]
\textbf{Q.58} \quad {\footnotesize\textsc{[NAT $\cdot$ 1 Mark(s) $\cdot$ No Neg]}} \hfill {\footnotesize\textbf{GATE 2018 CS} \quad \tiny\texttt{2018\_CS\_GA\_Q18}}
\vspace{2pt}\hrule height 0.3pt\vspace{4pt}

The chromatic number of the graph containing a 7-cycle with a single cross-chord is \_\_\_\_\_\_\_\_.

\par\vspace{6pt}\noindent\textbf{Your Answer:} \fbox{\rule{0pt}{1.3ex}\hspace{3cm}}

\tcblower

\textbf{Official Key:} \textbf{\large 3 to 3} \hfill \textit{\small [Easy]}
\vspace{2pt}\hrule height 0.3pt\vspace{3pt}
\textbf{Topic:} Discrete Mathematics \\
\textbf{Subtopic:} \textit{Graphs: connectivity, matching, colouring}
\vspace{2pt}\hrule height 0.3pt\vspace{3pt}
\textbf{Guide / Hint:} \\
{\small The given graph contains an odd cycle, which prevents it from being 2-colorable (bipartite), meaning $\chi(G) \ge 3$. A valid 3-coloring can be directly constructed without any conflicts between adjacent vertices, so $\chi(G) = 3$.}

\end{tcolorbox}
\vspace{6pt}\hrule height 0.35pt\vspace{6pt}
\noindent
\begin{tcolorbox}[
    blanker,
    sidebyside,
    sidebyside align=top,
    lefthand width=0.55\textwidth,
    top=4pt, bottom=4pt, left=0pt, right=0pt,
    middle=16pt,
    underlay={
        \draw[dashed, black!80, line width=0.8pt, dash pattern=on 4pt off 3pt]
            ([xshift=0.55\textwidth+8pt]frame.north west) --
            ([xshift=0.55\textwidth+8pt]frame.south west);
    }
]
\textbf{Q.59} \quad {\footnotesize\textsc{[NAT $\cdot$ 1 Mark(s) $\cdot$ No Neg]}} \hfill {\footnotesize\textbf{GATE 2018 CS} \quad \tiny\texttt{2018\_CS\_GA\_Q19}}
\vspace{2pt}\hrule height 0.3pt\vspace{4pt}

Let $G$ be a finite group on 84 elements. The size of a largest possible proper subgroup of $G$ is \_\_\_\_\_\_\_\_.

\par\vspace{6pt}\noindent\textbf{Your Answer:} \fbox{\rule{0pt}{1.3ex}\hspace{3cm}}

\tcblower

\textbf{Official Key:} \textbf{\large 42 to 42} \hfill \textit{\small [Easy]}
\vspace{2pt}\hrule height 0.3pt\vspace{3pt}
\textbf{Topic:} Discrete Mathematics \\
\textbf{Subtopic:} \textit{Monoids, Groups}
\vspace{2pt}\hrule height 0.3pt\vspace{3pt}
\textbf{Guide / Hint:} \\
{\small By Lagrange's Theorem, the order of any subgroup $H$ of a finite group $G$ must divide $|G| = 84$. A proper subgroup cannot equal the whole group, so $|H| < 84$. The largest proper divisor of 84 is $\frac{84}{2} = 42$.}

\end{tcolorbox}
\vspace{6pt}\hrule height 0.35pt\vspace{6pt}
\noindent
\begin{tcolorbox}[
    blanker,
    sidebyside,
    sidebyside align=top,
    lefthand width=0.55\textwidth,
    top=4pt, bottom=4pt, left=0pt, right=0pt,
    middle=16pt,
    underlay={
        \draw[dashed, black!80, line width=0.8pt, dash pattern=on 4pt off 3pt]
            ([xshift=0.55\textwidth+8pt]frame.north west) --
            ([xshift=0.55\textwidth+8pt]frame.south west);
    }
]
\textbf{Q.60} \quad {\footnotesize\textsc{[NAT $\cdot$ 1 Mark(s) $\cdot$ No Neg]}} \hfill {\footnotesize\textbf{GATE 2022 CS} \quad \tiny\texttt{2022\_CS\_CS\_Q30}}
\vspace{2pt}\hrule height 0.3pt\vspace{4pt}

Consider a simple undirected graph of 10 vertices. If the graph is disconnected, then the maximum number of edges it can have is \_\_\_\_\_\_\_\_\_\_\_\. (Answer in integer)

\par\vspace{6pt}\noindent\textbf{Your Answer:} \fbox{\rule{0pt}{1.3ex}\hspace{3cm}}

\tcblower

\textbf{Official Key:} \textbf{\large 36 to 36} \hfill \textit{\small [Easy]}
\vspace{2pt}\hrule height 0.3pt\vspace{3pt}
\textbf{Topic:} Discrete Mathematics \\
\textbf{Subtopic:} \textit{Graphs: connectivity, matching, colouring}
\vspace{2pt}\hrule height 0.3pt\vspace{3pt}
\textbf{Guide / Hint:} \\
{\small A disconnected graph on $n = 10$ vertices with the maximum number of edges has one isolated vertex and one complete component $K_{n-1} = K_9$. The maximum number of edges is $\binom{9}{2} = \frac{9 \times 8}{2} = 36$.}

\end{tcolorbox}
\vspace{6pt}\hrule height 0.35pt\vspace{6pt}
\noindent
\begin{tcolorbox}[
    blanker,
    sidebyside,
    sidebyside align=top,
    lefthand width=0.55\textwidth,
    top=4pt, bottom=4pt, left=0pt, right=0pt,
    middle=16pt,
    underlay={
        \draw[dashed, black!80, line width=0.8pt, dash pattern=on 4pt off 3pt]
            ([xshift=0.55\textwidth+8pt]frame.north west) --
            ([xshift=0.55\textwidth+8pt]frame.south west);
    }
]
\textbf{Q.61} \quad {\footnotesize\textsc{[NAT $\cdot$ 1 Mark(s) $\cdot$ No Neg]}} \hfill {\footnotesize\textbf{GATE 2024 CS1} \quad \tiny\texttt{2024\_CS1\_CS\_Q34}}
\vspace{2pt}\hrule height 0.3pt\vspace{4pt}

The number of spanning trees in a complete graph of 4 vertices labelled A, B, C, and D is \_\_\_\_\_\_\_\_\_\. (Answer in integer)

\par\vspace{6pt}\noindent\textbf{Your Answer:} \fbox{\rule{0pt}{1.3ex}\hspace{3cm}}

\tcblower

\textbf{Official Key:} \textbf{\large 16 to 16} \hfill \textit{\small [Easy]}
\vspace{2pt}\hrule height 0.3pt\vspace{3pt}
\textbf{Topic:} Discrete Mathematics \\
\textbf{Subtopic:} \textit{Combinatorics: counting, recurrence relations, generating functions}
\vspace{2pt}\hrule height 0.3pt\vspace{3pt}
\textbf{Guide / Hint:} \\
{\small By Cayley's theorem, the number of spanning trees in a complete labelled graph $K_n$ on $n$ vertices is given by $n^{n-2}$. For $n = 4$ labelled vertices A, B, C, D, the number of spanning trees is $4^{4-2} = 4^2 = 16$.}

\end{tcolorbox}
\vspace{6pt}\hrule height 0.35pt\vspace{6pt}
\noindent
\begin{tcolorbox}[
    blanker,
    sidebyside,
    sidebyside align=top,
    lefthand width=0.55\textwidth,
    top=4pt, bottom=4pt, left=0pt, right=0pt,
    middle=16pt,
    underlay={
        \draw[dashed, black!80, line width=0.8pt, dash pattern=on 4pt off 3pt]
            ([xshift=0.55\textwidth+8pt]frame.north west) --
            ([xshift=0.55\textwidth+8pt]frame.south west);
    }
]
\textbf{Q.62} \quad {\footnotesize\textsc{[NAT $\cdot$ 2 Mark(s) $\cdot$ No Neg]}} \hfill {\footnotesize\textbf{GATE 2024 CS1} \quad \tiny\texttt{2024\_CS1\_CS\_Q60}}
\vspace{2pt}\hrule height 0.3pt\vspace{4pt}

The number of edges present in the forest generated by the DFS traversal of an undirected graph $G$ with 100 vertices is 40. The number of connected components in $G$ is \_\_\_\_\_\_\_\_\_\. (Answer in integer)

\par\vspace{6pt}\noindent\textbf{Your Answer:} \fbox{\rule{0pt}{1.3ex}\hspace{3cm}}

\tcblower

\textbf{Official Key:} \textbf{\large 60 to 60} \hfill \textit{\small [Easy]}
\vspace{2pt}\hrule height 0.3pt\vspace{3pt}
\textbf{Topic:} Discrete Mathematics \\
\textbf{Subtopic:} \textit{Graphs: connectivity, matching, colouring}
\vspace{2pt}\hrule height 0.3pt\vspace{3pt}
\textbf{Guide / Hint:} \\
{\small A DFS traversal of an undirected graph with $V$ vertices and $k$ connected components produces a spanning forest containing $k$ trees. A tree with $v_i$ vertices contains $v_i - 1$ edges, so the total number of tree edges in the forest is $\sum (v_i - 1) = V - k$. Given $V = 100$ and $E = 40$, we have $40 = 100 - k \implies k = 60$ connected components.}

\end{tcolorbox}
\vspace{6pt}\hrule height 0.35pt\vspace{6pt}
\noindent
\begin{tcolorbox}[
    blanker,
    sidebyside,
    sidebyside align=top,
    lefthand width=0.55\textwidth,
    top=4pt, bottom=4pt, left=0pt, right=0pt,
    middle=16pt,
    underlay={
        \draw[dashed, black!80, line width=0.8pt, dash pattern=on 4pt off 3pt]
            ([xshift=0.55\textwidth+8pt]frame.north west) --
            ([xshift=0.55\textwidth+8pt]frame.south west);
    }
]
\textbf{Q.63} \quad {\footnotesize\textsc{[MCQ $\cdot$ 1 Mark(s) $\cdot$ Neg: 0.33]}} \hfill {\footnotesize\textbf{GATE 2024 CS2} \quad \tiny\texttt{2024\_CS2\_CS\_Q12}}
\vspace{2pt}\hrule height 0.3pt\vspace{4pt}

Let $p$ and $q$ be the following propositions:
$p$: Fail grade can be given.
$q$: Student scores more than 50\% marks.
Consider the statement: ``Fail grade cannot be given when student scores more than 50\% marks.''
Which one of the following is the CORRECT representation of the above statement in propositional logic?

\par\vspace{6pt}\noindent
\begin{tabular}{@{}p{0.48\linewidth}@{\hspace{0.04\linewidth}}p{0.48\linewidth}@{}}
$\square$\ \textbf{(A)} $q \to \neg p$ & $\square$\ \textbf{(B)} $q \to p$ \\
\addlinespace[3pt]
$\square$\ \textbf{(C)} $p \to q$ & $\square$\ \textbf{(D)} $\neg p \to q$
\end{tabular}

\tcblower

\textbf{Official Key:} \textbf{\large A} \hfill \textit{\small [Easy]}
\vspace{2pt}\hrule height 0.3pt\vspace{3pt}
\textbf{Topic:} Discrete Mathematics \\
\textbf{Subtopic:} \textit{Propositional and first order logic}
\vspace{2pt}\hrule height 0.3pt\vspace{3pt}
\textbf{Guide / Hint:} \\
{\small The statement is: 'Fail grade cannot be given ($\neg p$) when student scores more than 50\% marks ($q$)', which means 'If $q$, then $\neg p$'. In propositional logic, this is directly formalized as $q \to \neg p$.}

\end{tcolorbox}
\vspace{6pt}\hrule height 0.35pt\vspace{6pt}
\noindent
\begin{tcolorbox}[
    blanker,
    sidebyside,
    sidebyside align=top,
    lefthand width=0.55\textwidth,
    top=4pt, bottom=4pt, left=0pt, right=0pt,
    middle=16pt,
    underlay={
        \draw[dashed, black!80, line width=0.8pt, dash pattern=on 4pt off 3pt]
            ([xshift=0.55\textwidth+8pt]frame.north west) --
            ([xshift=0.55\textwidth+8pt]frame.south west);
    }
]
\textbf{Q.64} \quad {\footnotesize\textsc{[MCQ $\cdot$ 1 Mark(s) $\cdot$ Neg: 0.33]}} \hfill {\footnotesize\textbf{GATE 2024 CS2} \quad \tiny\texttt{2024\_CS2\_CS\_Q15}}
\vspace{2pt}\hrule height 0.3pt\vspace{4pt}

Let $T(n)$ be the recurrence relation defined as follows:
\[ T(0) = 1, \quad T(1) = 2, \quad \text{and} \quad T(n) = 5T(n-1) - 6T(n-2) \quad \text{for } n \ge 2 \]
Which one of the following statements is TRUE?

\par\vspace{6pt}\noindent
\begin{tabular}{@{}p{0.48\linewidth}@{\hspace{0.04\linewidth}}p{0.48\linewidth}@{}}
$\square$\ \textbf{(A)} $T(n) = \Theta(2^n)$ & $\square$\ \textbf{(B)} $T(n) = \Theta(n 2^n)$ \\
\addlinespace[3pt]
$\square$\ \textbf{(C)} $T(n) = \Theta(3^n)$ & $\square$\ \textbf{(D)} $T(n) = \Theta(n 3^n)$
\end{tabular}

\tcblower

\textbf{Official Key:} \textbf{\large A} \hfill \textit{\small [Easy]}
\vspace{2pt}\hrule height 0.3pt\vspace{3pt}
\textbf{Topic:} Discrete Mathematics \\
\textbf{Subtopic:} \textit{Combinatorics: counting, recurrence relations, generating functions}
\vspace{2pt}\hrule height 0.3pt\vspace{3pt}
\textbf{Guide / Hint:} \\
{\small The characteristic equation is $r^2 - 5r + 6 = 0 \implies (r - 2)(r - 3) = 0$, with roots $r = 2, 3$. The general solution is $T(n) = c_1 2^n + c_2 3^n$. Using initial conditions $T(0) = 1, T(1) = 2$: $c_1 + c_2 = 1$ and $2c_1 + 3c_2 = 2 \implies c_1 = 1, c_2 = 0$. Thus, $T(n) = 2^n = \Theta(2^n)$.}

\end{tcolorbox}
\vspace{6pt}\hrule height 0.35pt\vspace{6pt}
\noindent
\begin{tcolorbox}[
    blanker,
    sidebyside,
    sidebyside align=top,
    lefthand width=0.55\textwidth,
    top=4pt, bottom=4pt, left=0pt, right=0pt,
    middle=16pt,
    underlay={
        \draw[dashed, black!80, line width=0.8pt, dash pattern=on 4pt off 3pt]
            ([xshift=0.55\textwidth+8pt]frame.north west) --
            ([xshift=0.55\textwidth+8pt]frame.south west);
    }
]
\textbf{Q.65} \quad {\footnotesize\textsc{[NAT $\cdot$ 2 Mark(s) $\cdot$ No Neg]}} \hfill {\footnotesize\textbf{GATE 2024 CS2} \quad \tiny\texttt{2024\_CS2\_CS\_Q60}}
\vspace{2pt}\hrule height 0.3pt\vspace{4pt}

The chromatic number of a graph is the minimum number of colours used in a proper colouring of the graph. The chromatic number of the given bipartite graph is \_\_\_\_\_\_\. (Answer in integer)

\par\vspace{6pt}\noindent\textbf{Your Answer:} \fbox{\rule{0pt}{1.3ex}\hspace{3cm}}

\tcblower

\textbf{Official Key:} \textbf{\large 2 to 2} \hfill \textit{\small [Easy]}
\vspace{2pt}\hrule height 0.3pt\vspace{3pt}
\textbf{Topic:} Discrete Mathematics \\
\textbf{Subtopic:} \textit{Graphs: connectivity, matching, colouring}
\vspace{2pt}\hrule height 0.3pt\vspace{3pt}
\textbf{Guide / Hint:} \\
{\small A graph has chromatic number 2 if and only if it is bipartite and contains at least one edge, which is equivalent to having no odd cycles. Inspection of the given graph verifies that it contains no odd-length cycles, meaning it can be properly colored with 2 colors.}

\end{tcolorbox}
\vspace{6pt}\hrule height 0.35pt\vspace{6pt}
\noindent
\begin{tcolorbox}[
    blanker,
    sidebyside,
    sidebyside align=top,
    lefthand width=0.55\textwidth,
    top=4pt, bottom=4pt, left=0pt, right=0pt,
    middle=16pt,
    underlay={
        \draw[dashed, black!80, line width=0.8pt, dash pattern=on 4pt off 3pt]
            ([xshift=0.55\textwidth+8pt]frame.north west) --
            ([xshift=0.55\textwidth+8pt]frame.south west);
    }
]
\textbf{Q.66} \quad {\footnotesize\textsc{[NAT $\cdot$ 2 Mark(s) $\cdot$ No Neg]}} \hfill {\footnotesize\textbf{GATE 2024 CS2} \quad \tiny\texttt{2024\_CS2\_CS\_Q63}}
\vspace{2pt}\hrule height 0.3pt\vspace{4pt}

Let $Z_n$ be the group of integers $\{0, 1, 2, \dots, n-1\}$ with addition modulo $n$ as the group operation. The number of elements in the group $Z_2 \times Z_3 \times Z_4$ that are their own inverses is \_\_\_\_\_\_\. (Answer in integer)

\par\vspace{6pt}\noindent\textbf{Your Answer:} \fbox{\rule{0pt}{1.3ex}\hspace{3cm}}

\tcblower

\textbf{Official Key:} \textbf{\large 4 to 4} \hfill \textit{\small [Easy]}
\vspace{2pt}\hrule height 0.3pt\vspace{3pt}
\textbf{Topic:} Discrete Mathematics \\
\textbf{Subtopic:} \textit{Monoids, Groups}
\vspace{2pt}\hrule height 0.3pt\vspace{3pt}
\textbf{Guide / Hint:} \\
{\small In an additive group, an element $x$ is its own inverse if $x + x \equiv 0 \pmod n \implies 2x \equiv 0 \pmod n$. In $\mathbb{Z}_2$, elements satisfying $2x \equiv 0$ are $\{0, 1\}$ (2 elements). In $\mathbb{Z}_3$, only $\{0\}$ satisfies $2x \equiv 0 \pmod 3$ (1 element). In $\mathbb{Z}_4$, $\{0, 2\}$ satisfy $2x \equiv 0 \pmod 4$ (2 elements). In the direct product group $\mathbb{Z}_2 \times \mathbb{Z}_3 \times \mathbb{Z}_4$, total self-inverse elements $= 2 \times 1 \times 2 = 4$.}

\end{tcolorbox}
\vspace{6pt}\hrule height 0.35pt\vspace{6pt}
\noindent
\begin{tcolorbox}[
    blanker,
    sidebyside,
    sidebyside align=top,
    lefthand width=0.55\textwidth,
    top=4pt, bottom=4pt, left=0pt, right=0pt,
    middle=16pt,
    underlay={
        \draw[dashed, black!80, line width=0.8pt, dash pattern=on 4pt off 3pt]
            ([xshift=0.55\textwidth+8pt]frame.north west) --
            ([xshift=0.55\textwidth+8pt]frame.south west);
    }
]
\textbf{Q.67} \quad {\footnotesize\textsc{[MCQ $\cdot$ 1 Mark(s) $\cdot$ Neg: 0.33]}} \hfill {\footnotesize\textbf{GATE 2025 CS2} \quad \tiny\texttt{2025\_CS2\_CS\_Q15}}
\vspace{2pt}\hrule height 0.3pt\vspace{4pt}

Let $P(x)$ be an arbitrary predicate over the domain of natural numbers. Which ONE of the following statements is TRUE?

\par\vspace{6pt}\noindent
\begin{tabular}{@{}p{0.98\linewidth}@{}}
$\square$\ \textbf{(A)} $(P(0) \land (\forall x [P(x) \implies P(x+1)])) \implies (\forall x P(x))$ \\
\addlinespace[3pt]
$\square$\ \textbf{(B)} $(P(0) \land (\forall x [P(x) \implies P(x-1)])) \implies (\forall x P(x))$ \\
\addlinespace[3pt]
$\square$\ \textbf{(C)} $(P(1000) \land (\forall x [P(x) \implies P(x-1)])) \implies (\forall x P(x))$ \\
\addlinespace[3pt]
$\square$\ \textbf{(D)} $(P(1000) \land (\forall x [P(x) \implies P(x+1)])) \implies (\forall x P(x))$
\end{tabular}

\tcblower

\textbf{Official Key:} \textbf{\large A} \hfill \textit{\small [Easy]}
\vspace{2pt}\hrule height 0.3pt\vspace{3pt}
\textbf{Topic:} Discrete Mathematics \\
\textbf{Subtopic:} \textit{Propositional and first order logic: mathematical induction}
\vspace{2pt}\hrule height 0.3pt\vspace{3pt}
\textbf{Guide / Hint:} \\
{\small Option A is the formal statement of the Principle of Mathematical Induction over $\mathbb{N}$: if the base case $P(0)$ holds and for all $x$, $P(x) \implies P(x+1)$, then $\forall x, P(x)$ is true.}

\end{tcolorbox}
\vspace{6pt}\hrule height 0.35pt\vspace{6pt}
\vspace{8pt}
\subsection*{Digital Logic}
\vspace{-4pt}
\noindent
\begin{tcolorbox}[
    blanker,
    sidebyside,
    sidebyside align=top,
    lefthand width=0.55\textwidth,
    top=4pt, bottom=4pt, left=0pt, right=0pt,
    middle=16pt,
    underlay={
        \draw[dashed, black!80, line width=0.8pt, dash pattern=on 4pt off 3pt]
            ([xshift=0.55\textwidth+8pt]frame.north west) --
            ([xshift=0.55\textwidth+8pt]frame.south west);
    }
]
\textbf{Q.68} \quad {\footnotesize\textsc{[MCQ $\cdot$ 1 Mark(s) $\cdot$ Neg: 0.33]}} \hfill {\footnotesize\textbf{GATE 2018 CS} \quad \tiny\texttt{2018\_CS\_GA\_Q04\_v2}}
\vspace{2pt}\hrule height 0.3pt\vspace{4pt}

Let $\oplus$ and $\odot$ denote the Exclusive OR and Exclusive NOR operations, respectively. Which one of the following is NOT CORRECT?

\par\vspace{6pt}\noindent
\begin{tabular}{@{}p{0.98\linewidth}@{}}
$\square$\ \textbf{(A)} $\overline{P \oplus Q} = P \odot Q$ \\
\addlinespace[3pt]
$\square$\ \textbf{(B)} $\overline{P \odot Q} = P \oplus Q$ \\
\addlinespace[3pt]
$\square$\ \textbf{(C)} $\overline{P} \oplus Q = P \odot Q$ \\
\addlinespace[3pt]
$\square$\ \textbf{(D)} $(P \oplus \overline{P}) \oplus Q = (P \odot \overline{P}) \odot \overline{Q}$
\end{tabular}

\tcblower

\textbf{Official Key:} \textbf{\large D} \hfill \textit{\small [Easy]}
\vspace{2pt}\hrule height 0.3pt\vspace{3pt}
\textbf{Topic:} Digital Logic \\
\textbf{Subtopic:} \textit{Boolean algebra}
\vspace{2pt}\hrule height 0.3pt\vspace{3pt}
\textbf{Guide / Hint:} \\
{\small Recall that $P \oplus \overline{P} = 1$ and $1 \oplus Q = \overline{Q}$. On the right-hand side of (D), $P \odot \overline{P} = 0$, and $0 \odot \overline{Q} = \overline{0 \oplus \overline{Q}} = \overline{\overline{Q}} = Q$. Since $\overline{Q} \neq Q$, option (D) is NOT correct.}

\end{tcolorbox}
\vspace{6pt}\hrule height 0.35pt\vspace{6pt}
\noindent
\begin{tcolorbox}[
    blanker,
    sidebyside,
    sidebyside align=top,
    lefthand width=0.55\textwidth,
    top=4pt, bottom=4pt, left=0pt, right=0pt,
    middle=16pt,
    underlay={
        \draw[dashed, black!80, line width=0.8pt, dash pattern=on 4pt off 3pt]
            ([xshift=0.55\textwidth+8pt]frame.north west) --
            ([xshift=0.55\textwidth+8pt]frame.south west);
    }
]
\textbf{Q.69} \quad {\footnotesize\textsc{[NAT $\cdot$ 1 Mark(s) $\cdot$ No Neg]}} \hfill {\footnotesize\textbf{GATE 2023 CS} \quad \tiny\texttt{2023\_CS\_CS\_Q32}}
\vspace{2pt}\hrule height 0.3pt\vspace{4pt}

A particular number is written as 132 in radix-4 representation. The same number in radix-5 representation is \_\_\_\_\_\_\_\_\. (Answer in integer)

\par\vspace{6pt}\noindent\textbf{Your Answer:} \fbox{\rule{0pt}{1.3ex}\hspace{3cm}}

\tcblower

\textbf{Official Key:} \textbf{\large 110 to 110} \hfill \textit{\small [Easy]}
\vspace{2pt}\hrule height 0.3pt\vspace{3pt}
\textbf{Topic:} Digital Logic \\
\textbf{Subtopic:} \textit{Number representations and computer arithmetic (fixed and floating point)}
\vspace{2pt}\hrule height 0.3pt\vspace{3pt}
\textbf{Guide / Hint:} \\
{\small Converting $132_4$ to base 10: $1 \times 4^2 + 3 \times 4^1 + 2 \times 4^0 = 16 + 12 + 2 = 30_{10}$. Converting $30_{10}$ to radix 5: $30 = 6 \times 5 + 0$, $6 = 1 \times 5 + 1$, $1 = 0 \times 5 + 1$. Reading remainders from bottom to top yields $110_5$.}

\end{tcolorbox}
\vspace{6pt}\hrule height 0.35pt\vspace{6pt}
\noindent
\begin{tcolorbox}[
    blanker,
    sidebyside,
    sidebyside align=top,
    lefthand width=0.55\textwidth,
    top=4pt, bottom=4pt, left=0pt, right=0pt,
    middle=16pt,
    underlay={
        \draw[dashed, black!80, line width=0.8pt, dash pattern=on 4pt off 3pt]
            ([xshift=0.55\textwidth+8pt]frame.north west) --
            ([xshift=0.55\textwidth+8pt]frame.south west);
    }
]
\textbf{Q.70} \quad {\footnotesize\textsc{[MSQ $\cdot$ 2 Mark(s) $\cdot$ No Neg]}} \hfill {\footnotesize\textbf{GATE 2024 CS1} \quad \tiny\texttt{2024\_CS1\_CS\_Q47}}
\vspace{2pt}\hrule height 0.3pt\vspace{4pt}

Consider a Boolean expression given by $F(X, Y, Z) = \sum(3, 5, 6, 7)$. Which of the following statements is/are CORRECT?

\par\vspace{6pt}\noindent
\begin{tabular}{@{}p{0.98\linewidth}@{}}
$\square$\ \textbf{(A)} $F(X, Y, Z) = \prod(0, 1, 2, 4)$ \\
\addlinespace[3pt]
$\square$\ \textbf{(B)} $F(X, Y, Z) = XY + YZ + XZ$ \\
\addlinespace[3pt]
$\square$\ \textbf{(C)} $F(X, Y, Z)$ is independent of input $Y$ \\
\addlinespace[3pt]
$\square$\ \textbf{(D)} $F(X, Y, Z)$ is independent of input $X$
\end{tabular}

\tcblower

\textbf{Official Key:} \textbf{\large A;B} \hfill \textit{\small [Easy]}
\vspace{2pt}\hrule height 0.3pt\vspace{3pt}
\textbf{Topic:} Digital Logic \\
\textbf{Subtopic:} \textit{Minimization}
\vspace{2pt}\hrule height 0.3pt\vspace{3pt}
\textbf{Guide / Hint:} \\
{\small The maxterm representation is the product of maxterms not in the minterm list: $\prod(0, 1, 2, 4)$ (A is TRUE). Grouping minterms $\Sigma(3, 5, 6, 7)$ on a 3-variable K-map: pairs are $(3, 7) \implies YZ$, $(5, 7) \implies XZ$, and $(6, 7) \implies XY$, giving minimal SOP $F(X, Y, Z) = XY + YZ + XZ$ (B is TRUE).}

\end{tcolorbox}
\vspace{6pt}\hrule height 0.35pt\vspace{6pt}
\noindent
\begin{tcolorbox}[
    blanker,
    sidebyside,
    sidebyside align=top,
    lefthand width=0.55\textwidth,
    top=4pt, bottom=4pt, left=0pt, right=0pt,
    middle=16pt,
    underlay={
        \draw[dashed, black!80, line width=0.8pt, dash pattern=on 4pt off 3pt]
            ([xshift=0.55\textwidth+8pt]frame.north west) --
            ([xshift=0.55\textwidth+8pt]frame.south west);
    }
]
\textbf{Q.71} \quad {\footnotesize\textsc{[MCQ $\cdot$ 1 Mark(s) $\cdot$ Neg: 0.33]}} \hfill {\footnotesize\textbf{GATE 2024 CS2} \quad \tiny\texttt{2024\_CS2\_CS\_Q14}}
\vspace{2pt}\hrule height 0.3pt\vspace{4pt}

The format of a single-precision floating-point number as per the IEEE 754 standard is:
Sign (1 bit), Exponent (8 bits), Mantissa (23 bits).
Choose the largest floating-point number among the following options.

\par\vspace{6pt}\noindent
\begin{tabular}{@{}p{0.98\linewidth}@{}}
$\square$\ \textbf{(A)} Sign: 0, Exponent: 0111 1111, Mantissa: 1111 1111 1111 1111 1111 111 \\
\addlinespace[3pt]
$\square$\ \textbf{(B)} Sign: 0, Exponent: 1111 1110, Mantissa: 1111 1111 1111 1111 1111 111 \\
\addlinespace[3pt]
$\square$\ \textbf{(C)} Sign: 0, Exponent: 1111 1111, Mantissa: 1111 1111 1111 1111 1111 111 \\
\addlinespace[3pt]
$\square$\ \textbf{(D)} Sign: 0, Exponent: 0111 1111, Mantissa: 0000 0000 0000 0000 0000 000
\end{tabular}

\tcblower

\textbf{Official Key:} \textbf{\large B} \hfill \textit{\small [Easy]}
\vspace{2pt}\hrule height 0.3pt\vspace{3pt}
\textbf{Topic:} Digital Logic \\
\textbf{Subtopic:} \textit{Number representations and computer arithmetic (fixed and floating point)}
\vspace{2pt}\hrule height 0.3pt\vspace{3pt}
\textbf{Guide / Hint:} \\
{\small In IEEE 754 single precision: sign bit 0 means positive. The exponent field 11111111 is reserved for $\pm\infty$ and NaN. The largest finite representable number has the maximum non-reserved biased exponent, which is $254 = 11111110_2$, with all mantissa bits set to 1 ($111\dots1$), corresponding to option B.}

\end{tcolorbox}
\vspace{6pt}\hrule height 0.35pt\vspace{6pt}
\noindent
\begin{tcolorbox}[
    blanker,
    sidebyside,
    sidebyside align=top,
    lefthand width=0.55\textwidth,
    top=4pt, bottom=4pt, left=0pt, right=0pt,
    middle=16pt,
    underlay={
        \draw[dashed, black!80, line width=0.8pt, dash pattern=on 4pt off 3pt]
            ([xshift=0.55\textwidth+8pt]frame.north west) --
            ([xshift=0.55\textwidth+8pt]frame.south west);
    }
]
\textbf{Q.72} \quad {\footnotesize\textsc{[MSQ $\cdot$ 1 Mark(s) $\cdot$ No Neg]}} \hfill {\footnotesize\textbf{GATE 2024 CS2} \quad \tiny\texttt{2024\_CS2\_CS\_Q30}}
\vspace{2pt}\hrule height 0.3pt\vspace{4pt}

For a Boolean variable $x$, which of the following statements is/are FALSE?

\par\vspace{6pt}\noindent
\begin{tabular}{@{}p{0.48\linewidth}@{\hspace{0.04\linewidth}}p{0.48\linewidth}@{}}
$\square$\ \textbf{(A)} $x \cdot 1 = x$ & $\square$\ \textbf{(B)} $x + 1 = x$ \\
\addlinespace[3pt]
$\square$\ \textbf{(C)} $x \cdot x = 0$ & $\square$\ \textbf{(D)} $x + \overline{x} = 1$
\end{tabular}

\tcblower

\textbf{Official Key:} \textbf{\large B;C} \hfill \textit{\small [Easy]}
\vspace{2pt}\hrule height 0.3pt\vspace{3pt}
\textbf{Topic:} Digital Logic \\
\textbf{Subtopic:} \textit{Boolean algebra}
\vspace{2pt}\hrule height 0.3pt\vspace{3pt}
\textbf{Guide / Hint:} \\
{\small In Boolean algebra: $x + 1 = 1$ (null element property of OR), so $x + 1 = x$ is FALSE (B is false). Also, $x \cdot x = x$ (idempotent property of AND), so $x \cdot x = 0$ is FALSE (C is false).}

\end{tcolorbox}
\vspace{6pt}\hrule height 0.35pt\vspace{6pt}
\noindent
\begin{tcolorbox}[
    blanker,
    sidebyside,
    sidebyside align=top,
    lefthand width=0.55\textwidth,
    top=4pt, bottom=4pt, left=0pt, right=0pt,
    middle=16pt,
    underlay={
        \draw[dashed, black!80, line width=0.8pt, dash pattern=on 4pt off 3pt]
            ([xshift=0.55\textwidth+8pt]frame.north west) --
            ([xshift=0.55\textwidth+8pt]frame.south west);
    }
]
\textbf{Q.73} \quad {\footnotesize\textsc{[MSQ $\cdot$ 2 Mark(s) $\cdot$ No Neg]}} \hfill {\footnotesize\textbf{GATE 2024 CS2} \quad \tiny\texttt{2024\_CS2\_CS\_Q49}}
\vspace{2pt}\hrule height 0.3pt\vspace{4pt}

Which of the following is/are EQUAL to 224 in radix-5 (i.e., base-5) notation?

\par\vspace{6pt}\noindent
\begin{tabular}{@{}p{0.48\linewidth}@{\hspace{0.04\linewidth}}p{0.48\linewidth}@{}}
$\square$\ \textbf{(A)} 64 in radix-10 & $\square$\ \textbf{(B)} 100 in radix-8 \\
\addlinespace[3pt]
$\square$\ \textbf{(C)} 50 in radix-16 & $\square$\ \textbf{(D)} 121 in radix-7
\end{tabular}

\tcblower

\textbf{Official Key:} \textbf{\large A;B;D} \hfill \textit{\small [Easy]}
\vspace{2pt}\hrule height 0.3pt\vspace{3pt}
\textbf{Topic:} Digital Logic \\
\textbf{Subtopic:} \textit{Number representations and computer arithmetic (fixed and floating point)}
\vspace{2pt}\hrule height 0.3pt\vspace{3pt}
\textbf{Guide / Hint:} \\
{\small Convert $224_5$ to decimal: $2(5^2) + 2(5^1) + 4(5^0) = 50 + 10 + 4 = 64_{10}$ (A is true). In base 8: $64_{10} = 100_8$ (B is true). In base 16: $64_{10} = 40_{16} \neq 50_{16}$ (C is false). In base 7: $64_{10} = 1(7^2) + 2(7^1) + 1(7^0) = 121_7$ (D is true).}

\end{tcolorbox}
\vspace{6pt}\hrule height 0.35pt\vspace{6pt}
\noindent
\begin{tcolorbox}[
    blanker,
    sidebyside,
    sidebyside align=top,
    lefthand width=0.55\textwidth,
    top=4pt, bottom=4pt, left=0pt, right=0pt,
    middle=16pt,
    underlay={
        \draw[dashed, black!80, line width=0.8pt, dash pattern=on 4pt off 3pt]
            ([xshift=0.55\textwidth+8pt]frame.north west) --
            ([xshift=0.55\textwidth+8pt]frame.south west);
    }
]
\textbf{Q.74} \quad {\footnotesize\textsc{[MSQ $\cdot$ 1 Mark(s) $\cdot$ No Neg]}} \hfill {\footnotesize\textbf{GATE 2025 CS1} \quad \tiny\texttt{2025\_CS1\_CS\_Q25}}
\vspace{2pt}\hrule height 0.3pt\vspace{4pt}

The number $-6$ can be represented as $1010_2$ in 4-bit 2's complement representation. Which of the following is/are CORRECT 2's complement representation(s) of $-6$?

\par\vspace{6pt}\noindent
\begin{tabular}{@{}p{0.48\linewidth}@{\hspace{0.04\linewidth}}p{0.48\linewidth}@{}}
$\square$\ \textbf{(A)} 1000 1010 in 8-bits & $\square$\ \textbf{(B)} 1111 1010 in 8-bits \\
\addlinespace[3pt]
$\square$\ \textbf{(C)} 1000 0000 0000 1010 in 16-bits & $\square$\ \textbf{(D)} 1111 1111 1111 1010 in 16-bits
\end{tabular}

\tcblower

\textbf{Official Key:} \textbf{\large B;D} \hfill \textit{\small [Easy]}
\vspace{2pt}\hrule height 0.3pt\vspace{3pt}
\textbf{Topic:} Digital Logic \\
\textbf{Subtopic:} \textit{Number representations and computer arithmetic (fixed and floating point)}
\vspace{2pt}\hrule height 0.3pt\vspace{3pt}
\textbf{Guide / Hint:} \\
{\small Sign extension in two's complement preserves the numerical value by replicating the most significant bit (sign bit). Since the 4-bit representation of $-6$ is $1010$ (sign bit 1), extending to 8 bits gives $1111\,1010$ and extending to 16 bits gives $1111\,1111\,1111\,1010$.}

\end{tcolorbox}
\vspace{6pt}\hrule height 0.35pt\vspace{6pt}
\vspace{8pt}
\subsection*{Computer Organization and Architecture}
\vspace{-4pt}
\noindent
\begin{tcolorbox}[
    blanker,
    sidebyside,
    sidebyside align=top,
    lefthand width=0.55\textwidth,
    top=4pt, bottom=4pt, left=0pt, right=0pt,
    middle=16pt,
    underlay={
        \draw[dashed, black!80, line width=0.8pt, dash pattern=on 4pt off 3pt]
            ([xshift=0.55\textwidth+8pt]frame.north west) --
            ([xshift=0.55\textwidth+8pt]frame.south west);
    }
]
\textbf{Q.75} \quad {\footnotesize\textsc{[MCQ $\cdot$ 1 Mark(s) $\cdot$ Neg: 0.33]}} \hfill {\footnotesize\textbf{GATE 2018 CS} \quad \tiny\texttt{2018\_CS\_GA\_Q05\_v2}}
\vspace{2pt}\hrule height 0.3pt\vspace{4pt}

Consider the following processor design characteristics.
\begin{enumerate}[label=\Roman*.]
    \item Register-to-register arithmetic operations only
    \item Fixed-length instruction format
    \item Hardwired control unit
\end{enumerate}
Which of the characteristics above are used in the design of a RISC processor?

\par\vspace{6pt}\noindent
\begin{tabular}{@{}p{0.48\linewidth}@{\hspace{0.04\linewidth}}p{0.48\linewidth}@{}}
$\square$\ \textbf{(A)} I and II only & $\square$\ \textbf{(B)} II and III only \\
\addlinespace[3pt]
$\square$\ \textbf{(C)} I and III only & $\square$\ \textbf{(D)} I, II and III
\end{tabular}

\tcblower

\textbf{Official Key:} \textbf{\large D} \hfill \textit{\small [Easy]}
\vspace{2pt}\hrule height 0.3pt\vspace{3pt}
\textbf{Topic:} Computer Organization and Architecture \\
\textbf{Subtopic:} \textit{Machine instructions and addressing modes}
\vspace{2pt}\hrule height 0.3pt\vspace{3pt}
\textbf{Guide / Hint:} \\
{\small RISC (Reduced Instruction Set Computer) architectures feature load-store designs where arithmetic operations operate only on registers (I), fixed-length instruction encodings for streamlined pipelining (II), and hardwired control units for single-cycle instruction execution (III). Hence, all three characteristics I, II, and III are typical of RISC processors.}

\end{tcolorbox}
\vspace{6pt}\hrule height 0.35pt\vspace{6pt}
\noindent
\begin{tcolorbox}[
    blanker,
    sidebyside,
    sidebyside align=top,
    lefthand width=0.55\textwidth,
    top=4pt, bottom=4pt, left=0pt, right=0pt,
    middle=16pt,
    underlay={
        \draw[dashed, black!80, line width=0.8pt, dash pattern=on 4pt off 3pt]
            ([xshift=0.55\textwidth+8pt]frame.north west) --
            ([xshift=0.55\textwidth+8pt]frame.south west);
    }
]
\textbf{Q.76} \quad {\footnotesize\textsc{[MCQ $\cdot$ 1 Mark(s) $\cdot$ Neg: 0.33]}} \hfill {\footnotesize\textbf{GATE 2018 CS} \quad \tiny\texttt{2018\_CS\_GA\_Q09}}
\vspace{2pt}\hrule height 0.3pt\vspace{4pt}

The following are some events that occur after a device controller issues an interrupt while process $L$ is under execution.
\begin{itemize}[leftmargin=*, nosep]
    \item[P.] The processor pushes the process status of $L$ on to the control stack.
    \item[Q.] The processor finishes the execution of the current instruction.
    \item[R.] The processor executes the interrupt service routine.
    \item[S.] The processor pops the process status of $L$ from the control stack.
    \item[T.] The processor fetches the new PC value from the interrupt vector.
\end{itemize}
Which one of the following is the correct order of events?

\par\vspace{6pt}\noindent
\begin{tabular}{@{}p{0.48\linewidth}@{\hspace{0.04\linewidth}}p{0.48\linewidth}@{}}
$\square$\ \textbf{(A)} Q, P, T, R, S & $\square$\ \textbf{(B)} P, Q, T, R, S \\
\addlinespace[3pt]
$\square$\ \textbf{(C)} Q, P, R, T, S & $\square$\ \textbf{(D)} P, Q, R, T, S
\end{tabular}

\tcblower

\textbf{Official Key:} \textbf{\large A} \hfill \textit{\small [Easy]}
\vspace{2pt}\hrule height 0.3pt\vspace{3pt}
\textbf{Topic:} Computer Organization and Architecture \\
\textbf{Subtopic:} \textit{I/O interface: interrupt and DMA mode}
\vspace{2pt}\hrule height 0.3pt\vspace{3pt}
\textbf{Guide / Hint:} \\
{\small Upon receiving an interrupt, the processor finishes the execution of the current instruction (Q), saves the processor status and return PC onto the stack (P), loads the starting address of the ISR from the interrupt vector table (T), executes the ISR (R), and finally restores the saved status and resumes process $L$ (S). Hence the order is QPTRS.}

\end{tcolorbox}
\vspace{6pt}\hrule height 0.35pt\vspace{6pt}
\noindent
\begin{tcolorbox}[
    blanker,
    sidebyside,
    sidebyside align=top,
    lefthand width=0.55\textwidth,
    top=4pt, bottom=4pt, left=0pt, right=0pt,
    middle=16pt,
    underlay={
        \draw[dashed, black!80, line width=0.8pt, dash pattern=on 4pt off 3pt]
            ([xshift=0.55\textwidth+8pt]frame.north west) --
            ([xshift=0.55\textwidth+8pt]frame.south west);
    }
]
\textbf{Q.77} \quad {\footnotesize\textsc{[MCQ $\cdot$ 1 Mark(s) $\cdot$ Neg: 0.33]}} \hfill {\footnotesize\textbf{GATE 2022 CS} \quad \tiny\texttt{2022\_CS\_CS\_Q17}}
\vspace{2pt}\hrule height 0.3pt\vspace{4pt}

Which one of the following facilitates transfer of bulk data from hard disk to main memory with the highest throughput?

\par\vspace{6pt}\noindent
\begin{tabular}{@{}p{0.48\linewidth}@{\hspace{0.04\linewidth}}p{0.48\linewidth}@{}}
$\square$\ \textbf{(A)} DMA based I/O transfer & $\square$\ \textbf{(B)} Interrupt driven I/O transfer \\
\addlinespace[3pt]
$\square$\ \textbf{(C)} Polling based I/O transfer & $\square$\ \textbf{(D)} Programmed I/O transfer
\end{tabular}

\tcblower

\textbf{Official Key:} \textbf{\large A} \hfill \textit{\small [Easy]}
\vspace{2pt}\hrule height 0.3pt\vspace{3pt}
\textbf{Topic:} Computer Organization and Architecture \\
\textbf{Subtopic:} \textit{I/O interface: interrupt and DMA mode}
\vspace{2pt}\hrule height 0.3pt\vspace{3pt}
\textbf{Guide / Hint:} \\
{\small Direct Memory Access (DMA) enables high-speed bulk data transfer directly between the disk controller and main memory without continuous CPU intervention, achieving the highest possible data throughput compared to programmed or interrupt-driven I/O.}

\end{tcolorbox}
\vspace{6pt}\hrule height 0.35pt\vspace{6pt}
\noindent
\begin{tcolorbox}[
    blanker,
    sidebyside,
    sidebyside align=top,
    lefthand width=0.55\textwidth,
    top=4pt, bottom=4pt, left=0pt, right=0pt,
    middle=16pt,
    underlay={
        \draw[dashed, black!80, line width=0.8pt, dash pattern=on 4pt off 3pt]
            ([xshift=0.55\textwidth+8pt]frame.north west) --
            ([xshift=0.55\textwidth+8pt]frame.south west);
    }
]
\textbf{Q.78} \quad {\footnotesize\textsc{[NAT $\cdot$ 1 Mark(s) $\cdot$ No Neg]}} \hfill {\footnotesize\textbf{GATE 2023 CS} \quad \tiny\texttt{2023\_CS\_CS\_Q33}}
\vspace{2pt}\hrule height 0.3pt\vspace{4pt}

Consider a 3-stage pipelined processor having a delay of 10 ns (nanoseconds), 20 ns, and 14 ns, for the first, second, and the third stages, respectively. Assume that there is no other delay and the processor does not suffer from any pipeline hazards. Also assume that one instruction is fetched every cycle.
The total execution time for executing 100 instructions on this processor is \_\_\_\_\_\_\_\_\_ ns. (Answer in integer)

\par\vspace{6pt}\noindent\textbf{Your Answer:} \fbox{\rule{0pt}{1.3ex}\hspace{3cm}}

\tcblower

\textbf{Official Key:} \textbf{\large 2040 to 2040} \hfill \textit{\small [Easy]}
\vspace{2pt}\hrule height 0.3pt\vspace{3pt}
\textbf{Topic:} Computer Organization and Architecture \\
\textbf{Subtopic:} \textit{Instruction pipelining, pipeline hazards}
\vspace{2pt}\hrule height 0.3pt\vspace{3pt}
\textbf{Guide / Hint:} \\
{\small The clock period of the pipeline is governed by the longest stage delay: $T = \max(10, 20, 14) = 20\text{ ns}$. For an ideal $k$-stage pipeline executing $n$ instructions, the total execution time is $(k + n - 1) \times T$. For $k = 3$ and $n = 100$: $(3 + 100 - 1) \times 20 = 102 \times 20 = 2040\text{ ns}$.}

\end{tcolorbox}
\vspace{6pt}\hrule height 0.35pt\vspace{6pt}
\noindent
\begin{tcolorbox}[
    blanker,
    sidebyside,
    sidebyside align=top,
    lefthand width=0.55\textwidth,
    top=4pt, bottom=4pt, left=0pt, right=0pt,
    middle=16pt,
    underlay={
        \draw[dashed, black!80, line width=0.8pt, dash pattern=on 4pt off 3pt]
            ([xshift=0.55\textwidth+8pt]frame.north west) --
            ([xshift=0.55\textwidth+8pt]frame.south west);
    }
]
\textbf{Q.79} \quad {\footnotesize\textsc{[NAT $\cdot$ 2 Mark(s) $\cdot$ No Neg]}} \hfill {\footnotesize\textbf{GATE 2023 CS} \quad \tiny\texttt{2023\_CS\_CS\_Q64}}
\vspace{2pt}\hrule height 0.3pt\vspace{4pt}

An 8-way set associative cache of size 64 KB (1 KB = 1024 bytes) is used in a system with 32-bit address. The address is sub-divided into TAG, INDEX, and BLOCK OFFSET.
The number of bits in the TAG is \_\_\_\_\_\_\_\_\. (Answer in integer)

\par\vspace{6pt}\noindent\textbf{Your Answer:} \fbox{\rule{0pt}{1.3ex}\hspace{3cm}}

\tcblower

\textbf{Official Key:} \textbf{\large 19 to 19} \hfill \textit{\small [Easy]}
\vspace{2pt}\hrule height 0.3pt\vspace{3pt}
\textbf{Topic:} Computer Organization and Architecture \\
\textbf{Subtopic:} \textit{Memory hierarchy: cache, main memory and secondary storage}
\vspace{2pt}\hrule height 0.3pt\vspace{3pt}
\textbf{Guide / Hint:} \\
{\small Address size $= 32\text{ bits}$. Cache size $= 64\text{ KB} = 2^{16}\text{ bytes}$. For an 8-way set-associative cache, the combined capacity of index and offset is $\frac{\text{Cache Size}}{\text{Associativity}} = \frac{64\text{ KB}}{8} = 8\text{ KB} = 2^{13}\text{ bytes}$, requiring $13\text{ bits}$ for INDEX + BLOCK OFFSET. Therefore, the number of TAG bits is $32 - 13 = 19\text{ bits}$.}

\end{tcolorbox}
\vspace{6pt}\hrule height 0.35pt\vspace{6pt}
\noindent
\begin{tcolorbox}[
    blanker,
    sidebyside,
    sidebyside align=top,
    lefthand width=0.55\textwidth,
    top=4pt, bottom=4pt, left=0pt, right=0pt,
    middle=16pt,
    underlay={
        \draw[dashed, black!80, line width=0.8pt, dash pattern=on 4pt off 3pt]
            ([xshift=0.55\textwidth+8pt]frame.north west) --
            ([xshift=0.55\textwidth+8pt]frame.south west);
    }
]
\textbf{Q.80} \quad {\footnotesize\textsc{[MCQ $\cdot$ 1 Mark(s) $\cdot$ Neg: 0.33]}} \hfill {\footnotesize\textbf{GATE 2024 CS1} \quad \tiny\texttt{2024\_CS1\_CS\_Q15}}
\vspace{2pt}\hrule height 0.3pt\vspace{4pt}

Which one of the following statements is FALSE?

\par\vspace{6pt}\noindent
\begin{tabular}{@{}p{0.98\linewidth}@{}}
$\square$\ \textbf{(A)} In the cycle stealing mode of DMA, one word of data is transferred between an I/O device and main memory in a stolen cycle \\
\addlinespace[3pt]
$\square$\ \textbf{(B)} For bulk data transfer, the burst mode of DMA has a higher throughput than the cycle stealing mode \\
\addlinespace[3pt]
$\square$\ \textbf{(C)} Programmed I/O mechanism has a better CPU utilization than the interrupt driven I/O mechanism \\
\addlinespace[3pt]
$\square$\ \textbf{(D)} The CPU can start executing an interrupt service routine faster with vectored interrupts than with non-vectored interrupts
\end{tabular}

\tcblower

\textbf{Official Key:} \textbf{\large C} \hfill \textit{\small [Easy]}
\vspace{2pt}\hrule height 0.3pt\vspace{3pt}
\textbf{Topic:} Computer Organization and Architecture \\
\textbf{Subtopic:} \textit{I/O interface: interrupt and DMA mode}
\vspace{2pt}\hrule height 0.3pt\vspace{3pt}
\textbf{Guide / Hint:} \\
{\small Programmed I/O requires the CPU to stay in a busy-waiting polling loop to check device readiness, wasting processor cycles and yielding poor CPU utilization. Interrupt-driven I/O frees the CPU to execute other tasks while I/O operations are underway, providing significantly better CPU utilization. Thus, statement (C) is FALSE.}

\end{tcolorbox}
\vspace{6pt}\hrule height 0.35pt\vspace{6pt}
\noindent
\begin{tcolorbox}[
    blanker,
    sidebyside,
    sidebyside align=top,
    lefthand width=0.55\textwidth,
    top=4pt, bottom=4pt, left=0pt, right=0pt,
    middle=16pt,
    underlay={
        \draw[dashed, black!80, line width=0.8pt, dash pattern=on 4pt off 3pt]
            ([xshift=0.55\textwidth+8pt]frame.north west) --
            ([xshift=0.55\textwidth+8pt]frame.south west);
    }
]
\textbf{Q.81} \quad {\footnotesize\textsc{[MSQ $\cdot$ 1 Mark(s) $\cdot$ No Neg]}} \hfill {\footnotesize\textbf{GATE 2024 CS1} \quad \tiny\texttt{2024\_CS1\_CS\_Q30}}
\vspace{2pt}\hrule height 0.3pt\vspace{4pt}

Consider a 5-stage pipelined processor with Instruction Fetch (IF), Instruction Decode (ID), Execute (EX), Memory Access (MEM), and Register Writeback (WB) stages. Which of the following statements about forwarding is/are CORRECT?

\par\vspace{6pt}\noindent
\begin{tabular}{@{}p{0.98\linewidth}@{}}
$\square$\ \textbf{(A)} In a pipelined execution, forwarding means the result from a source stage of an earlier instruction is passed on to the destination stage of a later instruction \\
\addlinespace[3pt]
$\square$\ \textbf{(B)} In forwarding, data from the output of the MEM stage can be passed on to the input of the EX stage of the next instruction \\
\addlinespace[3pt]
$\square$\ \textbf{(C)} Forwarding cannot prevent all pipeline stalls \\
\addlinespace[3pt]
$\square$\ \textbf{(D)} Forwarding does not require any extra hardware to retrieve the data from the pipeline stages
\end{tabular}

\tcblower

\textbf{Official Key:} \textbf{\large A;C} \hfill \textit{\small [Easy]}
\vspace{2pt}\hrule height 0.3pt\vspace{3pt}
\textbf{Topic:} Computer Organization and Architecture \\
\textbf{Subtopic:} \textit{Instruction pipelining, pipeline hazards}
\vspace{2pt}\hrule height 0.3pt\vspace{3pt}
\textbf{Guide / Hint:} \\
{\small Forwarding (bypassing) routes intermediate results directly from pipeline registers of an earlier instruction to execution inputs of a subsequent instruction (A is TRUE). While forwarding resolves many ALU data dependencies without stalls, it cannot prevent a load-use stall (where the loaded data is needed immediately in the subsequent instruction's EX stage), requiring at least one stall cycle (C is TRUE).}

\end{tcolorbox}
\vspace{6pt}\hrule height 0.35pt\vspace{6pt}
\noindent
\begin{tcolorbox}[
    blanker,
    sidebyside,
    sidebyside align=top,
    lefthand width=0.55\textwidth,
    top=4pt, bottom=4pt, left=0pt, right=0pt,
    middle=16pt,
    underlay={
        \draw[dashed, black!80, line width=0.8pt, dash pattern=on 4pt off 3pt]
            ([xshift=0.55\textwidth+8pt]frame.north west) --
            ([xshift=0.55\textwidth+8pt]frame.south west);
    }
]
\textbf{Q.82} \quad {\footnotesize\textsc{[NAT $\cdot$ 2 Mark(s) $\cdot$ No Neg]}} \hfill {\footnotesize\textbf{GATE 2024 CS1} \quad \tiny\texttt{2024\_CS1\_CS\_Q54}}
\vspace{2pt}\hrule height 0.3pt\vspace{4pt}

Consider a 512 GB hard disk with 32 storage surfaces. There are 4096 sectors per track and each sector holds 1024 bytes of data. The number of cylinders in the hard disk is \_\_\_\_\_\_\_\_\_\. (Answer in integer)

\par\vspace{6pt}\noindent\textbf{Your Answer:} \fbox{\rule{0pt}{1.3ex}\hspace{3cm}}

\tcblower

\textbf{Official Key:} \textbf{\large 4096 to 4096} \hfill \textit{\small [Easy]}
\vspace{2pt}\hrule height 0.3pt\vspace{3pt}
\textbf{Topic:} Computer Organization and Architecture \\
\textbf{Subtopic:} \textit{Memory hierarchy: cache, main memory and secondary storage}
\vspace{2pt}\hrule height 0.3pt\vspace{3pt}
\textbf{Guide / Hint:} \\
{\small Disk capacity $= 512\text{ GB} = 512 \times 2^{30}\text{ bytes} = 2^{39}$ bytes. One cylinder comprises 32 surfaces $\times 4096$ sectors/track $\times 1024$ bytes/sector $= 2^5 \times 2^{12} \times 2^{10} = 2^{27}$ bytes. Number of cylinders $= \frac{2^{39}}{2^{27}} = 2^{12} = 4096$.}

\end{tcolorbox}
\vspace{6pt}\hrule height 0.35pt\vspace{6pt}
\noindent
\begin{tcolorbox}[
    blanker,
    sidebyside,
    sidebyside align=top,
    lefthand width=0.55\textwidth,
    top=4pt, bottom=4pt, left=0pt, right=0pt,
    middle=16pt,
    underlay={
        \draw[dashed, black!80, line width=0.8pt, dash pattern=on 4pt off 3pt]
            ([xshift=0.55\textwidth+8pt]frame.north west) --
            ([xshift=0.55\textwidth+8pt]frame.south west);
    }
]
\textbf{Q.83} \quad {\footnotesize\textsc{[MSQ $\cdot$ 1 Mark(s) $\cdot$ No Neg]}} \hfill {\footnotesize\textbf{GATE 2024 CS2} \quad \tiny\texttt{2024\_CS2\_CS\_Q31}}
\vspace{2pt}\hrule height 0.3pt\vspace{4pt}

An instruction format has the following structure:
Instruction Number: Opcode destination reg, source reg-1, source reg-2
Consider the following sequence of instructions to be executed in a pipelined processor:
\begin{verbatim}
I1: DIV R3, R1, R2
I2: SUB R5, R3, R4
I3: ADD R3, R5, R6
I4: MUL R7, R3, R8
\end{verbatim}
Which of the following dependencies is/are TRUE?

\par\vspace{6pt}\noindent
\begin{tabular}{@{}p{0.98\linewidth}@{}}
$\square$\ \textbf{(A)} There is a RAW dependency on R3 between I1 and I2 \\
\addlinespace[3pt]
$\square$\ \textbf{(B)} There is a WAR dependency on R3 between I1 and I3 \\
\addlinespace[3pt]
$\square$\ \textbf{(C)} There is a RAW dependency on R3 between I2 and I3 \\
\addlinespace[3pt]
$\square$\ \textbf{(D)} There is a WAW dependency on R3 between I3 and I4
\end{tabular}

\tcblower

\textbf{Official Key:} \textbf{\large A} \hfill \textit{\small [Easy]}
\vspace{2pt}\hrule height 0.3pt\vspace{3pt}
\textbf{Topic:} Computer Organization and Architecture \\
\textbf{Subtopic:} \textit{Instruction pipelining, pipeline hazards}
\vspace{2pt}\hrule height 0.3pt\vspace{3pt}
\textbf{Guide / Hint:} \\
{\small In I1: `DIV R3, R1, R2` (writes R3); I2: `SUB R5, R3, R4` (reads R3). Since I2 reads R3 after I1 writes R3, this is a Read-After-Write (RAW) data dependency on R3 between I1 and I2 (A is true). I3 writes R3, so between I1 (write) and I3 (write) is WAW, and between I2 (read) and I3 (write) is WAR.}

\end{tcolorbox}
\vspace{6pt}\hrule height 0.35pt\vspace{6pt}
\noindent
\begin{tcolorbox}[
    blanker,
    sidebyside,
    sidebyside align=top,
    lefthand width=0.55\textwidth,
    top=4pt, bottom=4pt, left=0pt, right=0pt,
    middle=16pt,
    underlay={
        \draw[dashed, black!80, line width=0.8pt, dash pattern=on 4pt off 3pt]
            ([xshift=0.55\textwidth+8pt]frame.north west) --
            ([xshift=0.55\textwidth+8pt]frame.south west);
    }
]
\textbf{Q.84} \quad {\footnotesize\textsc{[MCQ $\cdot$ 1 Mark(s) $\cdot$ Neg: 0.33]}} \hfill {\footnotesize\textbf{GATE 2025 CS1} \quad \tiny\texttt{2025\_CS1\_CS\_Q11}}
\vspace{2pt}\hrule height 0.3pt\vspace{4pt}

Suppose a program is running on a non-pipelined single processor computer system. The computer is connected to an external device that can interrupt the processor asynchronously. The processor needs to execute the interrupt service routine (ISR) to serve this interrupt. The following steps (not necessarily in order) are taken by the processor when the interrupt arrives:
\begin{enumerate}[label=(\roman*)]
\item The processor saves the content of the program counter.
\item The program counter is loaded with the start address of the ISR.
\item The processor finishes the present instruction.
\end{enumerate}
Which ONE of the following is the CORRECT sequence of steps?

\par\vspace{6pt}\noindent
\begin{tabular}{@{}p{0.48\linewidth}@{\hspace{0.04\linewidth}}p{0.48\linewidth}@{}}
$\square$\ \textbf{(A)} (iii), (i), (ii) & $\square$\ \textbf{(B)} (i), (iii), (ii) \\
\addlinespace[3pt]
$\square$\ \textbf{(C)} (i), (ii), (iii) & $\square$\ \textbf{(D)} (iii), (ii), (i)
\end{tabular}

\tcblower

\textbf{Official Key:} \textbf{\large A} \hfill \textit{\small [Easy]}
\vspace{2pt}\hrule height 0.3pt\vspace{3pt}
\textbf{Topic:} Computer Organization and Architecture \\
\textbf{Subtopic:} \textit{I/O interface: interrupt and DMA mode}
\vspace{2pt}\hrule height 0.3pt\vspace{3pt}
\textbf{Guide / Hint:} \\
{\small When an asynchronous interrupt arrives, the processor first finishes the execution of the current instruction (iii). Next, it saves the current program counter (and processor state) onto the stack (i). Finally, it loads the starting address of the interrupt service routine into the program counter (ii).}

\end{tcolorbox}
\vspace{6pt}\hrule height 0.35pt\vspace{6pt}
\noindent
\begin{tcolorbox}[
    blanker,
    sidebyside,
    sidebyside align=top,
    lefthand width=0.55\textwidth,
    top=4pt, bottom=4pt, left=0pt, right=0pt,
    middle=16pt,
    underlay={
        \draw[dashed, black!80, line width=0.8pt, dash pattern=on 4pt off 3pt]
            ([xshift=0.55\textwidth+8pt]frame.north west) --
            ([xshift=0.55\textwidth+8pt]frame.south west);
    }
]
\textbf{Q.85} \quad {\footnotesize\textsc{[MSQ $\cdot$ 1 Mark(s) $\cdot$ No Neg]}} \hfill {\footnotesize\textbf{GATE 2025 CS2} \quad \tiny\texttt{2025\_CS2\_CS\_Q28}}
\vspace{2pt}\hrule height 0.3pt\vspace{4pt}

Which of the following is/are part of an Instruction Set Architecture of a processor?

\par\vspace{6pt}\noindent
\begin{tabular}{@{}p{0.98\linewidth}@{}}
$\square$\ \textbf{(A)} The size of the cache memory \\
\addlinespace[3pt]
$\square$\ \textbf{(B)} The clock frequency of the processor \\
\addlinespace[3pt]
$\square$\ \textbf{(C)} The number of cache memory levels \\
\addlinespace[3pt]
$\square$\ \textbf{(D)} The total number of registers
\end{tabular}

\tcblower

\textbf{Official Key:} \textbf{\large D} \hfill \textit{\small [Easy]}
\vspace{2pt}\hrule height 0.3pt\vspace{3pt}
\textbf{Topic:} Computer Organization and Architecture \\
\textbf{Subtopic:} \textit{Machine instructions and addressing modes: ISA}
\vspace{2pt}\hrule height 0.3pt\vspace{3pt}
\textbf{Guide / Hint:} \\
{\small The Instruction Set Architecture (ISA) defines programmer-visible specifications, including the register set/count and instruction encoding. Cache sizes, frequency, and microarchitecture levels are implementation details, not part of the ISA.}

\end{tcolorbox}
\vspace{6pt}\hrule height 0.35pt\vspace{6pt}
\vspace{8pt}
\subsection*{Programming and Data Structures}
\vspace{-4pt}
\noindent
\begin{tcolorbox}[
    blanker,
    sidebyside,
    sidebyside align=top,
    lefthand width=0.55\textwidth,
    top=4pt, bottom=4pt, left=0pt, right=0pt,
    middle=16pt,
    underlay={
        \draw[dashed, black!80, line width=0.8pt, dash pattern=on 4pt off 3pt]
            ([xshift=0.55\textwidth+8pt]frame.north west) --
            ([xshift=0.55\textwidth+8pt]frame.south west);
    }
]
\textbf{Q.86} \quad {\footnotesize\textsc{[MCQ $\cdot$ 1 Mark(s) $\cdot$ Neg: 0.33]}} \hfill {\footnotesize\textbf{GATE 2018 CS} \quad \tiny\texttt{2018\_CS\_GA\_Q02\_v2}}
\vspace{2pt}\hrule height 0.3pt\vspace{4pt}

Consider the following C program.
\begin{lstlisting}[language=C]
#include <stdio.h>
struct Ournode {
    char x, y, z;
};
\par\medskip
int main() {
    struct Ournode p = {'1', '0', 'a' + 2};
    struct Ournode *q = &p;
    printf("%c, %c", *((char*)q + 1), *((char*)q + 2));
    return 0;
}
\end{lstlisting}
The output of this program is

\par\vspace{6pt}\noindent
\begin{tabular}{@{}p{0.48\linewidth}@{\hspace{0.04\linewidth}}p{0.48\linewidth}@{}}
$\square$\ \textbf{(A)} 0, c & $\square$\ \textbf{(B)} 0, a+2 \\
\addlinespace[3pt]
$\square$\ \textbf{(C)} '0', 'p+2' & $\square$\ \textbf{(D)} 1, 0
\end{tabular}

\tcblower

\textbf{Official Key:} \textbf{\large A} \hfill \textit{\small [Easy]}
\vspace{2pt}\hrule height 0.3pt\vspace{3pt}
\textbf{Topic:} Programming and Data Structures \\
\textbf{Subtopic:} \textit{Programming in C}
\vspace{2pt}\hrule height 0.3pt\vspace{3pt}
\textbf{Guide / Hint:} \\
{\small In `struct Ournode`, `char x, y, z;` are stored contiguously in memory. Since `p = {'1', '0', 'a'+2}`, `p.x = '1'`, `p.y = '0'`, and `p.z = 'c'`. Pointer `q` points to `p`; casting `q` to `(char*)` gives byte-level access where `*((char*)q + 1)` accesses `p.y` ('0') and `*((char*)q + 2)` accesses `p.z` ('c').}

\end{tcolorbox}
\vspace{6pt}\hrule height 0.35pt\vspace{6pt}
\noindent
\begin{tcolorbox}[
    blanker,
    sidebyside,
    sidebyside align=top,
    lefthand width=0.55\textwidth,
    top=4pt, bottom=4pt, left=0pt, right=0pt,
    middle=16pt,
    underlay={
        \draw[dashed, black!80, line width=0.8pt, dash pattern=on 4pt off 3pt]
            ([xshift=0.55\textwidth+8pt]frame.north west) --
            ([xshift=0.55\textwidth+8pt]frame.south west);
    }
]
\textbf{Q.87} \quad {\footnotesize\textsc{[MCQ $\cdot$ 1 Mark(s) $\cdot$ Neg: 0.33]}} \hfill {\footnotesize\textbf{GATE 2022 CS} \quad \tiny\texttt{2022\_CS\_CS\_Q15}}
\vspace{2pt}\hrule height 0.3pt\vspace{4pt}

Consider the problem of reversing a singly linked list. Which one of the following statements is TRUE about the time complexity of algorithms that solve the above problem in $\mathcal{O}(1)$ space?

\par\vspace{6pt}\noindent
\begin{tabular}{@{}p{0.98\linewidth}@{}}
$\square$\ \textbf{(A)} The best algorithm for the problem takes $\Theta(n)$ time in the worst case. \\
\addlinespace[3pt]
$\square$\ \textbf{(B)} The best algorithm for the problem takes $\Theta(n \log n)$ time in the worst case. \\
\addlinespace[3pt]
$\square$\ \textbf{(C)} The best algorithm for the problem takes $\Theta(n^2)$ time in the worst case. \\
\addlinespace[3pt]
$\square$\ \textbf{(D)} It is not possible to reverse a singly linked list in $\mathcal{O}(1)$ space.
\end{tabular}

\tcblower

\textbf{Official Key:} \textbf{\large A} \hfill \textit{\small [Easy]}
\vspace{2pt}\hrule height 0.3pt\vspace{3pt}
\textbf{Topic:} Programming and Data Structures \\
\textbf{Subtopic:} \textit{Arrays, stacks, queues, linked lists}
\vspace{2pt}\hrule height 0.3pt\vspace{3pt}
\textbf{Guide / Hint:} \\
{\small A singly linked list can be reversed in-place in $\mathcal{O}(1)$ auxiliary space using three moving pointers (`prev`, `curr`, `next`) by reversing each pointer sequentially in a single pass of $n$ nodes, achieving optimal $\Theta(n)$ worst-case time.}

\end{tcolorbox}
\vspace{6pt}\hrule height 0.35pt\vspace{6pt}
\noindent
\begin{tcolorbox}[
    blanker,
    sidebyside,
    sidebyside align=top,
    lefthand width=0.55\textwidth,
    top=4pt, bottom=4pt, left=0pt, right=0pt,
    middle=16pt,
    underlay={
        \draw[dashed, black!80, line width=0.8pt, dash pattern=on 4pt off 3pt]
            ([xshift=0.55\textwidth+8pt]frame.north west) --
            ([xshift=0.55\textwidth+8pt]frame.south west);
    }
]
\textbf{Q.88} \quad {\footnotesize\textsc{[MCQ $\cdot$ 1 Mark(s) $\cdot$ Neg: 0.33]}} \hfill {\footnotesize\textbf{GATE 2022 CS} \quad \tiny\texttt{2022\_CS\_CS\_Q21}}
\vspace{2pt}\hrule height 0.3pt\vspace{4pt}

What is printed by the following ANSI C program?
\begin{lstlisting}[language=C]
#include <stdio.h>
int main(int argc, char *argv[]) {
    int x = 1, z[2] = {10, 11};
    int *p = NULL;
    p = &x;
    *p = 10;
    p = &z[1];
    *(&z[0] + 1) += 3;
    printf("%d, %d, %d\n", x, z[0], z[1]);
    return 0;
}
\end{lstlisting}

\par\vspace{6pt}\noindent
\begin{tabular}{@{}p{0.48\linewidth}@{\hspace{0.04\linewidth}}p{0.48\linewidth}@{}}
$\square$\ \textbf{(A)} 1, 10, 11 & $\square$\ \textbf{(B)} 1, 10, 14 \\
\addlinespace[3pt]
$\square$\ \textbf{(C)} 10, 14, 11 & $\square$\ \textbf{(D)} 10, 10, 14
\end{tabular}

\tcblower

\textbf{Official Key:} \textbf{\large D} \hfill \textit{\small [Easy]}
\vspace{2pt}\hrule height 0.3pt\vspace{3pt}
\textbf{Topic:} Programming and Data Structures \\
\textbf{Subtopic:} \textit{Programming in C}
\vspace{2pt}\hrule height 0.3pt\vspace{3pt}
\textbf{Guide / Hint:} \\
{\small Initially $x = 1, z[0] = 10, z[1] = 11$. The assignment `p = \&x; *p = 10;` sets $x = 10$. The line `*(\&z[0] + 1) += 3;` targets `z[1]`, updating $z[1] = 11 + 3 = 14$. The array element $z[0]$ remains 10. The output printed is `10, 10, 14`.}

\end{tcolorbox}
\vspace{6pt}\hrule height 0.35pt\vspace{6pt}
\noindent
\begin{tcolorbox}[
    blanker,
    sidebyside,
    sidebyside align=top,
    lefthand width=0.55\textwidth,
    top=4pt, bottom=4pt, left=0pt, right=0pt,
    middle=16pt,
    underlay={
        \draw[dashed, black!80, line width=0.8pt, dash pattern=on 4pt off 3pt]
            ([xshift=0.55\textwidth+8pt]frame.north west) --
            ([xshift=0.55\textwidth+8pt]frame.south west);
    }
]
\textbf{Q.89} \quad {\footnotesize\textsc{[MCQ $\cdot$ 2 Mark(s) $\cdot$ Neg: 0.66]}} \hfill {\footnotesize\textbf{GATE 2022 CS} \quad \tiny\texttt{2022\_CS\_CS\_Q43}}
\vspace{2pt}\hrule height 0.3pt\vspace{4pt}

What is printed by the following ANSI C program?
\begin{lstlisting}[language=C]
#include <stdio.h>
int main(int argc, char *argv[]) {
    int a[3][3][3] = {
        {1, 2, 3, 4, 5, 6, 7, 8, 9},
        {10, 11, 12, 13, 14, 15, 16, 17, 18},
        {19, 20, 21, 22, 23, 24, 25, 26, 27}
    };
    int i = 0, j = 0, k = 0;
    for(i = 0; i < 3; i++) {
        for(k = 0; k < 3; k++)
            printf("%d ", a[i][j][k]);
        printf("\n");
    }
    return 0;
}
\end{lstlisting}

\par\vspace{6pt}\noindent
\begin{tabular}{@{}p{0.48\linewidth}@{\hspace{0.04\linewidth}}p{0.48\linewidth}@{}}
$\square$\ \textbf{(A)} 1 2 3 \\ 10 11 12 \\ 19 20 21 & $\square$\ \textbf{(B)} 1 4 7 \\ 10 13 16 \\ 19 22 25 \\
\addlinespace[3pt]
$\square$\ \textbf{(C)} 1 2 3 \\ 4 5 6 \\ 7 8 9 & $\square$\ \textbf{(D)} 1 2 3 \\ 13 14 15 \\ 25 26 27
\end{tabular}

\tcblower

\textbf{Official Key:} \textbf{\large A} \hfill \textit{\small [Easy]}
\vspace{2pt}\hrule height 0.3pt\vspace{3pt}
\textbf{Topic:} Programming and Data Structures \\
\textbf{Subtopic:} \textit{Programming in C}
\vspace{2pt}\hrule height 0.3pt\vspace{3pt}
\textbf{Guide / Hint:} \\
{\small In the nested loops, index $j = 0$ is fixed throughout. For $i = 0$, $k$ iterates from 0 to 2, printing `a[0][0][0], a[0][0][1], a[0][0][2]` ($1\ 2\ 3$). For $i = 1$, it prints `a[1][0][0], a[1][0][1], a[1][0][2]` ($10\ 11\ 12$). For $i = 2$, it prints `a[2][0][0], a[2][0][1], a[2][0][2]` ($19\ 20\ 21$).}

\end{tcolorbox}
\vspace{6pt}\hrule height 0.35pt\vspace{6pt}
\noindent
\begin{tcolorbox}[
    blanker,
    sidebyside,
    sidebyside align=top,
    lefthand width=0.55\textwidth,
    top=4pt, bottom=4pt, left=0pt, right=0pt,
    middle=16pt,
    underlay={
        \draw[dashed, black!80, line width=0.8pt, dash pattern=on 4pt off 3pt]
            ([xshift=0.55\textwidth+8pt]frame.north west) --
            ([xshift=0.55\textwidth+8pt]frame.south west);
    }
]
\textbf{Q.90} \quad {\footnotesize\textsc{[MCQ $\cdot$ 1 Mark(s) $\cdot$ Neg: 0.33]}} \hfill {\footnotesize\textbf{GATE 2023 CS} \quad \tiny\texttt{2023\_CS\_CS\_Q12}}
\vspace{2pt}\hrule height 0.3pt\vspace{4pt}

Which one of the following sequences when stored in an array at locations $A[1], \dots, A[10]$ forms a max-heap?

\par\vspace{6pt}\noindent
\begin{tabular}{@{}p{0.98\linewidth}@{}}
$\square$\ \textbf{(A)} 23, 17, 10, 6, 13, 14, 1, 5, 7, 12 \\
\addlinespace[3pt]
$\square$\ \textbf{(B)} 23, 17, 14, 7, 13, 10, 1, 5, 6, 12 \\
\addlinespace[3pt]
$\square$\ \textbf{(C)} 23, 17, 14, 6, 13, 10, 1, 5, 7, 15 \\
\addlinespace[3pt]
$\square$\ \textbf{(D)} 23, 14, 17, 1, 10, 13, 16, 12, 7, 5
\end{tabular}

\tcblower

\textbf{Official Key:} \textbf{\large B} \hfill \textit{\small [Easy]}
\vspace{2pt}\hrule height 0.3pt\vspace{3pt}
\textbf{Topic:} Programming and Data Structures \\
\textbf{Subtopic:} \textit{Trees, binary search trees, binary heaps}
\vspace{2pt}\hrule height 0.3pt\vspace{3pt}
\textbf{Guide / Hint:} \\
{\small In a 1-indexed max-heap, every node at index $i$ must satisfy $A[i] \ge A[2i]$ and $A[i] \ge A[2i+1]$. For sequence B $[23, 17, 14, 7, 13, 10, 1, 5, 6, 12]$: $A[1]=23 \ge A[2]=17, A[3]=14$; $A[2]=17 \ge A[4]=7, A[5]=13$; $A[3]=14 \ge A[6]=10, A[7]=1$; $A[4]=7 \ge A[8]=5, A[9]=6$; $A[5]=13 \ge A[10]=12$. Every parent is greater than or equal to its children, forming a valid max-heap.}

\end{tcolorbox}
\vspace{6pt}\hrule height 0.35pt\vspace{6pt}
\noindent
\begin{tcolorbox}[
    blanker,
    sidebyside,
    sidebyside align=top,
    lefthand width=0.55\textwidth,
    top=4pt, bottom=4pt, left=0pt, right=0pt,
    middle=16pt,
    underlay={
        \draw[dashed, black!80, line width=0.8pt, dash pattern=on 4pt off 3pt]
            ([xshift=0.55\textwidth+8pt]frame.north west) --
            ([xshift=0.55\textwidth+8pt]frame.south west);
    }
]
\textbf{Q.91} \quad {\footnotesize\textsc{[MCQ $\cdot$ 1 Mark(s) $\cdot$ Neg: 0.33]}} \hfill {\footnotesize\textbf{GATE 2023 CS} \quad \tiny\texttt{2023\_CS\_CS\_Q13}}
\vspace{2pt}\hrule height 0.3pt\vspace{4pt}

Let `SLLdel` be a function that deletes a node in a singly-linked list given a pointer to the node and a pointer to the head of the list. Similarly, let `DLLdel` be another function that deletes a node in a doubly-linked list given a pointer to the node and a pointer to the head of the list.
Let $n$ denote the number of nodes in each of the linked lists. Which one of the following choices is TRUE about the worst-case time complexity of `SLLdel` and `DLLdel`?

\par\vspace{6pt}\noindent
\begin{tabular}{@{}p{0.98\linewidth}@{}}
$\square$\ \textbf{(A)} `SLLdel` is $\mathcal{O}(1)$ and `DLLdel` is $\mathcal{O}(n)$ \\
\addlinespace[3pt]
$\square$\ \textbf{(B)} Both `SLLdel` and `DLLdel` are $\mathcal{O}(\log n)$ \\
\addlinespace[3pt]
$\square$\ \textbf{(C)} Both `SLLdel` and `DLLdel` are $\mathcal{O}(1)$ \\
\addlinespace[3pt]
$\square$\ \textbf{(D)} `SLLdel` is $\mathcal{O}(n)$ and `DLLdel` is $\mathcal{O}(1)$
\end{tabular}

\tcblower

\textbf{Official Key:} \textbf{\large D} \hfill \textit{\small [Easy]}
\vspace{2pt}\hrule height 0.3pt\vspace{3pt}
\textbf{Topic:} Programming and Data Structures \\
\textbf{Subtopic:} \textit{Arrays, stacks, queues, linked lists}
\vspace{2pt}\hrule height 0.3pt\vspace{3pt}
\textbf{Guide / Hint:} \\
{\small In a singly-linked list, deleting a given node requires finding its predecessor node to bypass it, requiring an $\mathcal{O}(n)$ traversal from the head in the worst case. In a doubly-linked list, each node directly points to its predecessor (`node->prev`), allowing immediate pointer adjustments in $\mathcal{O}(1)$ time.}

\end{tcolorbox}
\vspace{6pt}\hrule height 0.35pt\vspace{6pt}
\noindent
\begin{tcolorbox}[
    blanker,
    sidebyside,
    sidebyside align=top,
    lefthand width=0.55\textwidth,
    top=4pt, bottom=4pt, left=0pt, right=0pt,
    middle=16pt,
    underlay={
        \draw[dashed, black!80, line width=0.8pt, dash pattern=on 4pt off 3pt]
            ([xshift=0.55\textwidth+8pt]frame.north west) --
            ([xshift=0.55\textwidth+8pt]frame.south west);
    }
]
\textbf{Q.92} \quad {\footnotesize\textsc{[NAT $\cdot$ 1 Mark(s) $\cdot$ No Neg]}} \hfill {\footnotesize\textbf{GATE 2023 CS} \quad \tiny\texttt{2023\_CS\_CS\_Q35}}
\vspace{2pt}\hrule height 0.3pt\vspace{4pt}

The integer value printed by the ANSI-C program given below is \_\_\_\_\_\_\_\_\_\_\. (Answer in integer)
\begin{lstlisting}[language=C]
#include <stdio.h>
int funcp() {
    static int x = 1;
    x++;
    return x;
}
int main() {
    int x, y;
    x = funcp();
    y = funcp() + x;
    printf("%d\n", (x + y));
    return 0;
}
\end{lstlisting}

\par\vspace{6pt}\noindent\textbf{Your Answer:} \fbox{\rule{0pt}{1.3ex}\hspace{3cm}}

\tcblower

\textbf{Official Key:} \textbf{\large 7 to 7} \hfill \textit{\small [Easy]}
\vspace{2pt}\hrule height 0.3pt\vspace{3pt}
\textbf{Topic:} Programming and Data Structures \\
\textbf{Subtopic:} \textit{Programming in C}
\vspace{2pt}\hrule height 0.3pt\vspace{3pt}
\textbf{Guide / Hint:} \\
{\small In `funcp()`, the static variable $x$ persists across invocations. On first call `x = funcp()`, static $x$ increments from 1 to 2 and returns 2, setting local $x = 2$. In `y = funcp() + x`, `funcp()` increments static $x$ from 2 to 3 and returns 3, so $y = 3 + 2 = 5$. Finally, `printf` prints $x + y = 2 + 5 = 7$.}

\end{tcolorbox}
\vspace{6pt}\hrule height 0.35pt\vspace{6pt}
\noindent
\begin{tcolorbox}[
    blanker,
    sidebyside,
    sidebyside align=top,
    lefthand width=0.55\textwidth,
    top=4pt, bottom=4pt, left=0pt, right=0pt,
    middle=16pt,
    underlay={
        \draw[dashed, black!80, line width=0.8pt, dash pattern=on 4pt off 3pt]
            ([xshift=0.55\textwidth+8pt]frame.north west) --
            ([xshift=0.55\textwidth+8pt]frame.south west);
    }
]
\textbf{Q.93} \quad {\footnotesize\textsc{[MCQ $\cdot$ 2 Mark(s) $\cdot$ Neg: 0.66]}} \hfill {\footnotesize\textbf{GATE 2023 CS} \quad \tiny\texttt{2023\_CS\_CS\_Q46}}
\vspace{2pt}\hrule height 0.3pt\vspace{4pt}

Let $A$ be a priority queue for maintaining a set of elements. Suppose $A$ is implemented using a max-heap data structure. The operation $\text{Extract-Max}(A)$ extracts and deletes the maximum element from $A$. The operation $\text{Insert}(A, \text{key})$ inserts a new element $\text{key}$ in $A$. The properties of a max-heap are preserved at the end of each of these operations.
When $A$ contains $n$ elements, which one of the following statements about the worst case running time of these two operations is TRUE?

\par\vspace{6pt}\noindent
\begin{tabular}{@{}p{0.98\linewidth}@{}}
$\square$\ \textbf{(A)} Both `Extract-Max(A)` and `Insert(A,key)` run in $\mathcal{O}(1)$. \\
\addlinespace[3pt]
$\square$\ \textbf{(B)} Both `Extract-Max(A)` and `Insert(A,key)` run in $\mathcal{O}(\log n)$. \\
\addlinespace[3pt]
$\square$\ \textbf{(C)} `Extract-Max(A)` runs in $\mathcal{O}(1)$ whereas `Insert(A,key)` runs in $\mathcal{O}(n)$. \\
\addlinespace[3pt]
$\square$\ \textbf{(D)} `Extract-Max(A)` runs in $\mathcal{O}(1)$ whereas `Insert(A,key)` runs in $\mathcal{O}(\log n)$.
\end{tabular}

\tcblower

\textbf{Official Key:} \textbf{\large B} \hfill \textit{\small [Easy]}
\vspace{2pt}\hrule height 0.3pt\vspace{3pt}
\textbf{Topic:} Programming and Data Structures \\
\textbf{Subtopic:} \textit{Trees, binary search trees, binary heaps}
\vspace{2pt}\hrule height 0.3pt\vspace{3pt}
\textbf{Guide / Hint:} \\
{\small In a binary max-heap of $n$ elements, `Extract-Max` replaces the root with the last element and performs `maxHeapify` downward along the height $\lfloor \log_2 n \rfloor$, running in $\mathcal{O}(\log n)$ worst-case time. `Insert` places the new key at the next leaf position and sifts upward along the height, also running in $\mathcal{O}(\log n)$ worst-case time.}

\end{tcolorbox}
\vspace{6pt}\hrule height 0.35pt\vspace{6pt}
\noindent
\begin{tcolorbox}[
    blanker,
    sidebyside,
    sidebyside align=top,
    lefthand width=0.55\textwidth,
    top=4pt, bottom=4pt, left=0pt, right=0pt,
    middle=16pt,
    underlay={
        \draw[dashed, black!80, line width=0.8pt, dash pattern=on 4pt off 3pt]
            ([xshift=0.55\textwidth+8pt]frame.north west) --
            ([xshift=0.55\textwidth+8pt]frame.south west);
    }
]
\textbf{Q.94} \quad {\footnotesize\textsc{[MCQ $\cdot$ 1 Mark(s) $\cdot$ Neg: 0.33]}} \hfill {\footnotesize\textbf{GATE 2024 CS1} \quad \tiny\texttt{2024\_CS1\_CS\_Q18}}
\vspace{2pt}\hrule height 0.3pt\vspace{4pt}

Consider the following C program:
\begin{lstlisting}[language=C]
#include <stdio.h>
int main() {
    int a = 6;
    int b = 0;
    while(a < 10) {
        a = a / 12 + 1;
        a += b;
    }
    printf("%d", a);
    return 0;
}
\end{lstlisting}
Which one of the following statements is CORRECT?

\par\vspace{6pt}\noindent
\begin{tabular}{@{}p{0.98\linewidth}@{}}
$\square$\ \textbf{(A)} The program prints 9 as output \\
\addlinespace[3pt]
$\square$\ \textbf{(B)} The program prints 10 as output \\
\addlinespace[3pt]
$\square$\ \textbf{(C)} The program gets stuck in an infinite loop \\
\addlinespace[3pt]
$\square$\ \textbf{(D)} The program prints 6 as output
\end{tabular}

\tcblower

\textbf{Official Key:} \textbf{\large C} \hfill \textit{\small [Easy]}
\vspace{2pt}\hrule height 0.3pt\vspace{3pt}
\textbf{Topic:} Programming and Data Structures \\
\textbf{Subtopic:} \textit{Programming in C}
\vspace{2pt}\hrule height 0.3pt\vspace{3pt}
\textbf{Guide / Hint:} \\
{\small Initially $a = 6, b = 0$. In the loop, integer division $a/12 = 6/12 = 0$, so $a = 0 + 1 = 1$, and $a += b \implies a = 1$. In subsequent iterations, $a=1 < 10$, $1/12 + 1 = 1$, so $a$ remains 1 in every iteration. The loop condition $a < 10$ remains true indefinitely, causing an infinite loop.}

\end{tcolorbox}
\vspace{6pt}\hrule height 0.35pt\vspace{6pt}
\noindent
\begin{tcolorbox}[
    blanker,
    sidebyside,
    sidebyside align=top,
    lefthand width=0.55\textwidth,
    top=4pt, bottom=4pt, left=0pt, right=0pt,
    middle=16pt,
    underlay={
        \draw[dashed, black!80, line width=0.8pt, dash pattern=on 4pt off 3pt]
            ([xshift=0.55\textwidth+8pt]frame.north west) --
            ([xshift=0.55\textwidth+8pt]frame.south west);
    }
]
\textbf{Q.95} \quad {\footnotesize\textsc{[MCQ $\cdot$ 1 Mark(s) $\cdot$ Neg: 0.33]}} \hfill {\footnotesize\textbf{GATE 2024 CS1} \quad \tiny\texttt{2024\_CS1\_CS\_Q19}}
\vspace{2pt}\hrule height 0.3pt\vspace{4pt}

Consider the following C program:
\begin{lstlisting}[language=C]
#include <stdio.h>
void fX();
int main() {
    fX();
    return 0;
}
void fX() {
    char a;
    if((a = getchar()) != '\n')
        fX();
    if(a != '\n')
        putchar(a);
}
\end{lstlisting}
Assume that the input to the program from the command line is `1234` followed by a newline character. Which one of the following statements is CORRECT?

\par\vspace{6pt}\noindent
\begin{tabular}{@{}p{0.98\linewidth}@{}}
$\square$\ \textbf{(A)} The program will not terminate \\
\addlinespace[3pt]
$\square$\ \textbf{(B)} The program will terminate with no output \\
\addlinespace[3pt]
$\square$\ \textbf{(C)} The program will terminate with 4321 as output \\
\addlinespace[3pt]
$\square$\ \textbf{(D)} The program will terminate with 1234 as output
\end{tabular}

\tcblower

\textbf{Official Key:} \textbf{\large C} \hfill \textit{\small [Easy]}
\vspace{2pt}\hrule height 0.3pt\vspace{3pt}
\textbf{Topic:} Programming and Data Structures \\
\textbf{Subtopic:} \textit{Recursion}
\vspace{2pt}\hrule height 0.3pt\vspace{3pt}
\textbf{Guide / Hint:} \\
{\small The function `fX()` recursively reads each character before printing. As each character (`'1'`, `'2'`, `'3'`, `'4'`) is read, it is saved on the call stack before `fX()` recurses. When `\n` is reached, recursion terminates and the characters are printed during stack unwinding in reverse (LIFO) order as `4321`.}

\end{tcolorbox}
\vspace{6pt}\hrule height 0.35pt\vspace{6pt}
\noindent
\begin{tcolorbox}[
    blanker,
    sidebyside,
    sidebyside align=top,
    lefthand width=0.55\textwidth,
    top=4pt, bottom=4pt, left=0pt, right=0pt,
    middle=16pt,
    underlay={
        \draw[dashed, black!80, line width=0.8pt, dash pattern=on 4pt off 3pt]
            ([xshift=0.55\textwidth+8pt]frame.north west) --
            ([xshift=0.55\textwidth+8pt]frame.south west);
    }
]
\textbf{Q.96} \quad {\footnotesize\textsc{[MSQ $\cdot$ 1 Mark(s) $\cdot$ No Neg]}} \hfill {\footnotesize\textbf{GATE 2024 CS2} \quad \tiny\texttt{2024\_CS2\_CS\_Q33}}
\vspace{2pt}\hrule height 0.3pt\vspace{4pt}

Consider the following C function definition:
\begin{verbatim}
int fX(char *a){
    char *b = a;
    while(*b)
        b++;
    return b - a;
}
\end{verbatim}
Which of the following statements is/are TRUE?

\par\vspace{6pt}\noindent
\begin{tabular}{@{}p{0.98\linewidth}@{}}
$\square$\ \textbf{(A)} The function call fX("abcd") will always return a value \\
\addlinespace[3pt]
$\square$\ \textbf{(B)} Assuming a character array c is declared as char c[] = "abcd" in main(), the function call fX(c) will always return a value \\
\addlinespace[3pt]
$\square$\ \textbf{(C)} The code of the function will not compile \\
\addlinespace[3pt]
$\square$\ \textbf{(D)} Assuming a character pointer c is declared as char *c = "abcd" in main(), the function call fX(c) will always return a value
\end{tabular}

\tcblower

\textbf{Official Key:} \textbf{\large A;B;D} \hfill \textit{\small [Easy]}
\vspace{2pt}\hrule height 0.3pt\vspace{3pt}
\textbf{Topic:} Programming and Data Structures \\
\textbf{Subtopic:} \textit{Programming in C}
\vspace{2pt}\hrule height 0.3pt\vspace{3pt}
\textbf{Guide / Hint:} \\
{\small The function `fX` calculates the string length by advancing pointer `b` until the null terminator `\0` is reached, then returning pointer difference `b - a`. String literals in C are null-terminated, so calls with string literals or character arrays (`fX("abcd")`, `fX(c)`) correctly compute 4 and return (A, B, and D are true).}

\end{tcolorbox}
\vspace{6pt}\hrule height 0.35pt\vspace{6pt}
\noindent
\begin{tcolorbox}[
    blanker,
    sidebyside,
    sidebyside align=top,
    lefthand width=0.55\textwidth,
    top=4pt, bottom=4pt, left=0pt, right=0pt,
    middle=16pt,
    underlay={
        \draw[dashed, black!80, line width=0.8pt, dash pattern=on 4pt off 3pt]
            ([xshift=0.55\textwidth+8pt]frame.north west) --
            ([xshift=0.55\textwidth+8pt]frame.south west);
    }
]
\textbf{Q.97} \quad {\footnotesize\textsc{[MCQ $\cdot$ 2 Mark(s) $\cdot$ Neg: 0.66]}} \hfill {\footnotesize\textbf{GATE 2024 CS2} \quad \tiny\texttt{2024\_CS2\_CS\_Q36}}
\vspace{2pt}\hrule height 0.3pt\vspace{4pt}

What is the output of the following C program?
\begin{verbatim}
#include <stdio.h>
int main() {
    double a[2] = {20.0, 25.0}, *p, *q;
    p = a;
    q = p + 1;
    printf("%d,%d", (int)(q - p), (int)(*q - *p));
    return 0;
}
\end{verbatim}

\par\vspace{6pt}\noindent
\begin{tabular}{@{}p{0.48\linewidth}@{\hspace{0.04\linewidth}}p{0.48\linewidth}@{}}
$\square$\ \textbf{(A)} 4,8 & $\square$\ \textbf{(B)} 1,5 \\
\addlinespace[3pt]
$\square$\ \textbf{(C)} 8,5 & $\square$\ \textbf{(D)} 1,8
\end{tabular}

\tcblower

\textbf{Official Key:} \textbf{\large B} \hfill \textit{\small [Easy]}
\vspace{2pt}\hrule height 0.3pt\vspace{3pt}
\textbf{Topic:} Programming and Data Structures \\
\textbf{Subtopic:} \textit{Programming in C}
\vspace{2pt}\hrule height 0.3pt\vspace{3pt}
\textbf{Guide / Hint:} \\
{\small In pointer arithmetic, pointer subtraction `q - p` yields the number of elements between pointers, so $(q - p) = 1$. The expression `*q - *p` calculates the difference between pointed double values: $25.0 - 20.0 = 5.0$, cast to integer is 5. Thus, `1,5` is printed.}

\end{tcolorbox}
\vspace{6pt}\hrule height 0.35pt\vspace{6pt}
\noindent
\begin{tcolorbox}[
    blanker,
    sidebyside,
    sidebyside align=top,
    lefthand width=0.55\textwidth,
    top=4pt, bottom=4pt, left=0pt, right=0pt,
    middle=16pt,
    underlay={
        \draw[dashed, black!80, line width=0.8pt, dash pattern=on 4pt off 3pt]
            ([xshift=0.55\textwidth+8pt]frame.north west) --
            ([xshift=0.55\textwidth+8pt]frame.south west);
    }
]
\textbf{Q.98} \quad {\footnotesize\textsc{[MCQ $\cdot$ 2 Mark(s) $\cdot$ Neg: 0.66]}} \hfill {\footnotesize\textbf{GATE 2024 CS2} \quad \tiny\texttt{2024\_CS2\_CS\_Q39}}
\vspace{2pt}\hrule height 0.3pt\vspace{4pt}

You are given a set $V$ of distinct integers. A binary search tree $T$ is created by inserting all elements of $V$ one by one, starting with an empty tree. The tree $T$ follows the convention that, at each node, all values stored in the left subtree of the node are smaller than the value stored at the node. Which of the following is TRUE?

\par\vspace{6pt}\noindent
\begin{tabular}{@{}p{0.98\linewidth}@{}}
$\square$\ \textbf{(A)} Inorder traversal of $T$ can be determined from $V$ \\
\addlinespace[3pt]
$\square$\ \textbf{(B)} Root node of $T$ can be determined from $V$ \\
\addlinespace[3pt]
$\square$\ \textbf{(C)} Preorder traversal of $T$ can be determined from $V$ \\
\addlinespace[3pt]
$\square$\ \textbf{(D)} Postorder traversal of $T$ can be determined from $V$
\end{tabular}

\tcblower

\textbf{Official Key:} \textbf{\large A} \hfill \textit{\small [Easy]}
\vspace{2pt}\hrule height 0.3pt\vspace{3pt}
\textbf{Topic:} Programming and Data Structures \\
\textbf{Subtopic:} \textit{Trees, binary search trees, binary heaps}
\vspace{2pt}\hrule height 0.3pt\vspace{3pt}
\textbf{Guide / Hint:} \\
{\small For any binary search tree, the inorder traversal sequence always visits elements in strictly increasing sorted order regardless of insertion order or tree shape. Since the set of distinct elements $V$ is known, sorting $V$ directly determines the unique inorder traversal.}

\end{tcolorbox}
\vspace{6pt}\hrule height 0.35pt\vspace{6pt}
\noindent
\begin{tcolorbox}[
    blanker,
    sidebyside,
    sidebyside align=top,
    lefthand width=0.55\textwidth,
    top=4pt, bottom=4pt, left=0pt, right=0pt,
    middle=16pt,
    underlay={
        \draw[dashed, black!80, line width=0.8pt, dash pattern=on 4pt off 3pt]
            ([xshift=0.55\textwidth+8pt]frame.north west) --
            ([xshift=0.55\textwidth+8pt]frame.south west);
    }
]
\textbf{Q.99} \quad {\footnotesize\textsc{[MSQ $\cdot$ 1 Mark(s) $\cdot$ No Neg]}} \hfill {\footnotesize\textbf{GATE 2025 CS1} \quad \tiny\texttt{2025\_CS1\_CS\_Q26}}
\vspace{2pt}\hrule height 0.3pt\vspace{4pt}

Which of the following statement(s) is/are TRUE for any binary search tree (BST) having $n$ distinct integers?

\par\vspace{6pt}\noindent
\begin{tabular}{@{}p{0.98\linewidth}@{}}
$\square$\ \textbf{(A)} The maximum length of a path from the root node to any other node is $(n-1)$. \\
\addlinespace[3pt]
$\square$\ \textbf{(B)} An inorder traversal will always produce a sorted sequence of elements. \\
\addlinespace[3pt]
$\square$\ \textbf{(C)} Finding an element takes $O(\log_2 n)$ time in the worst case. \\
\addlinespace[3pt]
$\square$\ \textbf{(D)} Every BST is also a Min-Heap.
\end{tabular}

\tcblower

\textbf{Official Key:} \textbf{\large A;B} \hfill \textit{\small [Easy]}
\vspace{2pt}\hrule height 0.3pt\vspace{3pt}
\textbf{Topic:} Programming and Data Structures \\
\textbf{Subtopic:} \textit{Trees, binary search trees, binary heaps}
\vspace{2pt}\hrule height 0.3pt\vspace{3pt}
\textbf{Guide / Hint:} \\
{\small In a degenerate (skewed) binary search tree, the path length from root to the farthest leaf is $n - 1$ edges (A is true). Inorder traversal of any valid BST visits nodes in strictly increasing order (B is true). Worst-case search in an unbalanced BST takes $O(n)$ time, and BST property differs fundamentally from heap order.}

\end{tcolorbox}
\vspace{6pt}\hrule height 0.35pt\vspace{6pt}
\noindent
\begin{tcolorbox}[
    blanker,
    sidebyside,
    sidebyside align=top,
    lefthand width=0.55\textwidth,
    top=4pt, bottom=4pt, left=0pt, right=0pt,
    middle=16pt,
    underlay={
        \draw[dashed, black!80, line width=0.8pt, dash pattern=on 4pt off 3pt]
            ([xshift=0.55\textwidth+8pt]frame.north west) --
            ([xshift=0.55\textwidth+8pt]frame.south west);
    }
]
\textbf{Q.100} \quad {\footnotesize\textsc{[NAT $\cdot$ 1 Mark(s) $\cdot$ No Neg]}} \hfill {\footnotesize\textbf{GATE 2025 CS1} \quad \tiny\texttt{2025\_CS1\_CS\_Q34}}
\vspace{2pt}\hrule height 0.3pt\vspace{4pt}

Consider the following C program:
\par\medskip
\begin{verbatim}
#include <stdio.h>
void foo(int *p, int x){
    *p = x;
}
int main(){
    int *z;
    int a = 20, b = 25;
    z = &a;
    foo(z, b);
    printf("%d", a);
    return 0;
}
\end{verbatim}
The output of the given C program is \_\_\_\_\_\_\. (Answer in integer)

\par\vspace{6pt}\noindent\textbf{Your Answer:} \fbox{\rule{0pt}{1.3ex}\hspace{3cm}}

\tcblower

\textbf{Official Key:} \textbf{\large 25 to 25} \hfill \textit{\small [Easy]}
\vspace{2pt}\hrule height 0.3pt\vspace{3pt}
\textbf{Topic:} Programming and Data Structures \\
\textbf{Subtopic:} \textit{Programming in C}
\vspace{2pt}\hrule height 0.3pt\vspace{3pt}
\textbf{Guide / Hint:} \\
{\small In `main()`, pointer `z` is assigned the memory address of variable `a` (`\&a`). Function `foo(z, b)` dereferences `p` (which points to `a`) and assigns `x = b = 25` to `*p`. Thus, `a` is updated to 25, which `printf` prints.}

\end{tcolorbox}
\vspace{6pt}\hrule height 0.35pt\vspace{6pt}
\noindent
\begin{tcolorbox}[
    blanker,
    sidebyside,
    sidebyside align=top,
    lefthand width=0.55\textwidth,
    top=4pt, bottom=4pt, left=0pt, right=0pt,
    middle=16pt,
    underlay={
        \draw[dashed, black!80, line width=0.8pt, dash pattern=on 4pt off 3pt]
            ([xshift=0.55\textwidth+8pt]frame.north west) --
            ([xshift=0.55\textwidth+8pt]frame.south west);
    }
]
\textbf{Q.101} \quad {\footnotesize\textsc{[NAT $\cdot$ 1 Mark(s) $\cdot$ No Neg]}} \hfill {\footnotesize\textbf{GATE 2025 CS1} \quad \tiny\texttt{2025\_CS1\_CS\_Q35}}
\vspace{2pt}\hrule height 0.3pt\vspace{4pt}

The height of any rooted tree is defined as the maximum number of edges in the path from the root node to any leaf node. Suppose a Min-Heap $T$ stores 32 keys. The height of $T$ is \_\_\_\_\_\_\. (Answer in integer)

\par\vspace{6pt}\noindent\textbf{Your Answer:} \fbox{\rule{0pt}{1.3ex}\hspace{3cm}}

\tcblower

\textbf{Official Key:} \textbf{\large 5 to 5} \hfill \textit{\small [Easy]}
\vspace{2pt}\hrule height 0.3pt\vspace{3pt}
\textbf{Topic:} Programming and Data Structures \\
\textbf{Subtopic:} \textit{Trees, binary search trees, binary heaps}
\vspace{2pt}\hrule height 0.3pt\vspace{3pt}
\textbf{Guide / Hint:} \\
{\small A binary heap storing $n$ keys is an almost complete binary tree whose height in terms of edges is $\lfloor \log_2 n \rfloor$. For $n = 32$, the height is $\lfloor \log_2 32 \rfloor = 5$.}

\end{tcolorbox}
\vspace{6pt}\hrule height 0.35pt\vspace{6pt}
\noindent
\begin{tcolorbox}[
    blanker,
    sidebyside,
    sidebyside align=top,
    lefthand width=0.55\textwidth,
    top=4pt, bottom=4pt, left=0pt, right=0pt,
    middle=16pt,
    underlay={
        \draw[dashed, black!80, line width=0.8pt, dash pattern=on 4pt off 3pt]
            ([xshift=0.55\textwidth+8pt]frame.north west) --
            ([xshift=0.55\textwidth+8pt]frame.south west);
    }
]
\textbf{Q.102} \quad {\footnotesize\textsc{[MCQ $\cdot$ 1 Mark(s) $\cdot$ Neg: 0.33]}} \hfill {\footnotesize\textbf{GATE 2025 CS2} \quad \tiny\texttt{2025\_CS2\_CS\_Q13}}
\vspace{2pt}\hrule height 0.3pt\vspace{4pt}

Consider a binary tree $T$ in which every node has either zero or two children. Let $n > 0$ be the number of nodes in $T$. Which ONE of the following is the number of nodes in $T$ that have exactly two children?

\par\vspace{6pt}\noindent
\begin{tabular}{@{}p{0.48\linewidth}@{\hspace{0.04\linewidth}}p{0.48\linewidth}@{}}
$\square$\ \textbf{(A)} $\frac{n-2}{2}$ & $\square$\ \textbf{(B)} $\frac{n-1}{2}$ \\
\addlinespace[3pt]
$\square$\ \textbf{(C)} $\frac{n}{2}$ & $\square$\ \textbf{(D)} $\frac{n+1}{2}$
\end{tabular}

\tcblower

\textbf{Official Key:} \textbf{\large B} \hfill \textit{\small [Easy]}
\vspace{2pt}\hrule height 0.3pt\vspace{3pt}
\textbf{Topic:} Programming and Data Structures \\
\textbf{Subtopic:} \textit{Trees and strictly binary trees}
\vspace{2pt}\hrule height 0.3pt\vspace{3pt}
\textbf{Guide / Hint:} \\
{\small In any proper (strictly) binary tree with $L$ leaves and $I$ internal nodes (nodes with 2 children), $L = I + 1$. Since $n = L + I = 2I + 1$, the number of internal nodes is $I = \frac{n-1}{2}$.}

\end{tcolorbox}
\vspace{6pt}\hrule height 0.35pt\vspace{6pt}
\noindent
\begin{tcolorbox}[
    blanker,
    sidebyside,
    sidebyside align=top,
    lefthand width=0.55\textwidth,
    top=4pt, bottom=4pt, left=0pt, right=0pt,
    middle=16pt,
    underlay={
        \draw[dashed, black!80, line width=0.8pt, dash pattern=on 4pt off 3pt]
            ([xshift=0.55\textwidth+8pt]frame.north west) --
            ([xshift=0.55\textwidth+8pt]frame.south west);
    }
]
\textbf{Q.103} \quad {\footnotesize\textsc{[NAT $\cdot$ 1 Mark(s) $\cdot$ No Neg]}} \hfill {\footnotesize\textbf{GATE 2025 CS2} \quad \tiny\texttt{2025\_CS2\_CS\_Q33}}
\vspace{2pt}\hrule height 0.3pt\vspace{4pt}

Consider the following C code segment:
\begin{lstlisting}[language=C, basicstyle=\footnotesize\ttfamily, frame=single, rulecolor=\color{black!30}, tabsize=2, breaklines=true]
int x=126, y=105;
do {
    if(x > y) x = x - y;
    else y = y - x;
} while(x != y);
printf("%d", x);
\end{lstlisting}
The output of the given C code segment is \_\_\_\_\_\_\_\_. (Answer in integer)

\par\vspace{6pt}\noindent\textbf{Your Answer:} \fbox{\rule{0pt}{1.3ex}\hspace{3cm}}

\tcblower

\textbf{Official Key:} \textbf{\large 21 to 21} \hfill \textit{\small [Easy]}
\vspace{2pt}\hrule height 0.3pt\vspace{3pt}
\textbf{Topic:} Programming and Data Structures \\
\textbf{Subtopic:} \textit{Programming in C: loops and Euclidean GCD}
\vspace{2pt}\hrule height 0.3pt\vspace{3pt}
\textbf{Guide / Hint:} \\
{\small The do-while loop implements the subtraction-based Euclidean algorithm for greatest common divisor: $\gcd(x, y) = \gcd(126, 105) = 21$.}

\end{tcolorbox}
\vspace{6pt}\hrule height 0.35pt\vspace{6pt}
\noindent
\begin{tcolorbox}[
    blanker,
    sidebyside,
    sidebyside align=top,
    lefthand width=0.55\textwidth,
    top=4pt, bottom=4pt, left=0pt, right=0pt,
    middle=16pt,
    underlay={
        \draw[dashed, black!80, line width=0.8pt, dash pattern=on 4pt off 3pt]
            ([xshift=0.55\textwidth+8pt]frame.north west) --
            ([xshift=0.55\textwidth+8pt]frame.south west);
    }
]
\textbf{Q.104} \quad {\footnotesize\textsc{[NAT $\cdot$ 1 Mark(s) $\cdot$ No Neg]}} \hfill {\footnotesize\textbf{GATE 2025 CS2} \quad \tiny\texttt{2025\_CS2\_CS\_Q35}}
\vspace{2pt}\hrule height 0.3pt\vspace{4pt}

Suppose the values 10, $-4$, 15, 30, 20, 5, 60, 19 are inserted in that order into an initially empty binary search tree. Let $T$ be the resulting binary search tree. The number of edges in the path from the node containing 19 to the root node of $T$ is \_\_\_\_\_\_\_\_\_\_\_. (Answer in integer)

\par\vspace{6pt}\noindent\textbf{Your Answer:} \fbox{\rule{0pt}{1.3ex}\hspace{3cm}}

\tcblower

\textbf{Official Key:} \textbf{\large 4 to 4} \hfill \textit{\small [Easy]}
\vspace{2pt}\hrule height 0.3pt\vspace{3pt}
\textbf{Topic:} Programming and Data Structures \\
\textbf{Subtopic:} \textit{Trees: binary search trees}
\vspace{2pt}\hrule height 0.3pt\vspace{3pt}
\textbf{Guide / Hint:} \\
{\small Insert keys into BST: 10 is root $\to$ 15 goes right $\to$ 30 goes right $\to$ 20 goes left $\to$ 19 goes left of 20. The path from root is $10 \to 15 \to 30 \to 20 \to 19$, which contains 4 edges.}

\end{tcolorbox}
\vspace{6pt}\hrule height 0.35pt\vspace{6pt}
\vspace{8pt}
\subsection*{Algorithms}
\vspace{-4pt}
\noindent
\begin{tcolorbox}[
    blanker,
    sidebyside,
    sidebyside align=top,
    lefthand width=0.55\textwidth,
    top=4pt, bottom=4pt, left=0pt, right=0pt,
    middle=16pt,
    underlay={
        \draw[dashed, black!80, line width=0.8pt, dash pattern=on 4pt off 3pt]
            ([xshift=0.55\textwidth+8pt]frame.north west) --
            ([xshift=0.55\textwidth+8pt]frame.south west);
    }
]
\textbf{Q.105} \quad {\footnotesize\textsc{[MCQ $\cdot$ 1 Mark(s) $\cdot$ Neg: 0.33]}} \hfill {\footnotesize\textbf{GATE 2022 CS} \quad \tiny\texttt{2022\_CS\_CS\_Q16}}
\vspace{2pt}\hrule height 0.3pt\vspace{4pt}

Suppose we are given $n$ keys, $m$ hash table slots, and two simple uniform hash functions $h_1$ and $h_2$. Further suppose our hashing scheme uses $h_1$ for the odd keys and $h_2$ for the even keys. What is the expected number of keys in a slot?

\par\vspace{6pt}\noindent
\begin{tabular}{@{}p{0.48\linewidth}@{\hspace{0.04\linewidth}}p{0.48\linewidth}@{}}
$\square$\ \textbf{(A)} $m/n$ & $\square$\ \textbf{(B)} $n/m$ \\
\addlinespace[3pt]
$\square$\ \textbf{(C)} $2n/m$ & $\square$\ \textbf{(D)} $n / (2m)$
\end{tabular}

\tcblower

\textbf{Official Key:} \textbf{\large B} \hfill \textit{\small [Easy]}
\vspace{2pt}\hrule height 0.3pt\vspace{3pt}
\textbf{Topic:} Algorithms \\
\textbf{Subtopic:} \textit{Searching, sorting, hashing}
\vspace{2pt}\hrule height 0.3pt\vspace{3pt}
\textbf{Guide / Hint:} \\
{\small Let $n_1$ be the number of odd keys and $n_2$ be the number of even keys ($n_1 + n_2 = n$). Under simple uniform hashing, each key independently hashes to any slot with probability $1/m$. By linearity of expectation, the expected number of keys in any slot is $\frac{n_1}{m} + \frac{n_2}{m} = \frac{n}{m}$.}

\end{tcolorbox}
\vspace{6pt}\hrule height 0.35pt\vspace{6pt}
\noindent
\begin{tcolorbox}[
    blanker,
    sidebyside,
    sidebyside align=top,
    lefthand width=0.55\textwidth,
    top=4pt, bottom=4pt, left=0pt, right=0pt,
    middle=16pt,
    underlay={
        \draw[dashed, black!80, line width=0.8pt, dash pattern=on 4pt off 3pt]
            ([xshift=0.55\textwidth+8pt]frame.north west) --
            ([xshift=0.55\textwidth+8pt]frame.south west);
    }
]
\textbf{Q.106} \quad {\footnotesize\textsc{[MCQ $\cdot$ 1 Mark(s) $\cdot$ Neg: 0.33]}} \hfill {\footnotesize\textbf{GATE 2023 CS} \quad \tiny\texttt{2023\_CS\_CS\_Q20}}
\vspace{2pt}\hrule height 0.3pt\vspace{4pt}

An algorithm has to store several keys generated by an adversary in a hash table. The adversary is malicious who tries to maximize the number of collisions. Let $k$ be the number of keys, $m$ be the number of slots in the hash table, and $k > m$.
Which one of the following is the best hashing strategy to counteract the adversary?

\par\vspace{6pt}\noindent
\begin{tabular}{@{}p{0.98\linewidth}@{}}
$\square$\ \textbf{(A)} Division method, i.e., use the hash function $h(k) = k \bmod m$. \\
\addlinespace[3pt]
$\square$\ \textbf{(B)} Multiplication method, i.e., use the hash function $h(k) = \lfloor m(kA - \lfloor kA \rfloor)\rfloor$, where $A$ is a carefully chosen constant. \\
\addlinespace[3pt]
$\square$\ \textbf{(C)} Universal hashing method. \\
\addlinespace[3pt]
$\square$\ \textbf{(D)} If $k$ is a prime number, use Division method. Otherwise, use Multiplication method.
\end{tabular}

\tcblower

\textbf{Official Key:} \textbf{\large C} \hfill \textit{\small [Easy]}
\vspace{2pt}\hrule height 0.3pt\vspace{3pt}
\textbf{Topic:} Algorithms \\
\textbf{Subtopic:} \textit{Searching, sorting, hashing}
\vspace{2pt}\hrule height 0.3pt\vspace{3pt}
\textbf{Guide / Hint:} \\
{\small In universal hashing, the hash function is chosen uniformly at random from a carefully designed universal collection of hash functions independently of the keys. Because the adversary cannot predict which function will be selected at runtime, any fixed choice of malicious keys cannot guarantee worst-case collision behavior.}

\end{tcolorbox}
\vspace{6pt}\hrule height 0.35pt\vspace{6pt}
\noindent
\begin{tcolorbox}[
    blanker,
    sidebyside,
    sidebyside align=top,
    lefthand width=0.55\textwidth,
    top=4pt, bottom=4pt, left=0pt, right=0pt,
    middle=16pt,
    underlay={
        \draw[dashed, black!80, line width=0.8pt, dash pattern=on 4pt off 3pt]
            ([xshift=0.55\textwidth+8pt]frame.north west) --
            ([xshift=0.55\textwidth+8pt]frame.south west);
    }
]
\textbf{Q.107} \quad {\footnotesize\textsc{[MSQ $\cdot$ 1 Mark(s) $\cdot$ No Neg]}} \hfill {\footnotesize\textbf{GATE 2023 CS} \quad \tiny\texttt{2023\_CS\_CS\_Q29}}
\vspace{2pt}\hrule height 0.3pt\vspace{4pt}

Let $f$ and $g$ be functions of natural numbers given by $f(n) = n$ and $g(n) = n^2$.
Which of the following statements is/are TRUE?

\par\vspace{6pt}\noindent
\begin{tabular}{@{}p{0.48\linewidth}@{\hspace{0.04\linewidth}}p{0.48\linewidth}@{}}
$\square$\ \textbf{(A)} $f \in \mathcal{O}(g)$ & $\square$\ \textbf{(B)} $f \in \Omega(g)$ \\
\addlinespace[3pt]
$\square$\ \textbf{(C)} $f \in o(g)$ & $\square$\ \textbf{(D)} $f \in \Theta(g)$
\end{tabular}

\tcblower

\textbf{Official Key:} \textbf{\large A, C} \hfill \textit{\small [Easy]}
\vspace{2pt}\hrule height 0.3pt\vspace{3pt}
\textbf{Topic:} Algorithms \\
\textbf{Subtopic:} \textit{Asymptotic worst case time and space complexity}
\vspace{2pt}\hrule height 0.3pt\vspace{3pt}
\textbf{Guide / Hint:} \\
{\small For $f(n) = n$ and $g(n) = n^2$, we evaluate $\lim_{n \to \infty} \frac{f(n)}{g(n)} = \lim_{n \to \infty} \frac{1}{n} = 0$. By definition, if the limit is 0, $f(n) \in o(g(n))$, which strictly implies $f(n) \in \mathcal{O}(g(n))$. Hence, both A and C are TRUE.}

\end{tcolorbox}
\vspace{6pt}\hrule height 0.35pt\vspace{6pt}
\noindent
\begin{tcolorbox}[
    blanker,
    sidebyside,
    sidebyside align=top,
    lefthand width=0.55\textwidth,
    top=4pt, bottom=4pt, left=0pt, right=0pt,
    middle=16pt,
    underlay={
        \draw[dashed, black!80, line width=0.8pt, dash pattern=on 4pt off 3pt]
            ([xshift=0.55\textwidth+8pt]frame.north west) --
            ([xshift=0.55\textwidth+8pt]frame.south west);
    }
]
\textbf{Q.108} \quad {\footnotesize\textsc{[MSQ $\cdot$ 2 Mark(s) $\cdot$ No Neg]}} \hfill {\footnotesize\textbf{GATE 2023 CS} \quad \tiny\texttt{2023\_CS\_CS\_Q54}}
\vspace{2pt}\hrule height 0.3pt\vspace{4pt}

Consider functions Function 1 and Function 2 expressed in pseudocode as follows:
\begin{lstlisting}
// Function 1
while n > 1 do
    for i = 1 to n do
        x = x + 1;
    end for
    n = floor(n / 2);
end while
\par\medskip
// Function 2
for i = 1 to 100 * n do
    x = x + 1;
end for
\end{lstlisting}
Let $f_1(n)$ and $f_2(n)$ denote the number of times the statement ``$x = x + 1$'' is executed in Function 1 and Function 2, respectively.
Which of the following statements is/are TRUE?

\par\vspace{6pt}\noindent
\begin{tabular}{@{}p{0.48\linewidth}@{\hspace{0.04\linewidth}}p{0.48\linewidth}@{}}
$\square$\ \textbf{(A)} $f_1(n) \in \Theta(f_2(n))$ & $\square$\ \textbf{(B)} $f_1(n) \in o(f_2(n))$ \\
\addlinespace[3pt]
$\square$\ \textbf{(C)} $f_1(n) \in \omega(f_2(n))$ & $\square$\ \textbf{(D)} $f_1(n) \in \mathcal{O}(n)$
\end{tabular}

\tcblower

\textbf{Official Key:} \textbf{\large A, D} \hfill \textit{\small [Easy]}
\vspace{2pt}\hrule height 0.3pt\vspace{3pt}
\textbf{Topic:} Algorithms \\
\textbf{Subtopic:} \textit{Asymptotic worst case time and space complexity}
\vspace{2pt}\hrule height 0.3pt\vspace{3pt}
\textbf{Guide / Hint:} \\
{\small In Function 1, the inner loop executes $n + \lfloor n/2 \rfloor + \lfloor n/4 \rfloor + \dots \approx 2n$ times, which gives $f_1(n) = \Theta(n)$. In Function 2, the loop runs $100n$ times, which gives $f_2(n) = \Theta(n)$. Because both functions have exact linear asymptotic growth, $f_1(n) \in \Theta(f_2(n))$ (A is TRUE) and $f_1(n) \in \mathcal{O}(n)$ (D is TRUE).}

\end{tcolorbox}
\vspace{6pt}\hrule height 0.35pt\vspace{6pt}
\noindent
\begin{tcolorbox}[
    blanker,
    sidebyside,
    sidebyside align=top,
    lefthand width=0.55\textwidth,
    top=4pt, bottom=4pt, left=0pt, right=0pt,
    middle=16pt,
    underlay={
        \draw[dashed, black!80, line width=0.8pt, dash pattern=on 4pt off 3pt]
            ([xshift=0.55\textwidth+8pt]frame.north west) --
            ([xshift=0.55\textwidth+8pt]frame.south west);
    }
]
\textbf{Q.109} \quad {\footnotesize\textsc{[MCQ $\cdot$ 1 Mark(s) $\cdot$ Neg: 0.33]}} \hfill {\footnotesize\textbf{GATE 2024 CS1} \quad \tiny\texttt{2024\_CS1\_CS\_Q17}}
\vspace{2pt}\hrule height 0.3pt\vspace{4pt}

Given an integer array of size $N$, we want to check if the array is sorted (in either ascending or descending order). An algorithm solves this problem by making a single pass through the array and comparing each element of the array only with its adjacent elements. The worst-case time complexity of this algorithm is

\par\vspace{6pt}\noindent
\begin{tabular}{@{}p{0.98\linewidth}@{}}
$\square$\ \textbf{(A)} both $\mathcal{O}(N)$ and $\Omega(N)$ \\
\addlinespace[3pt]
$\square$\ \textbf{(B)} $\mathcal{O}(N)$ but not $\Omega(N)$ \\
\addlinespace[3pt]
$\square$\ \textbf{(C)} $\Omega(N)$ but not $\mathcal{O}(N)$ \\
\addlinespace[3pt]
$\square$\ \textbf{(D)} neither $\mathcal{O}(N)$ nor $\Omega(N)$
\end{tabular}

\tcblower

\textbf{Official Key:} \textbf{\large A} \hfill \textit{\small [Easy]}
\vspace{2pt}\hrule height 0.3pt\vspace{3pt}
\textbf{Topic:} Algorithms \\
\textbf{Subtopic:} \textit{Asymptotic worst case time and space complexity}
\vspace{2pt}\hrule height 0.3pt\vspace{3pt}
\textbf{Guide / Hint:} \\
{\small In the worst case (such as when the array is sorted or the disorder occurs at the very end), a single-pass algorithm that compares adjacent elements must traverse the entire array, performing $N-1$ comparisons. The worst-case runtime is asymptotically tight at $\Theta(N)$, meaning it is both $\mathcal{O}(N)$ and $\Omega(N)$.}

\end{tcolorbox}
\vspace{6pt}\hrule height 0.35pt\vspace{6pt}
\noindent
\begin{tcolorbox}[
    blanker,
    sidebyside,
    sidebyside align=top,
    lefthand width=0.55\textwidth,
    top=4pt, bottom=4pt, left=0pt, right=0pt,
    middle=16pt,
    underlay={
        \draw[dashed, black!80, line width=0.8pt, dash pattern=on 4pt off 3pt]
            ([xshift=0.55\textwidth+8pt]frame.north west) --
            ([xshift=0.55\textwidth+8pt]frame.south west);
    }
]
\textbf{Q.110} \quad {\footnotesize\textsc{[NAT $\cdot$ 1 Mark(s) $\cdot$ No Neg]}} \hfill {\footnotesize\textbf{GATE 2024 CS2} \quad \tiny\texttt{2024\_CS2\_CS\_Q35}}
\vspace{2pt}\hrule height 0.3pt\vspace{4pt}

Let $A$ be an array containing integer values. The distance of $A$ is defined as the minimum number of elements in $A$ that must be replaced with another integer so that the resulting array is sorted in non-decreasing order. The distance of the array $[2, 5, 3, 1, 4, 2, 6]$ is \_\_\_\_\_\_\. (Answer in integer)

\par\vspace{6pt}\noindent\textbf{Your Answer:} \fbox{\rule{0pt}{1.3ex}\hspace{3cm}}

\tcblower

\textbf{Official Key:} \textbf{\large 3 to 3} \hfill \textit{\small [Easy]}
\vspace{2pt}\hrule height 0.3pt\vspace{3pt}
\textbf{Topic:} Algorithms \\
\textbf{Subtopic:} \textit{Searching, sorting, hashing}
\vspace{2pt}\hrule height 0.3pt\vspace{3pt}
\textbf{Guide / Hint:} \\
{\small To minimize replacements to make array $A$ sorted in non-decreasing order, keep the longest non-decreasing subsequence (LNDS). For $A = [2, 5, 3, 1, 4, 2, 6]$, a longest non-decreasing subsequence is $[2, 3, 4, 6]$ of length 4. The minimum elements to replace is $n - \text{length(LNDS)} = 7 - 4 = 3$.}

\end{tcolorbox}
\vspace{6pt}\hrule height 0.35pt\vspace{6pt}
\noindent
\begin{tcolorbox}[
    blanker,
    sidebyside,
    sidebyside align=top,
    lefthand width=0.55\textwidth,
    top=4pt, bottom=4pt, left=0pt, right=0pt,
    middle=16pt,
    underlay={
        \draw[dashed, black!80, line width=0.8pt, dash pattern=on 4pt off 3pt]
            ([xshift=0.55\textwidth+8pt]frame.north west) --
            ([xshift=0.55\textwidth+8pt]frame.south west);
    }
]
\textbf{Q.111} \quad {\footnotesize\textsc{[MCQ $\cdot$ 1 Mark(s) $\cdot$ Neg: 0.33]}} \hfill {\footnotesize\textbf{GATE 2025 CS1} \quad \tiny\texttt{2025\_CS1\_CS\_Q18}}
\vspace{2pt}\hrule height 0.3pt\vspace{4pt}

Let $G$ be any undirected graph with positive edge weights, and $T$ be a minimum spanning tree of $G$. For any two vertices, $u$ and $v$, let $d_1(u, v)$ and $d_2(u, v)$ be the shortest distances between $u$ and $v$ in $G$ and $T$, respectively. Which ONE of the options is CORRECT for all possible $G, T, u$ and $v$?

\par\vspace{6pt}\noindent
\begin{tabular}{@{}p{0.48\linewidth}@{\hspace{0.04\linewidth}}p{0.48\linewidth}@{}}
$\square$\ \textbf{(A)} $d_1(u, v) = d_2(u, v)$ & $\square$\ \textbf{(B)} $d_1(u, v) \le d_2(u, v)$ \\
\addlinespace[3pt]
$\square$\ \textbf{(C)} $d_1(u, v) \ge d_2(u, v)$ & $\square$\ \textbf{(D)} $d_1(u, v) \ne d_2(u, v)$
\end{tabular}

\tcblower

\textbf{Official Key:} \textbf{\large B} \hfill \textit{\small [Easy]}
\vspace{2pt}\hrule height 0.3pt\vspace{3pt}
\textbf{Topic:} Algorithms \\
\textbf{Subtopic:} \textit{Graph traversals, minimum spanning trees, shortest paths}
\vspace{2pt}\hrule height 0.3pt\vspace{3pt}
\textbf{Guide / Hint:} \\
{\small Since the minimum spanning tree $T$ is a subgraph of $G$ ($E(T) \subseteq E(G)$), every path between $u$ and $v$ in $T$ is also a valid path in $G$. Thus, the shortest path distance in $G$ cannot exceed that in $T$, guaranteeing $d_1(u, v) \le d_2(u, v)$.}

\end{tcolorbox}
\vspace{6pt}\hrule height 0.35pt\vspace{6pt}
\noindent
\begin{tcolorbox}[
    blanker,
    sidebyside,
    sidebyside align=top,
    lefthand width=0.55\textwidth,
    top=4pt, bottom=4pt, left=0pt, right=0pt,
    middle=16pt,
    underlay={
        \draw[dashed, black!80, line width=0.8pt, dash pattern=on 4pt off 3pt]
            ([xshift=0.55\textwidth+8pt]frame.north west) --
            ([xshift=0.55\textwidth+8pt]frame.south west);
    }
]
\textbf{Q.112} \quad {\footnotesize\textsc{[NAT $\cdot$ 1 Mark(s) $\cdot$ No Neg]}} \hfill {\footnotesize\textbf{GATE 2025 CS1} \quad \tiny\texttt{2025\_CS1\_CS\_Q33}}
\vspace{2pt}\hrule height 0.3pt\vspace{4pt}

The pseudocode of a function fun() is given below:
\begin{verbatim}
fun(int A[0,...,n-1]){
    for i=0 to n-2
        for j=0 to n-i-2
            if (A[j] > A[j+1])
                then swap A[j] and A[j+1]
}
\end{verbatim}
Let $A[0, \ldots, 29]$ be an array storing 30 distinct integers in descending order. The number of swap operations that will be performed, if the function fun() is called with $A[0, \ldots, 29]$ as argument, is \_\_\_\_\_\_\. (Answer in integer)

\par\vspace{6pt}\noindent\textbf{Your Answer:} \fbox{\rule{0pt}{1.3ex}\hspace{3cm}}

\tcblower

\textbf{Official Key:} \textbf{\large 435 to 435} \hfill \textit{\small [Easy]}
\vspace{2pt}\hrule height 0.3pt\vspace{3pt}
\textbf{Topic:} Algorithms \\
\textbf{Subtopic:} \textit{Searching, sorting, hashing}
\vspace{2pt}\hrule height 0.3pt\vspace{3pt}
\textbf{Guide / Hint:} \\
{\small The function implements standard Bubble Sort. When sorting an array of $n = 30$ elements in strictly reverse order, every comparison triggers a swap operation. The total number of swaps is $\sum_{i=0}^{n-2} (n - 1 - i) = \frac{n(n-1)}{2} = \frac{30 \times 29}{2} = 435$.}

\end{tcolorbox}
\vspace{6pt}\hrule height 0.35pt\vspace{6pt}
\noindent
\begin{tcolorbox}[
    blanker,
    sidebyside,
    sidebyside align=top,
    lefthand width=0.55\textwidth,
    top=4pt, bottom=4pt, left=0pt, right=0pt,
    middle=16pt,
    underlay={
        \draw[dashed, black!80, line width=0.8pt, dash pattern=on 4pt off 3pt]
            ([xshift=0.55\textwidth+8pt]frame.north west) --
            ([xshift=0.55\textwidth+8pt]frame.south west);
    }
]
\textbf{Q.113} \quad {\footnotesize\textsc{[MCQ $\cdot$ 1 Mark(s) $\cdot$ Neg: 0.33]}} \hfill {\footnotesize\textbf{GATE 2025 CS2} \quad \tiny\texttt{2025\_CS2\_CS\_Q20}}
\vspace{2pt}\hrule height 0.3pt\vspace{4pt}

Consider an unordered list of $N$ distinct integers. What is the minimum number of element comparisons required to find an integer in the list that is NOT the largest in the list?

\par\vspace{6pt}\noindent
\begin{tabular}{@{}p{0.48\linewidth}@{\hspace{0.04\linewidth}}p{0.48\linewidth}@{}}
$\square$\ \textbf{(A)} 1 & $\square$\ \textbf{(B)} $N - 1$ \\
\addlinespace[3pt]
$\square$\ \textbf{(C)} $N$ & $\square$\ \textbf{(D)} $2N - 1$
\end{tabular}

\tcblower

\textbf{Official Key:} \textbf{\large A} \hfill \textit{\small [Easy]}
\vspace{2pt}\hrule height 0.3pt\vspace{3pt}
\textbf{Topic:} Algorithms \\
\textbf{Subtopic:} \textit{Searching and comparison complexity}
\vspace{2pt}\hrule height 0.3pt\vspace{3pt}
\textbf{Guide / Hint:} \\
{\small To find any element that is not the maximum, compare any two distinct elements in 1 comparison; the smaller of the two is guaranteed not to be the maximum.}

\end{tcolorbox}
\vspace{6pt}\hrule height 0.35pt\vspace{6pt}
\vspace{8pt}
\subsection*{Theory of Computation}
\vspace{-4pt}
\noindent
\begin{tcolorbox}[
    blanker,
    sidebyside,
    sidebyside align=top,
    lefthand width=0.55\textwidth,
    top=4pt, bottom=4pt, left=0pt, right=0pt,
    middle=16pt,
    underlay={
        \draw[dashed, black!80, line width=0.8pt, dash pattern=on 4pt off 3pt]
            ([xshift=0.55\textwidth+8pt]frame.north west) --
            ([xshift=0.55\textwidth+8pt]frame.south west);
    }
]
\textbf{Q.114} \quad {\footnotesize\textsc{[MCQ $\cdot$ 1 Mark(s) $\cdot$ Neg: 0.33]}} \hfill {\footnotesize\textbf{GATE 2018 CS} \quad \tiny\texttt{2018\_CS\_GA\_Q06\_v2}}
\vspace{2pt}\hrule height 0.3pt\vspace{4pt}

Let $N$ be an NFA with $n$ states. Let $k$ be the number of states of a minimal DFA which is equivalent to $N$. Which one of the following is necessarily true?

\par\vspace{6pt}\noindent
\begin{tabular}{@{}p{0.48\linewidth}@{\hspace{0.04\linewidth}}p{0.48\linewidth}@{}}
$\square$\ \textbf{(A)} $k \ge 2^n$ & $\square$\ \textbf{(B)} $k \ge n$ \\
\addlinespace[3pt]
$\square$\ \textbf{(C)} $k \le n^2$ & $\square$\ \textbf{(D)} $k \le 2^n$
\end{tabular}

\tcblower

\textbf{Official Key:} \textbf{\large D} \hfill \textit{\small [Easy]}
\vspace{2pt}\hrule height 0.3pt\vspace{3pt}
\textbf{Topic:} Theory of Computation \\
\textbf{Subtopic:} \textit{Regular expressions and finite automata}
\vspace{2pt}\hrule height 0.3pt\vspace{3pt}
\textbf{Guide / Hint:} \\
{\small By the subset construction algorithm, converting an $n$-state NFA to a DFA produces states that are subsets of the NFA states. Thus, the power set has $2^n$ possible states, ensuring that any equivalent minimal DFA has $k \le 2^n$ states.}

\end{tcolorbox}
\vspace{6pt}\hrule height 0.35pt\vspace{6pt}
\noindent
\begin{tcolorbox}[
    blanker,
    sidebyside,
    sidebyside align=top,
    lefthand width=0.55\textwidth,
    top=4pt, bottom=4pt, left=0pt, right=0pt,
    middle=16pt,
    underlay={
        \draw[dashed, black!80, line width=0.8pt, dash pattern=on 4pt off 3pt]
            ([xshift=0.55\textwidth+8pt]frame.north west) --
            ([xshift=0.55\textwidth+8pt]frame.south west);
    }
]
\textbf{Q.115} \quad {\footnotesize\textsc{[MCQ $\cdot$ 1 Mark(s) $\cdot$ Neg: 0.33]}} \hfill {\footnotesize\textbf{GATE 2018 CS} \quad \tiny\texttt{2018\_CS\_GA\_Q07\_v2}}
\vspace{2pt}\hrule height 0.3pt\vspace{4pt}

The set of all recursively enumerable languages is

\par\vspace{6pt}\noindent
\begin{tabular}{@{}p{0.98\linewidth}@{}}
$\square$\ \textbf{(A)} closed under complementation. \\
\addlinespace[3pt]
$\square$\ \textbf{(B)} closed under intersection. \\
\addlinespace[3pt]
$\square$\ \textbf{(C)} a subset of the set of all context-sensitive languages. \\
\addlinespace[3pt]
$\square$\ \textbf{(D)} an uncountable set.
\end{tabular}

\tcblower

\textbf{Official Key:} \textbf{\large B} \hfill \textit{\small [Easy]}
\vspace{2pt}\hrule height 0.3pt\vspace{3pt}
\textbf{Topic:} Theory of Computation \\
\textbf{Subtopic:} \textit{Regular and context-free languages, pumping lemma}
\vspace{2pt}\hrule height 0.3pt\vspace{3pt}
\textbf{Guide / Hint:} \\
{\small The set of recursively enumerable (RE) languages is closed under union, intersection, concatenation, and Kleene star, but not closed under complementation or set difference.}

\end{tcolorbox}
\vspace{6pt}\hrule height 0.35pt\vspace{6pt}
\noindent
\begin{tcolorbox}[
    blanker,
    sidebyside,
    sidebyside align=top,
    lefthand width=0.55\textwidth,
    top=4pt, bottom=4pt, left=0pt, right=0pt,
    middle=16pt,
    underlay={
        \draw[dashed, black!80, line width=0.8pt, dash pattern=on 4pt off 3pt]
            ([xshift=0.55\textwidth+8pt]frame.north west) --
            ([xshift=0.55\textwidth+8pt]frame.south west);
    }
]
\textbf{Q.116} \quad {\footnotesize\textsc{[MSQ $\cdot$ 1 Mark(s) $\cdot$ No Neg]}} \hfill {\footnotesize\textbf{GATE 2023 CS} \quad \tiny\texttt{2023\_CS\_CS\_Q24}}
\vspace{2pt}\hrule height 0.3pt\vspace{4pt}

Which of the following statements is/are CORRECT?

\par\vspace{6pt}\noindent
\begin{tabular}{@{}p{0.98\linewidth}@{}}
$\square$\ \textbf{(A)} The intersection of two regular languages is regular. \\
\addlinespace[3pt]
$\square$\ \textbf{(B)} The intersection of two context-free languages is context-free. \\
\addlinespace[3pt]
$\square$\ \textbf{(C)} The intersection of two recursive languages is recursive. \\
\addlinespace[3pt]
$\square$\ \textbf{(D)} The intersection of two recursively enumerable languages is recursively enumerable.
\end{tabular}

\tcblower

\textbf{Official Key:} \textbf{\large A, C, D} \hfill \textit{\small [Easy]}
\vspace{2pt}\hrule height 0.3pt\vspace{3pt}
\textbf{Topic:} Theory of Computation \\
\textbf{Subtopic:} \textit{Regular and context-free languages, pumping lemma}
\vspace{2pt}\hrule height 0.3pt\vspace{3pt}
\textbf{Guide / Hint:} \\
{\small Regular languages, recursive languages, and recursively enumerable languages are all closed under intersection (A, C, and D are CORRECT). Context-free languages are not closed under intersection; for example, $L_1 = \{a^n b^n c^m\}$ and $L_2 = \{a^m b^n c^n\}$ are both CFLs, but $L_1 \cap L_2 = \{a^n b^n c^n\}$ is not context-free.}

\end{tcolorbox}
\vspace{6pt}\hrule height 0.35pt\vspace{6pt}
\noindent
\begin{tcolorbox}[
    blanker,
    sidebyside,
    sidebyside align=top,
    lefthand width=0.55\textwidth,
    top=4pt, bottom=4pt, left=0pt, right=0pt,
    middle=16pt,
    underlay={
        \draw[dashed, black!80, line width=0.8pt, dash pattern=on 4pt off 3pt]
            ([xshift=0.55\textwidth+8pt]frame.north west) --
            ([xshift=0.55\textwidth+8pt]frame.south west);
    }
]
\textbf{Q.117} \quad {\footnotesize\textsc{[NAT $\cdot$ 2 Mark(s) $\cdot$ No Neg]}} \hfill {\footnotesize\textbf{GATE 2023 CS} \quad \tiny\texttt{2023\_CS\_CS\_Q63}}
\vspace{2pt}\hrule height 0.3pt\vspace{4pt}

Consider the language $L$ over the alphabet $\{0, 1\}$, given below:
\[ L = \{w \in \{0, 1\}^* \mid w \text{ does not contain three or more consecutive 1's}\}. \]
The minimum number of states in a Deterministic Finite-State Automaton (DFA) for $L$ is \_\_\_\_\_\_\_\_\. (Answer in integer)

\par\vspace{6pt}\noindent\textbf{Your Answer:} \fbox{\rule{0pt}{1.3ex}\hspace{3cm}}

\tcblower

\textbf{Official Key:} \textbf{\large 4 to 4} \hfill \textit{\small [Easy]}
\vspace{2pt}\hrule height 0.3pt\vspace{3pt}
\textbf{Topic:} Theory of Computation \\
\textbf{Subtopic:} \textit{Regular expressions and finite automata}
\vspace{2pt}\hrule height 0.3pt\vspace{3pt}
\textbf{Guide / Hint:} \\
{\small The language contains all binary strings not containing `111`. The minimal DFA tracks consecutive 1s: state $q_0$ (0 consecutive 1s, start), state $q_1$ (1 consecutive 1), state $q_2$ (2 consecutive 1s), and state $q_3$ (trap state reached on 3rd consecutive 1). States $q_0, q_1, q_2$ are accepting states. Thus, exactly 4 states are required.}

\end{tcolorbox}
\vspace{6pt}\hrule height 0.35pt\vspace{6pt}
\noindent
\begin{tcolorbox}[
    blanker,
    sidebyside,
    sidebyside align=top,
    lefthand width=0.55\textwidth,
    top=4pt, bottom=4pt, left=0pt, right=0pt,
    middle=16pt,
    underlay={
        \draw[dashed, black!80, line width=0.8pt, dash pattern=on 4pt off 3pt]
            ([xshift=0.55\textwidth+8pt]frame.north west) --
            ([xshift=0.55\textwidth+8pt]frame.south west);
    }
]
\textbf{Q.118} \quad {\footnotesize\textsc{[MCQ $\cdot$ 1 Mark(s) $\cdot$ Neg: 0.33]}} \hfill {\footnotesize\textbf{GATE 2025 CS1} \quad \tiny\texttt{2025\_CS1\_CS\_Q19}}
\vspace{2pt}\hrule height 0.3pt\vspace{4pt}

Consider the following context-free grammar $G$, where $S, A$, and $B$ are the variables (non-terminals), $a$ and $b$ are the terminal symbols, $S$ is the start variable, and the rules of $G$ are described as:
\[ S \to aaB \mid Abb \]
\[ A \to a \mid aA \]
\[ B \to b \mid bB \]
Which ONE of the languages $L(G)$ is accepted by $G$?

\par\vspace{6pt}\noindent
\begin{tabular}{@{}p{0.98\linewidth}@{}}
$\square$\ \textbf{(A)} $L(G) = \{a^2b^n \mid n \ge 1\} \cup \{a^nb^2 \mid n \ge 1\}$ \\
\addlinespace[3pt]
$\square$\ \textbf{(B)} $L(G) = \{a^nb^{2n} \mid n \ge 1\} \cup \{a^{2n}b^n \mid n \ge 1\}$ \\
\addlinespace[3pt]
$\square$\ \textbf{(C)} $L(G) = \{a^nb^n \mid n \ge 1\}$ \\
\addlinespace[3pt]
$\square$\ \textbf{(D)} $L(G) = \{a^{2n}b^{2n} \mid n \ge 1\}$
\end{tabular}

\tcblower

\textbf{Official Key:} \textbf{\large A} \hfill \textit{\small [Easy]}
\vspace{2pt}\hrule height 0.3pt\vspace{3pt}
\textbf{Topic:} Theory of Computation \\
\textbf{Subtopic:} \textit{Context-free grammars and push-down automata}
\vspace{2pt}\hrule height 0.3pt\vspace{3pt}
\textbf{Guide / Hint:} \\
{\small From the rules, $A$ generates $a^+$ (any $a^n$ with $n \ge 1$) and $B$ generates $b^+$ (any $b^n$ with $n \ge 1$). Thus, $aaB$ produces $a^2b^n$ ($n \ge 1$), and $Abb$ produces $a^nb^2$ ($n \ge 1$). Therefore, $L(G) = \{a^2b^n \mid n \ge 1\} \cup \{a^nb^2 \mid n \ge 1\}$.}

\end{tcolorbox}
\vspace{6pt}\hrule height 0.35pt\vspace{6pt}
\noindent
\begin{tcolorbox}[
    blanker,
    sidebyside,
    sidebyside align=top,
    lefthand width=0.55\textwidth,
    top=4pt, bottom=4pt, left=0pt, right=0pt,
    middle=16pt,
    underlay={
        \draw[dashed, black!80, line width=0.8pt, dash pattern=on 4pt off 3pt]
            ([xshift=0.55\textwidth+8pt]frame.north west) --
            ([xshift=0.55\textwidth+8pt]frame.south west);
    }
]
\textbf{Q.119} \quad {\footnotesize\textsc{[MSQ $\cdot$ 1 Mark(s) $\cdot$ No Neg]}} \hfill {\footnotesize\textbf{GATE 2025 CS1} \quad \tiny\texttt{2025\_CS1\_CS\_Q28}}
\vspace{2pt}\hrule height 0.3pt\vspace{4pt}

A regular language $L$ is accepted by a non-deterministic finite automaton (NFA) with $n$ states. Which of the following statement(s) is/are FALSE?

\par\vspace{6pt}\noindent
\begin{tabular}{@{}p{0.98\linewidth}@{}}
$\square$\ \textbf{(A)} $L$ may have an accepting NFA with $< n$ states. \\
\addlinespace[3pt]
$\square$\ \textbf{(B)} $L$ may have an accepting DFA with $< n$ states. \\
\addlinespace[3pt]
$\square$\ \textbf{(C)} There exists a DFA with $\le 2^n$ states that accepts $L$. \\
\addlinespace[3pt]
$\square$\ \textbf{(D)} Every DFA that accepts $L$ has $> 2^n$ states.
\end{tabular}

\tcblower

\textbf{Official Key:} \textbf{\large D} \hfill \textit{\small [Easy]}
\vspace{2pt}\hrule height 0.3pt\vspace{3pt}
\textbf{Topic:} Theory of Computation \\
\textbf{Subtopic:} \textit{Regular expressions and finite automata}
\vspace{2pt}\hrule height 0.3pt\vspace{3pt}
\textbf{Guide / Hint:} \\
{\small By the standard subset construction algorithm, every $n$-state NFA can be converted into an equivalent deterministic finite automaton (DFA) having at most $2^n$ states. Statement D asserts that every DFA must have strictly more than $2^n$ states, which is directly false.}

\end{tcolorbox}
\vspace{6pt}\hrule height 0.35pt\vspace{6pt}
\noindent
\begin{tcolorbox}[
    blanker,
    sidebyside,
    sidebyside align=top,
    lefthand width=0.55\textwidth,
    top=4pt, bottom=4pt, left=0pt, right=0pt,
    middle=16pt,
    underlay={
        \draw[dashed, black!80, line width=0.8pt, dash pattern=on 4pt off 3pt]
            ([xshift=0.55\textwidth+8pt]frame.north west) --
            ([xshift=0.55\textwidth+8pt]frame.south west);
    }
]
\textbf{Q.120} \quad {\footnotesize\textsc{[MCQ $\cdot$ 1 Mark(s) $\cdot$ Neg: 0.33]}} \hfill {\footnotesize\textbf{GATE 2025 CS2} \quad \tiny\texttt{2025\_CS2\_CS\_Q24}}
\vspace{2pt}\hrule height 0.3pt\vspace{4pt}

Which ONE of the following languages is accepted by a deterministic pushdown automaton?

\par\vspace{6pt}\noindent
\begin{tabular}{@{}p{0.98\linewidth}@{}}
$\square$\ \textbf{(A)} Any regular language. \\
\addlinespace[3pt]
$\square$\ \textbf{(B)} Any context-free language. \\
\addlinespace[3pt]
$\square$\ \textbf{(C)} Any language accepted by a non-deterministic pushdown automaton. \\
\addlinespace[3pt]
$\square$\ \textbf{(D)} Any decidable language.
\end{tabular}

\tcblower

\textbf{Official Key:} \textbf{\large A} \hfill \textit{\small [Easy]}
\vspace{2pt}\hrule height 0.3pt\vspace{3pt}
\textbf{Topic:} Theory of Computation \\
\textbf{Subtopic:} \textit{Context-free grammars and push-down automata}
\vspace{2pt}\hrule height 0.3pt\vspace{3pt}
\textbf{Guide / Hint:} \\
{\small Every regular language is recognized by a Deterministic Finite Automaton (DFA), which can be directly simulated by a DPDA without using stack operations.}

\end{tcolorbox}
\vspace{6pt}\hrule height 0.35pt\vspace{6pt}
\noindent
\begin{tcolorbox}[
    blanker,
    sidebyside,
    sidebyside align=top,
    lefthand width=0.55\textwidth,
    top=4pt, bottom=4pt, left=0pt, right=0pt,
    middle=16pt,
    underlay={
        \draw[dashed, black!80, line width=0.8pt, dash pattern=on 4pt off 3pt]
            ([xshift=0.55\textwidth+8pt]frame.north west) --
            ([xshift=0.55\textwidth+8pt]frame.south west);
    }
]
\textbf{Q.121} \quad {\footnotesize\textsc{[MCQ $\cdot$ 1 Mark(s) $\cdot$ Neg: 0.33]}} \hfill {\footnotesize\textbf{GATE 2025 CS2} \quad \tiny\texttt{2025\_CS2\_CS\_Q25}}
\vspace{2pt}\hrule height 0.3pt\vspace{4pt}

Let $G_1, G_2$ be Context Free Grammars (CFGs) and $R$ be a regular expression. For a grammar $G$, let $L(G)$ denote the language generated by $G$. Which ONE among the following questions is decidable?

\par\vspace{6pt}\noindent
\begin{tabular}{@{}p{0.98\linewidth}@{}}
$\square$\ \textbf{(A)} Is $L(G_1) = L(G_2)$? \\
\addlinespace[3pt]
$\square$\ \textbf{(B)} Is $L(G_1) \cap L(G_2) = \emptyset$? \\
\addlinespace[3pt]
$\square$\ \textbf{(C)} Is $L(G_1) = L(R)$? \\
\addlinespace[3pt]
$\square$\ \textbf{(D)} Is $L(G_1) = \emptyset$?
\end{tabular}

\tcblower

\textbf{Official Key:} \textbf{\large D} \hfill \textit{\small [Easy]}
\vspace{2pt}\hrule height 0.3pt\vspace{3pt}
\textbf{Topic:} Theory of Computation \\
\textbf{Subtopic:} \textit{Turing machines and undecidability}
\vspace{2pt}\hrule height 0.3pt\vspace{3pt}
\textbf{Guide / Hint:} \\
{\small For Context-Free Grammars, emptiness ($L(G) = \emptyset$) is decidable via finding generating non-terminals, whereas equivalence ($L(G_1) = L(G_2)$) and disjointness ($L(G_1) \cap L(G_2) = \emptyset$) are undecidable.}

\end{tcolorbox}
\vspace{6pt}\hrule height 0.35pt\vspace{6pt}
\vspace{8pt}
\subsection*{Compiler Design}
\vspace{-4pt}
\noindent
\begin{tcolorbox}[
    blanker,
    sidebyside,
    sidebyside align=top,
    lefthand width=0.55\textwidth,
    top=4pt, bottom=4pt, left=0pt, right=0pt,
    middle=16pt,
    underlay={
        \draw[dashed, black!80, line width=0.8pt, dash pattern=on 4pt off 3pt]
            ([xshift=0.55\textwidth+8pt]frame.north west) --
            ([xshift=0.55\textwidth+8pt]frame.south west);
    }
]
\textbf{Q.122} \quad {\footnotesize\textsc{[MCQ $\cdot$ 1 Mark(s) $\cdot$ Neg: 0.33]}} \hfill {\footnotesize\textbf{GATE 2018 CS} \quad \tiny\texttt{2018\_CS\_GA\_Q08\_v2}}
\vspace{2pt}\hrule height 0.3pt\vspace{4pt}

Which one of the following statements is FALSE?

\par\vspace{6pt}\noindent
\begin{tabular}{@{}p{0.98\linewidth}@{}}
$\square$\ \textbf{(A)} Context-free grammar can be used to specify both lexical and syntax rules. \\
\addlinespace[3pt]
$\square$\ \textbf{(B)} Type checking is done before parsing. \\
\addlinespace[3pt]
$\square$\ \textbf{(C)} High-level language programs can be converted to intermediate code directly. \\
\addlinespace[3pt]
$\square$\ \textbf{(D)} Arguments to a function can be passed using the program stack.
\end{tabular}

\tcblower

\textbf{Official Key:} \textbf{\large B} \hfill \textit{\small [Easy]}
\vspace{2pt}\hrule height 0.3pt\vspace{3pt}
\textbf{Topic:} Compiler Design \\
\textbf{Subtopic:} \textit{Lexical analysis, parsing, syntax-directed translation}
\vspace{2pt}\hrule height 0.3pt\vspace{3pt}
\textbf{Guide / Hint:} \\
{\small Type checking verifies semantic consistency according to language type rules and is performed during semantic analysis, which occurs after syntax analysis (parsing). Therefore, the statement that type checking is done during syntax analysis is FALSE.}

\end{tcolorbox}
\vspace{6pt}\hrule height 0.35pt\vspace{6pt}
\noindent
\begin{tcolorbox}[
    blanker,
    sidebyside,
    sidebyside align=top,
    lefthand width=0.55\textwidth,
    top=4pt, bottom=4pt, left=0pt, right=0pt,
    middle=16pt,
    underlay={
        \draw[dashed, black!80, line width=0.8pt, dash pattern=on 4pt off 3pt]
            ([xshift=0.55\textwidth+8pt]frame.north west) --
            ([xshift=0.55\textwidth+8pt]frame.south west);
    }
]
\textbf{Q.123} \quad {\footnotesize\textsc{[MCQ $\cdot$ 1 Mark(s) $\cdot$ Neg: 0.33]}} \hfill {\footnotesize\textbf{GATE 2022 CS} \quad \tiny\texttt{2022\_CS\_CS\_Q13}}
\vspace{2pt}\hrule height 0.3pt\vspace{4pt}

Which one of the following statements is TRUE?

\par\vspace{6pt}\noindent
\begin{tabular}{@{}p{0.98\linewidth}@{}}
$\square$\ \textbf{(A)} The LALR(1) parser for a grammar G cannot have reduce-reduce conflict if the LR(1) parser for G does not have reduce-reduce conflict. \\
\addlinespace[3pt]
$\square$\ \textbf{(B)} Symbol table is accessed only during the lexical analysis phase. \\
\addlinespace[3pt]
$\square$\ \textbf{(C)} Data flow analysis is necessary for run-time memory management. \\
\addlinespace[3pt]
$\square$\ \textbf{(D)} LR(1) parsing is sufficient for deterministic context-free languages.
\end{tabular}

\tcblower

\textbf{Official Key:} \textbf{\large D} \hfill \textit{\small [Easy]}
\vspace{2pt}\hrule height 0.3pt\vspace{3pt}
\textbf{Topic:} Compiler Design \\
\textbf{Subtopic:} \textit{Lexical analysis, parsing, syntax-directed translation}
\vspace{2pt}\hrule height 0.3pt\vspace{3pt}
\textbf{Guide / Hint:} \\
{\small LR(1) grammars describe exactly the family of Deterministic Context-Free Languages (DCFLs), meaning LR(1) parsing is both necessary and sufficient for any DCFL. LALR(1) can introduce reduce-reduce conflicts not present in LR(1), and symbol tables are accessed throughout compilation.}

\end{tcolorbox}
\vspace{6pt}\hrule height 0.35pt\vspace{6pt}
\noindent
\begin{tcolorbox}[
    blanker,
    sidebyside,
    sidebyside align=top,
    lefthand width=0.55\textwidth,
    top=4pt, bottom=4pt, left=0pt, right=0pt,
    middle=16pt,
    underlay={
        \draw[dashed, black!80, line width=0.8pt, dash pattern=on 4pt off 3pt]
            ([xshift=0.55\textwidth+8pt]frame.north west) --
            ([xshift=0.55\textwidth+8pt]frame.south west);
    }
]
\textbf{Q.124} \quad {\footnotesize\textsc{[MSQ $\cdot$ 1 Mark(s) $\cdot$ No Neg]}} \hfill {\footnotesize\textbf{GATE 2024 CS1} \quad \tiny\texttt{2024\_CS1\_CS\_Q26}}
\vspace{2pt}\hrule height 0.3pt\vspace{4pt}

Which of the following is/are Bottom-Up Parser(s)?

\par\vspace{6pt}\noindent
\begin{tabular}{@{}p{0.48\linewidth}@{\hspace{0.04\linewidth}}p{0.48\linewidth}@{}}
$\square$\ \textbf{(A)} Shift-reduce Parser & $\square$\ \textbf{(B)} Predictive Parser \\
\addlinespace[3pt]
$\square$\ \textbf{(C)} LL(1) Parser & $\square$\ \textbf{(D)} LR Parser
\end{tabular}

\tcblower

\textbf{Official Key:} \textbf{\large A;D} \hfill \textit{\small [Easy]}
\vspace{2pt}\hrule height 0.3pt\vspace{3pt}
\textbf{Topic:} Compiler Design \\
\textbf{Subtopic:} \textit{Lexical analysis, parsing, syntax-directed translation}
\vspace{2pt}\hrule height 0.3pt\vspace{3pt}
\textbf{Guide / Hint:} \\
{\small Bottom-up parsers reduce an input string back to the start symbol by finding handles. Shift-reduce parsers and LR parsers (SLR, CLR, LALR) are bottom-up parsers. In contrast, Predictive and LL(1) parsers construct the derivation tree from the root downward and are top-down parsers.}

\end{tcolorbox}
\vspace{6pt}\hrule height 0.35pt\vspace{6pt}
\noindent
\begin{tcolorbox}[
    blanker,
    sidebyside,
    sidebyside align=top,
    lefthand width=0.55\textwidth,
    top=4pt, bottom=4pt, left=0pt, right=0pt,
    middle=16pt,
    underlay={
        \draw[dashed, black!80, line width=0.8pt, dash pattern=on 4pt off 3pt]
            ([xshift=0.55\textwidth+8pt]frame.north west) --
            ([xshift=0.55\textwidth+8pt]frame.south west);
    }
]
\textbf{Q.125} \quad {\footnotesize\textsc{[MCQ $\cdot$ 1 Mark(s) $\cdot$ Neg: 0.33]}} \hfill {\footnotesize\textbf{GATE 2024 CS2} \quad \tiny\texttt{2024\_CS2\_CS\_Q21}}
\vspace{2pt}\hrule height 0.3pt\vspace{4pt}

Consider the following two sets:
\par\medskip
\begin{tabular}{ll}
\textbf{Set X} & \textbf{Set Y} \\
P. Lexical Analyzer & 1. Abstract Syntax Tree \\
Q. Syntax Analyzer & 2. Token \\
R. Intermediate Code Generator & 3. Parse Tree \\
S. Code Optimizer & 4. Constant Folding
\end{tabular}
\par\medskip
Which one of the following options is the CORRECT match from Set X to Set Y?

\par\vspace{6pt}\noindent
\begin{tabular}{@{}p{0.48\linewidth}@{\hspace{0.04\linewidth}}p{0.48\linewidth}@{}}
$\square$\ \textbf{(A)} P -- 4; Q -- 1; R -- 3; S -- 2 & $\square$\ \textbf{(B)} P -- 2; Q -- 3; R -- 1; S -- 4 \\
\addlinespace[3pt]
$\square$\ \textbf{(C)} P -- 2; Q -- 1; R -- 3; S -- 4 & $\square$\ \textbf{(D)} P -- 4; Q -- 3; R -- 2; S -- 1
\end{tabular}

\tcblower

\textbf{Official Key:} \textbf{\large B} \hfill \textit{\small [Easy]}
\vspace{2pt}\hrule height 0.3pt\vspace{3pt}
\textbf{Topic:} Compiler Design \\
\textbf{Subtopic:} \textit{Lexical analysis, parsing, syntax-directed translation}
\vspace{2pt}\hrule height 0.3pt\vspace{3pt}
\textbf{Guide / Hint:} \\
{\small Lexical Analyzer outputs tokens (P - 2). Syntax Analyzer constructs parse trees (Q - 3). Intermediate Code Generator generates abstract syntax trees or intermediate representations (R - 1). Code Optimizer performs constant folding (S - 4). Thus, the correct matching is P - 2; Q - 3; R - 1; S - 4.}

\end{tcolorbox}
\vspace{6pt}\hrule height 0.35pt\vspace{6pt}
\noindent
\begin{tcolorbox}[
    blanker,
    sidebyside,
    sidebyside align=top,
    lefthand width=0.55\textwidth,
    top=4pt, bottom=4pt, left=0pt, right=0pt,
    middle=16pt,
    underlay={
        \draw[dashed, black!80, line width=0.8pt, dash pattern=on 4pt off 3pt]
            ([xshift=0.55\textwidth+8pt]frame.north west) --
            ([xshift=0.55\textwidth+8pt]frame.south west);
    }
]
\textbf{Q.126} \quad {\footnotesize\textsc{[MCQ $\cdot$ 1 Mark(s) $\cdot$ Neg: 0.33]}} \hfill {\footnotesize\textbf{GATE 2025 CS1} \quad \tiny\texttt{2025\_CS1\_CS\_Q12}}
\vspace{2pt}\hrule height 0.3pt\vspace{4pt}

Which ONE of the following statements is FALSE regarding the symbol table?

\par\vspace{6pt}\noindent
\begin{tabular}{@{}p{0.98\linewidth}@{}}
$\square$\ \textbf{(A)} Symbol table is responsible for keeping track of the scope of variables. \\
\addlinespace[3pt]
$\square$\ \textbf{(B)} Symbol table can be implemented using a binary search tree. \\
\addlinespace[3pt]
$\square$\ \textbf{(C)} Symbol table is not required after the parsing phase. \\
\addlinespace[3pt]
$\square$\ \textbf{(D)} Symbol table is created during the lexical analysis phase.
\end{tabular}

\tcblower

\textbf{Official Key:} \textbf{\large C} \hfill \textit{\small [Easy]}
\vspace{2pt}\hrule height 0.3pt\vspace{3pt}
\textbf{Topic:} Compiler Design \\
\textbf{Subtopic:} \textit{Lexical analysis, parsing, syntax-directed translation}
\vspace{2pt}\hrule height 0.3pt\vspace{3pt}
\textbf{Guide / Hint:} \\
{\small The symbol table is actively used throughout all phases of the compiler, including semantic analysis, intermediate code generation, code optimization, and target code generation for type checking, memory allocation, and address resolution. Thus, claiming it is not required after parsing is false.}

\end{tcolorbox}
\vspace{6pt}\hrule height 0.35pt\vspace{6pt}
\noindent
\begin{tcolorbox}[
    blanker,
    sidebyside,
    sidebyside align=top,
    lefthand width=0.55\textwidth,
    top=4pt, bottom=4pt, left=0pt, right=0pt,
    middle=16pt,
    underlay={
        \draw[dashed, black!80, line width=0.8pt, dash pattern=on 4pt off 3pt]
            ([xshift=0.55\textwidth+8pt]frame.north west) --
            ([xshift=0.55\textwidth+8pt]frame.south west);
    }
]
\textbf{Q.127} \quad {\footnotesize\textsc{[MCQ $\cdot$ 1 Mark(s) $\cdot$ Neg: 0.33]}} \hfill {\footnotesize\textbf{GATE 2025 CS1} \quad \tiny\texttt{2025\_CS1\_CS\_Q13}}
\vspace{2pt}\hrule height 0.3pt\vspace{4pt}

Which ONE of the following techniques used in compiler code optimization uses live variable analysis?

\par\vspace{6pt}\noindent
\begin{tabular}{@{}p{0.98\linewidth}@{}}
$\square$\ \textbf{(A)} Run-time function call management \\
\addlinespace[3pt]
$\square$\ \textbf{(B)} Register assignment to variables \\
\addlinespace[3pt]
$\square$\ \textbf{(C)} Strength reduction \\
\addlinespace[3pt]
$\square$\ \textbf{(D)} Constant folding
\end{tabular}

\tcblower

\textbf{Official Key:} \textbf{\large B} \hfill \textit{\small [Easy]}
\vspace{2pt}\hrule height 0.3pt\vspace{3pt}
\textbf{Topic:} Compiler Design \\
\textbf{Subtopic:} \textit{Data flow analyses: constant propagation, liveness analysis, common sub expression elimination}
\vspace{2pt}\hrule height 0.3pt\vspace{3pt}
\textbf{Guide / Hint:} \\
{\small Live variable analysis determines whether the value in a variable could be read along some execution path before being redefined. This information is directly used in register allocation and assignment to construct interference graphs and allocate registers efficiently.}

\end{tcolorbox}
\vspace{6pt}\hrule height 0.35pt\vspace{6pt}
\noindent
\begin{tcolorbox}[
    blanker,
    sidebyside,
    sidebyside align=top,
    lefthand width=0.55\textwidth,
    top=4pt, bottom=4pt, left=0pt, right=0pt,
    middle=16pt,
    underlay={
        \draw[dashed, black!80, line width=0.8pt, dash pattern=on 4pt off 3pt]
            ([xshift=0.55\textwidth+8pt]frame.north west) --
            ([xshift=0.55\textwidth+8pt]frame.south west);
    }
]
\textbf{Q.128} \quad {\footnotesize\textsc{[MSQ $\cdot$ 2 Mark(s) $\cdot$ No Neg]}} \hfill {\footnotesize\textbf{GATE 2025 CS1} \quad \tiny\texttt{2025\_CS1\_CS\_Q46}}
\vspace{2pt}\hrule height 0.3pt\vspace{4pt}

Which of the following statement(s) is/are TRUE while computing First and Follow during top down parsing by a compiler?

\par\vspace{6pt}\noindent
\begin{tabular}{@{}p{0.98\linewidth}@{}}
$\square$\ \textbf{(A)} For a production $A \to \epsilon$, $\epsilon$ will be added to $First(A)$. \\
\addlinespace[3pt]
$\square$\ \textbf{(B)} If there is any input right end marker, it will be added to $First(S)$, where $S$ is the start symbol. \\
\addlinespace[3pt]
$\square$\ \textbf{(C)} For a production $A \to \epsilon$, $\epsilon$ will be added to $Follow(A)$. \\
\addlinespace[3pt]
$\square$\ \textbf{(D)} If there is any input right end marker, it will be added to $Follow(S)$, where $S$ is the start symbol.
\end{tabular}

\tcblower

\textbf{Official Key:} \textbf{\large A;D} \hfill \textit{\small [Easy]}
\vspace{2pt}\hrule height 0.3pt\vspace{3pt}
\textbf{Topic:} Compiler Design \\
\textbf{Subtopic:} \textit{Lexical analysis, parsing, syntax-directed translation}
\vspace{2pt}\hrule height 0.3pt\vspace{3pt}
\textbf{Guide / Hint:} \\
{\small By definition of FIRST and FOLLOW: if $A \to \epsilon$, then $\epsilon \in \text{FIRST}(A)$ (A is true). $\epsilon$ can never belong to any FOLLOW set (C is false). The right end marker (\$) is placed into $\text{FOLLOW}(S)$ for the start symbol $S$ to mark the end of input (D is true, B is false).}

\end{tcolorbox}
\vspace{6pt}\hrule height 0.35pt\vspace{6pt}
\noindent
\begin{tcolorbox}[
    blanker,
    sidebyside,
    sidebyside align=top,
    lefthand width=0.55\textwidth,
    top=4pt, bottom=4pt, left=0pt, right=0pt,
    middle=16pt,
    underlay={
        \draw[dashed, black!80, line width=0.8pt, dash pattern=on 4pt off 3pt]
            ([xshift=0.55\textwidth+8pt]frame.north west) --
            ([xshift=0.55\textwidth+8pt]frame.south west);
    }
]
\textbf{Q.129} \quad {\footnotesize\textsc{[MCQ $\cdot$ 1 Mark(s) $\cdot$ Neg: 0.33]}} \hfill {\footnotesize\textbf{GATE 2025 CS2} \quad \tiny\texttt{2025\_CS2\_CS\_Q21}}
\vspace{2pt}\hrule height 0.3pt\vspace{4pt}

Consider the following statements about the use of backpatching in a compiler for intermediate code generation:
(I) Backpatching can be used to generate code for Boolean expression in one pass.
(II) Backpatching can be used to generate code for flow-of-control statements in one pass.
Which ONE of the following options is CORRECT?

\par\vspace{6pt}\noindent
\begin{tabular}{@{}p{0.98\linewidth}@{}}
$\square$\ \textbf{(A)} Only (I) is correct. \\
\addlinespace[3pt]
$\square$\ \textbf{(B)} Only (II) is correct. \\
\addlinespace[3pt]
$\square$\ \textbf{(C)} Both (I) and (II) are correct. \\
\addlinespace[3pt]
$\square$\ \textbf{(D)} Neither (I) nor (II) is correct.
\end{tabular}

\tcblower

\textbf{Official Key:} \textbf{\large C} \hfill \textit{\small [Easy]}
\vspace{2pt}\hrule height 0.3pt\vspace{3pt}
\textbf{Topic:} Compiler Design \\
\textbf{Subtopic:} \textit{Intermediate code generation and backpatching}
\vspace{2pt}\hrule height 0.3pt\vspace{3pt}
\textbf{Guide / Hint:} \\
{\small Backpatching is a standard technique used in single-pass compilers to synthesize target labels for forward jumps in both Boolean expressions and flow-of-control statements.}

\end{tcolorbox}
\vspace{6pt}\hrule height 0.35pt\vspace{6pt}
\vspace{8pt}
\subsection*{Operating Systems}
\vspace{-4pt}
\noindent
\begin{tcolorbox}[
    blanker,
    sidebyside,
    sidebyside align=top,
    lefthand width=0.55\textwidth,
    top=4pt, bottom=4pt, left=0pt, right=0pt,
    middle=16pt,
    underlay={
        \draw[dashed, black!80, line width=0.8pt, dash pattern=on 4pt off 3pt]
            ([xshift=0.55\textwidth+8pt]frame.north west) --
            ([xshift=0.55\textwidth+8pt]frame.south west);
    }
]
\textbf{Q.130} \quad {\footnotesize\textsc{[NAT $\cdot$ 1 Mark(s) $\cdot$ No Neg]}} \hfill {\footnotesize\textbf{GATE 2018 CS} \quad \tiny\texttt{2018\_CS\_GA\_Q24}}
\vspace{2pt}\hrule height 0.3pt\vspace{4pt}

Consider a system with 3 processes that share 4 instances of the same resource type. Each process can request at most $K$ instances of the resource type at any time. The largest value of $K$ that will always avoid deadlock is \_\_\_\_\_\_\_\_.

\par\vspace{6pt}\noindent\textbf{Your Answer:} \fbox{\rule{0pt}{1.3ex}\hspace{3cm}}

\tcblower

\textbf{Official Key:} \textbf{\large 2 to 2} \hfill \textit{\small [Easy]}
\vspace{2pt}\hrule height 0.3pt\vspace{3pt}
\textbf{Topic:} Operating System \\
\textbf{Subtopic:} \textit{Deadlock}
\vspace{2pt}\hrule height 0.3pt\vspace{3pt}
\textbf{Guide / Hint:} \\
{\small For $n$ processes sharing $m$ instances of a single resource type where each process needs at most $K$ instances, deadlock is guaranteed to be avoided if $\sum_{i=1}^n (K - 1) < m \implies n(K - 1) \le m - 1$. Here $n = 3$ and $m = 4$, so $3(K - 1) \le 3 \implies K - 1 \le 1 \implies K \le 2$.}

\end{tcolorbox}
\vspace{6pt}\hrule height 0.35pt\vspace{6pt}
\noindent
\begin{tcolorbox}[
    blanker,
    sidebyside,
    sidebyside align=top,
    lefthand width=0.55\textwidth,
    top=4pt, bottom=4pt, left=0pt, right=0pt,
    middle=16pt,
    underlay={
        \draw[dashed, black!80, line width=0.8pt, dash pattern=on 4pt off 3pt]
            ([xshift=0.55\textwidth+8pt]frame.north west) --
            ([xshift=0.55\textwidth+8pt]frame.south west);
    }
]
\textbf{Q.131} \quad {\footnotesize\textsc{[MSQ $\cdot$ 1 Mark(s) $\cdot$ No Neg]}} \hfill {\footnotesize\textbf{GATE 2022 CS} \quad \tiny\texttt{2022\_CS\_CS\_Q26}}
\vspace{2pt}\hrule height 0.3pt\vspace{4pt}

Which of the following statements is/are TRUE with respect to deadlocks?

\par\vspace{6pt}\noindent
\begin{tabular}{@{}p{0.98\linewidth}@{}}
$\square$\ \textbf{(A)} Circular wait is a necessary condition for the formation of deadlock. \\
\addlinespace[3pt]
$\square$\ \textbf{(B)} In a system where each resource has more than one instance, a cycle in its wait-for graph indicates the presence of a deadlock. \\
\addlinespace[3pt]
$\square$\ \textbf{(C)} If the current allocation of resources to processes leads the system to unsafe state, then deadlock will necessarily occur. \\
\addlinespace[3pt]
$\square$\ \textbf{(D)} In the resource-allocation graph of a system, if every edge is an assignment edge, then the system is not in deadlock state.
\end{tabular}

\tcblower

\textbf{Official Key:} \textbf{\large A,D} \hfill \textit{\small [Easy]}
\vspace{2pt}\hrule height 0.3pt\vspace{3pt}
\textbf{Topic:} Operating System \\
\textbf{Subtopic:} \textit{Deadlock}
\vspace{2pt}\hrule height 0.3pt\vspace{3pt}
\textbf{Guide / Hint:} \\
{\small Circular wait is one of Coffman's four necessary conditions for deadlock (A is TRUE). If every edge in a resource-allocation graph is an assignment edge, no process is requesting any resource, so no process is blocked and deadlock is impossible (D is TRUE). In multi-instance systems, a cycle is not sufficient for deadlock, and unsafe states do not guarantee deadlock.}

\end{tcolorbox}
\vspace{6pt}\hrule height 0.35pt\vspace{6pt}
\noindent
\begin{tcolorbox}[
    blanker,
    sidebyside,
    sidebyside align=top,
    lefthand width=0.55\textwidth,
    top=4pt, bottom=4pt, left=0pt, right=0pt,
    middle=16pt,
    underlay={
        \draw[dashed, black!80, line width=0.8pt, dash pattern=on 4pt off 3pt]
            ([xshift=0.55\textwidth+8pt]frame.north west) --
            ([xshift=0.55\textwidth+8pt]frame.south west);
    }
]
\textbf{Q.132} \quad {\footnotesize\textsc{[MSQ $\cdot$ 1 Mark(s) $\cdot$ No Neg]}} \hfill {\footnotesize\textbf{GATE 2023 CS} \quad \tiny\texttt{2023\_CS\_CS\_Q22}}
\vspace{2pt}\hrule height 0.3pt\vspace{4pt}

Which one or more of the following need to be saved on a context switch from one thread (T1) of a process to another thread (T2) of the same process?

\par\vspace{6pt}\noindent
\begin{tabular}{@{}p{0.48\linewidth}@{\hspace{0.04\linewidth}}p{0.48\linewidth}@{}}
$\square$\ \textbf{(A)} Page table base register & $\square$\ \textbf{(B)} Stack pointer \\
\addlinespace[3pt]
$\square$\ \textbf{(C)} Program counter & $\square$\ \textbf{(D)} General purpose registers
\end{tabular}

\tcblower

\textbf{Official Key:} \textbf{\large B, C, D} \hfill \textit{\small [Easy]}
\vspace{2pt}\hrule height 0.3pt\vspace{3pt}
\textbf{Topic:} Operating System \\
\textbf{Subtopic:} \textit{System calls, processes, threads, inter-process communication}
\vspace{2pt}\hrule height 0.3pt\vspace{3pt}
\textbf{Guide / Hint:} \\
{\small Threads of the same process share the process address space, open files, and the page table base register (PTBR is not switched). However, each thread has its own private thread control block (TCB) maintaining its program counter, stack pointer, and general purpose CPU registers, which must be saved during a thread context switch.}

\end{tcolorbox}
\vspace{6pt}\hrule height 0.35pt\vspace{6pt}
\noindent
\begin{tcolorbox}[
    blanker,
    sidebyside,
    sidebyside align=top,
    lefthand width=0.55\textwidth,
    top=4pt, bottom=4pt, left=0pt, right=0pt,
    middle=16pt,
    underlay={
        \draw[dashed, black!80, line width=0.8pt, dash pattern=on 4pt off 3pt]
            ([xshift=0.55\textwidth+8pt]frame.north west) --
            ([xshift=0.55\textwidth+8pt]frame.south west);
    }
]
\textbf{Q.133} \quad {\footnotesize\textsc{[MSQ $\cdot$ 1 Mark(s) $\cdot$ No Neg]}} \hfill {\footnotesize\textbf{GATE 2023 CS} \quad \tiny\texttt{2023\_CS\_CS\_Q23}}
\vspace{2pt}\hrule height 0.3pt\vspace{4pt}

Which one or more of the following options guarantee that a computer system will transition from user mode to kernel mode?

\par\vspace{6pt}\noindent
\begin{tabular}{@{}p{0.48\linewidth}@{\hspace{0.04\linewidth}}p{0.48\linewidth}@{}}
$\square$\ \textbf{(A)} Function Call & $\square$\ \textbf{(B)} `malloc` Call \\
\addlinespace[3pt]
$\square$\ \textbf{(C)} Page Fault & $\square$\ \textbf{(D)} System Call
\end{tabular}

\tcblower

\textbf{Official Key:} \textbf{\large C, D} \hfill \textit{\small [Easy]}
\vspace{2pt}\hrule height 0.3pt\vspace{3pt}
\textbf{Topic:} Operating System \\
\textbf{Subtopic:} \textit{System calls, processes, threads, inter-process communication}
\vspace{2pt}\hrule height 0.3pt\vspace{3pt}
\textbf{Guide / Hint:} \\
{\small A transition to kernel mode is triggered by traps, exceptions, or hardware interrupts. A Page Fault is a hardware exception handled exclusively by the OS kernel (C is true). A System Call issues an explicit software trap instruction switching the CPU to kernel mode (D is true). Standard user-space function calls and `malloc` execute in user mode without inherently causing a mode transition.}

\end{tcolorbox}
\vspace{6pt}\hrule height 0.35pt\vspace{6pt}
\noindent
\begin{tcolorbox}[
    blanker,
    sidebyside,
    sidebyside align=top,
    lefthand width=0.55\textwidth,
    top=4pt, bottom=4pt, left=0pt, right=0pt,
    middle=16pt,
    underlay={
        \draw[dashed, black!80, line width=0.8pt, dash pattern=on 4pt off 3pt]
            ([xshift=0.55\textwidth+8pt]frame.north west) --
            ([xshift=0.55\textwidth+8pt]frame.south west);
    }
]
\textbf{Q.134} \quad {\footnotesize\textsc{[MSQ $\cdot$ 1 Mark(s) $\cdot$ No Neg]}} \hfill {\footnotesize\textbf{GATE 2023 CS} \quad \tiny\texttt{2023\_CS\_CS\_Q27}}
\vspace{2pt}\hrule height 0.3pt\vspace{4pt}

Which one or more of the following CPU scheduling algorithms can potentially cause starvation?

\par\vspace{6pt}\noindent
\begin{tabular}{@{}p{0.48\linewidth}@{\hspace{0.04\linewidth}}p{0.48\linewidth}@{}}
$\square$\ \textbf{(A)} First-in First-Out & $\square$\ \textbf{(B)} Round Robin \\
\addlinespace[3pt]
$\square$\ \textbf{(C)} Priority Scheduling & $\square$\ \textbf{(D)} Shortest Job First
\end{tabular}

\tcblower

\textbf{Official Key:} \textbf{\large A, C, D or C, D} \hfill \textit{\small [Easy]}
\vspace{2pt}\hrule height 0.3pt\vspace{3pt}
\textbf{Topic:} Operating System \\
\textbf{Subtopic:} \textit{CPU and I/O scheduling}
\vspace{2pt}\hrule height 0.3pt\vspace{3pt}
\textbf{Guide / Hint:} \\
{\small Priority Scheduling can cause starvation for low-priority processes if higher-priority processes arrive continuously (C). Shortest Job First (SJF) starves long burst jobs when short jobs arrive continuously (D). Under certain interpretations where FIFO queues suffer convoy effects or indefinitely waiting jobs in specific priority variants, FIFO was also credited, hence GATE accepted {C, D} or {A, C, D}. Round Robin provides bounded waiting and prevents starvation.}

\end{tcolorbox}
\vspace{6pt}\hrule height 0.35pt\vspace{6pt}
\noindent
\begin{tcolorbox}[
    blanker,
    sidebyside,
    sidebyside align=top,
    lefthand width=0.55\textwidth,
    top=4pt, bottom=4pt, left=0pt, right=0pt,
    middle=16pt,
    underlay={
        \draw[dashed, black!80, line width=0.8pt, dash pattern=on 4pt off 3pt]
            ([xshift=0.55\textwidth+8pt]frame.north west) --
            ([xshift=0.55\textwidth+8pt]frame.south west);
    }
]
\textbf{Q.135} \quad {\footnotesize\textsc{[MSQ $\cdot$ 1 Mark(s) $\cdot$ No Neg]}} \hfill {\footnotesize\textbf{GATE 2024 CS1} \quad \tiny\texttt{2024\_CS1\_CS\_Q24}}
\vspace{2pt}\hrule height 0.3pt\vspace{4pt}

Which of the following statements about threads is/are TRUE?

\par\vspace{6pt}\noindent
\begin{tabular}{@{}p{0.98\linewidth}@{}}
$\square$\ \textbf{(A)} Threads can only be implemented in kernel space \\
\addlinespace[3pt]
$\square$\ \textbf{(B)} Each thread has its own file descriptor table for open files \\
\addlinespace[3pt]
$\square$\ \textbf{(C)} All the threads belonging to a process share a common stack \\
\addlinespace[3pt]
$\square$\ \textbf{(D)} Threads belonging to a process are by default not protected from each other
\end{tabular}

\tcblower

\textbf{Official Key:} \textbf{\large D} \hfill \textit{\small [Easy]}
\vspace{2pt}\hrule height 0.3pt\vspace{3pt}
\textbf{Topic:} Operating System \\
\textbf{Subtopic:} \textit{System calls, processes, threads, inter-process communication}
\vspace{2pt}\hrule height 0.3pt\vspace{3pt}
\textbf{Guide / Hint:} \\
{\small All threads of a single process share the same address space (code, global data, heap) and open files. Because they run within the same virtual address space without memory protection barriers between them, one thread can corrupt another thread's memory; hence they are by default not protected from each other (D is TRUE). Each thread maintains its own private stack and registers.}

\end{tcolorbox}
\vspace{6pt}\hrule height 0.35pt\vspace{6pt}
\noindent
\begin{tcolorbox}[
    blanker,
    sidebyside,
    sidebyside align=top,
    lefthand width=0.55\textwidth,
    top=4pt, bottom=4pt, left=0pt, right=0pt,
    middle=16pt,
    underlay={
        \draw[dashed, black!80, line width=0.8pt, dash pattern=on 4pt off 3pt]
            ([xshift=0.55\textwidth+8pt]frame.north west) --
            ([xshift=0.55\textwidth+8pt]frame.south west);
    }
]
\textbf{Q.136} \quad {\footnotesize\textsc{[MSQ $\cdot$ 1 Mark(s) $\cdot$ No Neg]}} \hfill {\footnotesize\textbf{GATE 2024 CS1} \quad \tiny\texttt{2024\_CS1\_CS\_Q25}}
\vspace{2pt}\hrule height 0.3pt\vspace{4pt}

Which of the following process state transitions is/are NOT possible?

\par\vspace{6pt}\noindent
\begin{tabular}{@{}p{0.48\linewidth}@{\hspace{0.04\linewidth}}p{0.48\linewidth}@{}}
$\square$\ \textbf{(A)} Running to Ready & $\square$\ \textbf{(B)} Waiting to Running \\
\addlinespace[3pt]
$\square$\ \textbf{(C)} Ready to Waiting & $\square$\ \textbf{(D)} Running to Terminated
\end{tabular}

\tcblower

\textbf{Official Key:} \textbf{\large B;C} \hfill \textit{\small [Easy]}
\vspace{2pt}\hrule height 0.3pt\vspace{3pt}
\textbf{Topic:} Operating System \\
\textbf{Subtopic:} \textit{System calls, processes, threads, inter-process communication}
\vspace{2pt}\hrule height 0.3pt\vspace{3pt}
\textbf{Guide / Hint:} \\
{\small In the standard process state model, a process in the Waiting (Blocked) state transitions to Ready upon event/I/O completion; it cannot transition directly to Running without CPU dispatch (B is impossible). Similarly, only a Running process can initiate an I/O request and transition to Waiting; a Ready process cannot transition directly to Waiting (C is impossible).}

\end{tcolorbox}
\vspace{6pt}\hrule height 0.35pt\vspace{6pt}
\noindent
\begin{tcolorbox}[
    blanker,
    sidebyside,
    sidebyside align=top,
    lefthand width=0.55\textwidth,
    top=4pt, bottom=4pt, left=0pt, right=0pt,
    middle=16pt,
    underlay={
        \draw[dashed, black!80, line width=0.8pt, dash pattern=on 4pt off 3pt]
            ([xshift=0.55\textwidth+8pt]frame.north west) --
            ([xshift=0.55\textwidth+8pt]frame.south west);
    }
]
\textbf{Q.137} \quad {\footnotesize\textsc{[NAT $\cdot$ 2 Mark(s) $\cdot$ No Neg]}} \hfill {\footnotesize\textbf{GATE 2024 CS1} \quad \tiny\texttt{2024\_CS1\_CS\_Q62}}
\vspace{2pt}\hrule height 0.3pt\vspace{4pt}

Consider a memory management system that uses a page size of 2 KB. Assume that both the physical and virtual addresses start from 0. Assume that the pages 0, 1, 2, and 3 are stored in the page frames 1, 3, 2, and 0, respectively. The physical address (in decimal format) corresponding to the virtual address 2500 (in decimal format) is \_\_\_\_\_\_\_\_\_\. (Answer in integer)

\par\vspace{6pt}\noindent\textbf{Your Answer:} \fbox{\rule{0pt}{1.3ex}\hspace{3cm}}

\tcblower

\textbf{Official Key:} \textbf{\large 6596 to 6596} \hfill \textit{\small [Easy]}
\vspace{2pt}\hrule height 0.3pt\vspace{3pt}
\textbf{Topic:} Operating System \\
\textbf{Subtopic:} \textit{Memory management and virtual memory}
\vspace{2pt}\hrule height 0.3pt\vspace{3pt}
\textbf{Guide / Hint:} \\
{\small With page size $2\text{ KB} = 2048\text{ bytes}$, virtual address $2500$ has virtual page number $\lfloor 2500 / 2048 \rfloor = 1$ and offset $2500 \pmod{2048} = 452$. Page 1 maps to frame 3. The physical address is $(\text{frame } 3 \times 2048) + \text{offset } 452 = 6144 + 452 = 6596$.}

\end{tcolorbox}
\vspace{6pt}\hrule height 0.35pt\vspace{6pt}
\vspace{8pt}
\subsection*{Databases}
\vspace{-4pt}
\noindent
\begin{tcolorbox}[
    blanker,
    sidebyside,
    sidebyside align=top,
    lefthand width=0.55\textwidth,
    top=4pt, bottom=4pt, left=0pt, right=0pt,
    middle=16pt,
    underlay={
        \draw[dashed, black!80, line width=0.8pt, dash pattern=on 4pt off 3pt]
            ([xshift=0.55\textwidth+8pt]frame.north west) --
            ([xshift=0.55\textwidth+8pt]frame.south west);
    }
]
\textbf{Q.138} \quad {\footnotesize\textsc{[MCQ $\cdot$ 1 Mark(s) $\cdot$ Neg: 0.33]}} \hfill {\footnotesize\textbf{GATE 2018 CS} \quad \tiny\texttt{2018\_CS\_GA\_Q11}}
\vspace{2pt}\hrule height 0.3pt\vspace{4pt}

In an Entity-Relationship (ER) model, suppose $R$ is a many-to-one relationship from entity set $E_1$ to entity set $E_2$. Assume that $E_1$ has total participation in $R$, but $E_2$ has partial participation in $R$. Which one of the following statements is necessarily TRUE?

\par\vspace{6pt}\noindent
\begin{tabular}{@{}p{0.98\linewidth}@{}}
$\square$\ \textbf{(A)} Every entity in $E_1$ is associated with exactly one entity in $E_2$. \\
\addlinespace[3pt]
$\square$\ \textbf{(B)} Some entity in $E_1$ is associated with more than one entity in $E_2$. \\
\addlinespace[3pt]
$\square$\ \textbf{(C)} Every entity in $E_2$ is associated with at least one entity in $E_1$. \\
\addlinespace[3pt]
$\square$\ \textbf{(D)} Every entity in $E_2$ is associated with at most one entity in $E_1$.
\end{tabular}

\tcblower

\textbf{Official Key:} \textbf{\large A} \hfill \textit{\small [Easy]}
\vspace{2pt}\hrule height 0.3pt\vspace{3pt}
\textbf{Topic:} Databases \\
\textbf{Subtopic:} \textit{ER-model}
\vspace{2pt}\hrule height 0.3pt\vspace{3pt}
\textbf{Guide / Hint:} \\
{\small In an ER diagram, a many-to-one relationship from $E_1$ to $E_2$ means each entity in $E_1$ is related to at most one entity in $E_2$. Total participation of $E_1$ requires that every entity in $E_1$ participates in at least one relationship. Combining these, every entity in $E_1$ is associated with exactly one entity in $E_2$.}

\end{tcolorbox}
\vspace{6pt}\hrule height 0.35pt\vspace{6pt}
\noindent
\begin{tcolorbox}[
    blanker,
    sidebyside,
    sidebyside align=top,
    lefthand width=0.55\textwidth,
    top=4pt, bottom=4pt, left=0pt, right=0pt,
    middle=16pt,
    underlay={
        \draw[dashed, black!80, line width=0.8pt, dash pattern=on 4pt off 3pt]
            ([xshift=0.55\textwidth+8pt]frame.north west) --
            ([xshift=0.55\textwidth+8pt]frame.south west);
    }
]
\textbf{Q.139} \quad {\footnotesize\textsc{[MCQ $\cdot$ 1 Mark(s) $\cdot$ Neg: 0.33]}} \hfill {\footnotesize\textbf{GATE 2018 CS} \quad \tiny\texttt{2018\_CS\_GA\_Q12}}
\vspace{2pt}\hrule height 0.3pt\vspace{4pt}

Consider the following two tables and four queries in SQL.
\texttt{Book (isbn, bname)}
\texttt{Stock (isbn, copies)}
\par\medskip
Query 1: \texttt{SELECT B.isbn, S.copies FROM Book B INNER JOIN Stock S ON B.isbn = S.isbn;}
Query 2: \texttt{SELECT B.isbn, S.copies FROM Book B LEFT OUTER JOIN Stock S ON B.isbn = S.isbn;}
Query 3: \texttt{SELECT B.isbn, S.copies FROM Book B RIGHT OUTER JOIN Stock S ON B.isbn = S.isbn;}
Query 4: \texttt{SELECT B.isbn, S.copies FROM Book B FULL OUTER JOIN Stock S ON B.isbn = S.isbn;}
\par\medskip
Which one of the following statements is CORRECT?

\par\vspace{6pt}\noindent
\begin{tabular}{@{}p{0.98\linewidth}@{}}
$\square$\ \textbf{(A)} Query 1 and Query 2 always produce the same number of rows. \\
\addlinespace[3pt]
$\square$\ \textbf{(B)} Query 2 and Query 3 always produce the same number of rows. \\
\addlinespace[3pt]
$\square$\ \textbf{(C)} Query 1 always produces the maximum number of rows among all queries. \\
\addlinespace[3pt]
$\square$\ \textbf{(D)} All rows from Query 1, Query 2, and Query 3 are always present in the result of Query 4.
\end{tabular}

\tcblower

\textbf{Official Key:} \textbf{\large D} \hfill \textit{\small [Easy]}
\vspace{2pt}\hrule height 0.3pt\vspace{3pt}
\textbf{Topic:} Databases \\
\textbf{Subtopic:} \textit{Relational model: relational algebra, tuple calculus, SQL}
\vspace{2pt}\hrule height 0.3pt\vspace{3pt}
\textbf{Guide / Hint:} \\
{\small Query 1 uses natural join (inner join), Query 2 uses left outer join, Query 3 uses right outer join, and Query 4 uses full outer join. A full outer join preserves all tuples from both relation $Book$ and relation $Stock$ (filling missing values with NULLs). Therefore, every tuple output by Query 1, Query 2, and Query 3 is guaranteed to appear in Query 4.}

\end{tcolorbox}
\vspace{6pt}\hrule height 0.35pt\vspace{6pt}
\noindent
\begin{tcolorbox}[
    blanker,
    sidebyside,
    sidebyside align=top,
    lefthand width=0.55\textwidth,
    top=4pt, bottom=4pt, left=0pt, right=0pt,
    middle=16pt,
    underlay={
        \draw[dashed, black!80, line width=0.8pt, dash pattern=on 4pt off 3pt]
            ([xshift=0.55\textwidth+8pt]frame.north west) --
            ([xshift=0.55\textwidth+8pt]frame.south west);
    }
]
\textbf{Q.140} \quad {\footnotesize\textsc{[MCQ $\cdot$ 1 Mark(s) $\cdot$ Neg: 0.33]}} \hfill {\footnotesize\textbf{GATE 2022 CS} \quad \tiny\texttt{2022\_CS\_CS\_Q14}}
\vspace{2pt}\hrule height 0.3pt\vspace{4pt}

In a relational data model, which one of the following statements is TRUE?

\par\vspace{6pt}\noindent
\begin{tabular}{@{}p{0.98\linewidth}@{}}
$\square$\ \textbf{(A)} A relation with only two attributes is always in BCNF. \\
\addlinespace[3pt]
$\square$\ \textbf{(B)} If all attributes of a relation are prime attributes, then the relation is in BCNF. \\
\addlinespace[3pt]
$\square$\ \textbf{(C)} Every relation has at least one non-prime attribute. \\
\addlinespace[3pt]
$\square$\ \textbf{(D)} BCNF decompositions preserve functional dependencies.
\end{tabular}

\tcblower

\textbf{Official Key:} \textbf{\large A} \hfill \textit{\small [Easy]}
\vspace{2pt}\hrule height 0.3pt\vspace{3pt}
\textbf{Topic:} Databases \\
\textbf{Subtopic:} \textit{Integrity constraints, normal forms}
\vspace{2pt}\hrule height 0.3pt\vspace{3pt}
\textbf{Guide / Hint:} \\
{\small For any binary relation $R(A, B)$, the only possible non-trivial functional dependencies are $A \to B$ or $B \to A$. In either case, the left-hand side is a superkey of $R$, satisfying the definition of Boyce-Codd Normal Form (BCNF).}

\end{tcolorbox}
\vspace{6pt}\hrule height 0.35pt\vspace{6pt}
\noindent
\begin{tcolorbox}[
    blanker,
    sidebyside,
    sidebyside align=top,
    lefthand width=0.55\textwidth,
    top=4pt, bottom=4pt, left=0pt, right=0pt,
    middle=16pt,
    underlay={
        \draw[dashed, black!80, line width=0.8pt, dash pattern=on 4pt off 3pt]
            ([xshift=0.55\textwidth+8pt]frame.north west) --
            ([xshift=0.55\textwidth+8pt]frame.south west);
    }
]
\textbf{Q.141} \quad {\footnotesize\textsc{[NAT $\cdot$ 1 Mark(s) $\cdot$ No Neg]}} \hfill {\footnotesize\textbf{GATE 2022 CS} \quad \tiny\texttt{2022\_CS\_CS\_Q31}}
\vspace{2pt}\hrule height 0.3pt\vspace{4pt}

Consider a relation $R(A, B, C, D, E)$ with the following three functional dependencies:
\[ AB \to C; \quad BC \to D; \quad C \to E; \]
The number of superkeys in the relation $R$ is \_\_\_\_\_\_\_\_\_\_\_\. (Answer in integer)

\par\vspace{6pt}\noindent\textbf{Your Answer:} \fbox{\rule{0pt}{1.3ex}\hspace{3cm}}

\tcblower

\textbf{Official Key:} \textbf{\large 8 to 8} \hfill \textit{\small [Easy]}
\vspace{2pt}\hrule height 0.3pt\vspace{3pt}
\textbf{Topic:} Databases \\
\textbf{Subtopic:} \textit{Integrity constraints, normal forms}
\vspace{2pt}\hrule height 0.3pt\vspace{3pt}
\textbf{Guide / Hint:} \\
{\small Computing attribute closures: $(AB)^+ = \{A, B, C, D, E\}$. Attributes $A$ and $B$ never appear on any RHS, so $AB$ is the unique candidate key. Every superkey must contain $AB$, and can include any subset of the remaining 3 attributes $\{C, D, E\}$, giving $2^3 = 8$ superkeys.}

\end{tcolorbox}
\vspace{6pt}\hrule height 0.35pt\vspace{6pt}
\noindent
\begin{tcolorbox}[
    blanker,
    sidebyside,
    sidebyside align=top,
    lefthand width=0.55\textwidth,
    top=4pt, bottom=4pt, left=0pt, right=0pt,
    middle=16pt,
    underlay={
        \draw[dashed, black!80, line width=0.8pt, dash pattern=on 4pt off 3pt]
            ([xshift=0.55\textwidth+8pt]frame.north west) --
            ([xshift=0.55\textwidth+8pt]frame.south west);
    }
]
\textbf{Q.142} \quad {\footnotesize\textsc{[MCQ $\cdot$ 1 Mark(s) $\cdot$ Neg: 0.33]}} \hfill {\footnotesize\textbf{GATE 2023 CS} \quad \tiny\texttt{2023\_CS\_CS\_Q16}}
\vspace{2pt}\hrule height 0.3pt\vspace{4pt}

Which one of the options given below refers to the degree (or arity) of a relation in relational database systems?

\par\vspace{6pt}\noindent
\begin{tabular}{@{}p{0.98\linewidth}@{}}
$\square$\ \textbf{(A)} Number of attributes of its relation schema. \\
\addlinespace[3pt]
$\square$\ \textbf{(B)} Number of tuples stored in the relation. \\
\addlinespace[3pt]
$\square$\ \textbf{(C)} Number of entries in the relation. \\
\addlinespace[3pt]
$\square$\ \textbf{(D)} Number of distinct domains of its relation schema.
\end{tabular}

\tcblower

\textbf{Official Key:} \textbf{\large A} \hfill \textit{\small [Easy]}
\vspace{2pt}\hrule height 0.3pt\vspace{3pt}
\textbf{Topic:} Databases \\
\textbf{Subtopic:} \textit{Relational model: relational algebra, tuple calculus, SQL}
\vspace{2pt}\hrule height 0.3pt\vspace{3pt}
\textbf{Guide / Hint:} \\
{\small In relational database theory, the degree (or arity) of a relation is defined as the number of attributes (columns) in its relation schema, whereas the cardinality of a relation is the number of tuples (rows).}

\end{tcolorbox}
\vspace{6pt}\hrule height 0.35pt\vspace{6pt}
\noindent
\begin{tcolorbox}[
    blanker,
    sidebyside,
    sidebyside align=top,
    lefthand width=0.55\textwidth,
    top=4pt, bottom=4pt, left=0pt, right=0pt,
    middle=16pt,
    underlay={
        \draw[dashed, black!80, line width=0.8pt, dash pattern=on 4pt off 3pt]
            ([xshift=0.55\textwidth+8pt]frame.north west) --
            ([xshift=0.55\textwidth+8pt]frame.south west);
    }
]
\textbf{Q.143} \quad {\footnotesize\textsc{[NAT $\cdot$ 2 Mark(s) $\cdot$ No Neg]}} \hfill {\footnotesize\textbf{GATE 2023 CS} \quad \tiny\texttt{2023\_CS\_CS\_Q61}}
\vspace{2pt}\hrule height 0.3pt\vspace{4pt}

Consider the following table named `Student` in a relational database. The primary key of this table is `rollNum`.
\begin{center}
\begin{tabular}{|c|c|c|c|}
\hline
\textbf{rollNum} & \textbf{name} & \textbf{gender} & \textbf{marks} \\
\hline
1 & Naman & M & 62 \\
2 & Aliya & F & 70 \\
3 & Aliya & F & 80 \\
4 & James & M & 82 \\
5 & Swati & F & 65 \\
\hline
\end{tabular}
\end{center}
The SQL query below is executed on this database:
\begin{lstlisting}[language=SQL]
SELECT *
FROM Student
WHERE gender = 'F' AND marks > 65;
\end{lstlisting}
The number of rows returned by the query is \_\_\_\_\_\_\_\_\. (Answer in integer)

\par\vspace{6pt}\noindent\textbf{Your Answer:} \fbox{\rule{0pt}{1.3ex}\hspace{3cm}}

\tcblower

\textbf{Official Key:} \textbf{\large 2 to 2} \hfill \textit{\small [Easy]}
\vspace{2pt}\hrule height 0.3pt\vspace{3pt}
\textbf{Topic:} Databases \\
\textbf{Subtopic:} \textit{Relational model: relational algebra, tuple calculus, SQL}
\vspace{2pt}\hrule height 0.3pt\vspace{3pt}
\textbf{Guide / Hint:} \\
{\small The query evaluates the WHERE clause condition `gender = 'F' AND marks > 65`. Among female students: Aliya (rollNum 2, marks 70) and Aliya (rollNum 3, marks 80) satisfy both conditions, whereas Swati (rollNum 5, marks 65) fails the strict inequality. Thus, exactly 2 rows are returned.}

\end{tcolorbox}
\vspace{6pt}\hrule height 0.35pt\vspace{6pt}
\noindent
\begin{tcolorbox}[
    blanker,
    sidebyside,
    sidebyside align=top,
    lefthand width=0.55\textwidth,
    top=4pt, bottom=4pt, left=0pt, right=0pt,
    middle=16pt,
    underlay={
        \draw[dashed, black!80, line width=0.8pt, dash pattern=on 4pt off 3pt]
            ([xshift=0.55\textwidth+8pt]frame.north west) --
            ([xshift=0.55\textwidth+8pt]frame.south west);
    }
]
\textbf{Q.144} \quad {\footnotesize\textsc{[MCQ $\cdot$ 1 Mark(s) $\cdot$ Neg: 0.33]}} \hfill {\footnotesize\textbf{GATE 2024 CS1} \quad \tiny\texttt{2024\_CS1\_CS\_Q21}}
\vspace{2pt}\hrule height 0.3pt\vspace{4pt}

In a $\text{B}^+$ tree, the requirement of at least half-full (50\%) node occupancy is relaxed for which one of the following cases?

\par\vspace{6pt}\noindent
\begin{tabular}{@{}p{0.48\linewidth}@{\hspace{0.04\linewidth}}p{0.48\linewidth}@{}}
$\square$\ \textbf{(A)} Only the root node & $\square$\ \textbf{(B)} All leaf nodes \\
\addlinespace[3pt]
$\square$\ \textbf{(C)} All internal nodes & $\square$\ \textbf{(D)} Only the leftmost leaf node
\end{tabular}

\tcblower

\textbf{Official Key:} \textbf{\large A} \hfill \textit{\small [Easy]}
\vspace{2pt}\hrule height 0.3pt\vspace{3pt}
\textbf{Topic:} Databases \\
\textbf{Subtopic:} \textit{File organization, indexing (e.g., B and B+ trees)}
\vspace{2pt}\hrule height 0.3pt\vspace{3pt}
\textbf{Guide / Hint:} \\
{\small In a $\text{B}^+$ tree, every node except the root node must satisfy the minimum occupancy constraint of being at least half full ($	ext{at least } \lceil p/2 \rceil$ child pointers). The root node is the sole exception, and is permitted to hold as few as two pointers (or one key) when non-leaf, and one key when leaf.}

\end{tcolorbox}
\vspace{6pt}\hrule height 0.35pt\vspace{6pt}
\noindent
\begin{tcolorbox}[
    blanker,
    sidebyside,
    sidebyside align=top,
    lefthand width=0.55\textwidth,
    top=4pt, bottom=4pt, left=0pt, right=0pt,
    middle=16pt,
    underlay={
        \draw[dashed, black!80, line width=0.8pt, dash pattern=on 4pt off 3pt]
            ([xshift=0.55\textwidth+8pt]frame.north west) --
            ([xshift=0.55\textwidth+8pt]frame.south west);
    }
]
\textbf{Q.145} \quad {\footnotesize\textsc{[MSQ $\cdot$ 1 Mark(s) $\cdot$ No Neg]}} \hfill {\footnotesize\textbf{GATE 2024 CS1} \quad \tiny\texttt{2024\_CS1\_CS\_Q22}}
\vspace{2pt}\hrule height 0.3pt\vspace{4pt}

Which of the following statements about a relation $R$ in first normal form (1NF) is/are TRUE?

\par\vspace{6pt}\noindent
\begin{tabular}{@{}p{0.98\linewidth}@{}}
$\square$\ \textbf{(A)} $R$ can have a multi-attribute key \\
\addlinespace[3pt]
$\square$\ \textbf{(B)} $R$ cannot have a foreign key \\
\addlinespace[3pt]
$\square$\ \textbf{(C)} $R$ cannot have a composite attribute \\
\addlinespace[3pt]
$\square$\ \textbf{(D)} $R$ cannot have more than one candidate key
\end{tabular}

\tcblower

\textbf{Official Key:} \textbf{\large A;C} \hfill \textit{\small [Easy]}
\vspace{2pt}\hrule height 0.3pt\vspace{3pt}
\textbf{Topic:} Databases \\
\textbf{Subtopic:} \textit{Integrity constraints, normal forms}
\vspace{2pt}\hrule height 0.3pt\vspace{3pt}
\textbf{Guide / Hint:} \\
{\small A relation is in 1NF if and only if all attribute values are atomic, which explicitly forbids composite or multi-valued attributes (making (C) TRUE). 1NF places no limitation on candidate keys having multiple attributes (multi-attribute keys are fully permitted, making (A) TRUE) or the presence of foreign keys.}

\end{tcolorbox}
\vspace{6pt}\hrule height 0.35pt\vspace{6pt}
\noindent
\begin{tcolorbox}[
    blanker,
    sidebyside,
    sidebyside align=top,
    lefthand width=0.55\textwidth,
    top=4pt, bottom=4pt, left=0pt, right=0pt,
    middle=16pt,
    underlay={
        \draw[dashed, black!80, line width=0.8pt, dash pattern=on 4pt off 3pt]
            ([xshift=0.55\textwidth+8pt]frame.north west) --
            ([xshift=0.55\textwidth+8pt]frame.south west);
    }
]
\textbf{Q.146} \quad {\footnotesize\textsc{[NAT $\cdot$ 1 Mark(s) $\cdot$ No Neg]}} \hfill {\footnotesize\textbf{GATE 2024 CS1} \quad \tiny\texttt{2024\_CS1\_CS\_Q35}}
\vspace{2pt}\hrule height 0.3pt\vspace{4pt}

Consider the following two relations, $R(A,B)$ and $S(A,C)$:
\begin{center}
\begin{tabular}{|c|c|}
\multicolumn{2}{c}{$R$} \\
\hline
$A$ & $B$ \\
\hline
10 & 20 \\
20 & 30 \\
30 & 40 \\
30 & 50 \\
\hline
\end{tabular}
\qquad
\begin{tabular}{|c|c|}
\multicolumn{2}{c}{$S$} \\
\hline
$A$ & $C$ \\
\hline
10 & 90 \\
30 & 45 \\
40 & 80 \\
50 & 95 \\
\hline
\end{tabular}
\end{center}
The total number of tuples obtained by evaluating the following expression
\[ \sigma\_{B < C}(R \bowtie\_{R.A = S.A} S) \]
is \_\_\_\_\_\_\_\_\_\. (Answer in integer)

\par\vspace{6pt}\noindent\textbf{Your Answer:} \fbox{\rule{0pt}{1.3ex}\hspace{3cm}}

\tcblower

\textbf{Official Key:} \textbf{\large 2 to 2} \hfill \textit{\small [Easy]}
\vspace{2pt}\hrule height 0.3pt\vspace{3pt}
\textbf{Topic:} Databases \\
\textbf{Subtopic:} \textit{Relational model: relational algebra, tuple calculus, SQL}
\vspace{2pt}\hrule height 0.3pt\vspace{3pt}
\textbf{Guide / Hint:} \\
{\small Joining $R(A,B)$ and $S(A,C)$ on $R.A = S.A$ yields: $(10, 20, 90)$ for $A=10$, and $(30, 40, 45)$, $(30, 50, 45)$ for $A=30$. Applying selection $\sigma_{B < C}$: $20 < 90$ is true (tuple $(10, 20, 90)$), $40 < 45$ is true (tuple $(30, 40, 45)$), and $50 < 45$ is false. Thus, exactly 2 tuples satisfy the condition.}

\end{tcolorbox}
\vspace{6pt}\hrule height 0.35pt\vspace{6pt}
\noindent
\begin{tcolorbox}[
    blanker,
    sidebyside,
    sidebyside align=top,
    lefthand width=0.55\textwidth,
    top=4pt, bottom=4pt, left=0pt, right=0pt,
    middle=16pt,
    underlay={
        \draw[dashed, black!80, line width=0.8pt, dash pattern=on 4pt off 3pt]
            ([xshift=0.55\textwidth+8pt]frame.north west) --
            ([xshift=0.55\textwidth+8pt]frame.south west);
    }
]
\textbf{Q.147} \quad {\footnotesize\textsc{[MCQ $\cdot$ 1 Mark(s) $\cdot$ Neg: 0.33]}} \hfill {\footnotesize\textbf{GATE 2024 CS2} \quad \tiny\texttt{2024\_CS2\_CS\_Q19}}
\vspace{2pt}\hrule height 0.3pt\vspace{4pt}

Once the DBMS informs the user that a transaction has been successfully completed, its effect should persist even if the system crashes before all its changes are reflected on disk. This property is called:

\par\vspace{6pt}\noindent
\begin{tabular}{@{}p{0.48\linewidth}@{\hspace{0.04\linewidth}}p{0.48\linewidth}@{}}
$\square$\ \textbf{(A)} durability & $\square$\ \textbf{(B)} atomicity \\
\addlinespace[3pt]
$\square$\ \textbf{(C)} consistency & $\square$\ \textbf{(D)} isolation
\end{tabular}

\tcblower

\textbf{Official Key:} \textbf{\large A} \hfill \textit{\small [Easy]}
\vspace{2pt}\hrule height 0.3pt\vspace{3pt}
\textbf{Topic:} Databases \\
\textbf{Subtopic:} \textit{Transactions and concurrency control}
\vspace{2pt}\hrule height 0.3pt\vspace{3pt}
\textbf{Guide / Hint:} \\
{\small The Durability property of the ACID paradigm guarantees that once a transaction has committed successfully, its updates and effects survive permanently in non-volatile storage, even in the event of a subsequent system crash.}

\end{tcolorbox}
\vspace{6pt}\hrule height 0.35pt\vspace{6pt}
\noindent
\begin{tcolorbox}[
    blanker,
    sidebyside,
    sidebyside align=top,
    lefthand width=0.55\textwidth,
    top=4pt, bottom=4pt, left=0pt, right=0pt,
    middle=16pt,
    underlay={
        \draw[dashed, black!80, line width=0.8pt, dash pattern=on 4pt off 3pt]
            ([xshift=0.55\textwidth+8pt]frame.north west) --
            ([xshift=0.55\textwidth+8pt]frame.south west);
    }
]
\textbf{Q.148} \quad {\footnotesize\textsc{[MCQ $\cdot$ 1 Mark(s) $\cdot$ Neg: 0.33]}} \hfill {\footnotesize\textbf{GATE 2024 CS2} \quad \tiny\texttt{2024\_CS2\_CS\_Q20}}
\vspace{2pt}\hrule height 0.3pt\vspace{4pt}

In the context of owner and weak entity sets in the ER (Entity-Relationship) data model, which one of the following statements is TRUE?

\par\vspace{6pt}\noindent
\begin{tabular}{@{}p{0.98\linewidth}@{}}
$\square$\ \textbf{(A)} The weak entity set MUST have total participation in the identifying relationship \\
\addlinespace[3pt]
$\square$\ \textbf{(B)} The owner entity set MUST have total participation in the identifying relationship \\
\addlinespace[3pt]
$\square$\ \textbf{(C)} Both weak and owner entity sets MUST have total participation in the identifying relationship \\
\addlinespace[3pt]
$\square$\ \textbf{(D)} Neither weak entity set nor owner entity set MUST have total participation in the identifying relationship
\end{tabular}

\tcblower

\textbf{Official Key:} \textbf{\large A} \hfill \textit{\small [Easy]}
\vspace{2pt}\hrule height 0.3pt\vspace{3pt}
\textbf{Topic:} Databases \\
\textbf{Subtopic:} \textit{ER-model}
\vspace{2pt}\hrule height 0.3pt\vspace{3pt}
\textbf{Guide / Hint:} \\
{\small A weak entity cannot exist without its identifying (owner) entity. Therefore, every instance of the weak entity set must be associated with an owner entity instance via the identifying relationship, mandating total participation of the weak entity set.}

\end{tcolorbox}
\vspace{6pt}\hrule height 0.35pt\vspace{6pt}
\noindent
\begin{tcolorbox}[
    blanker,
    sidebyside,
    sidebyside align=top,
    lefthand width=0.55\textwidth,
    top=4pt, bottom=4pt, left=0pt, right=0pt,
    middle=16pt,
    underlay={
        \draw[dashed, black!80, line width=0.8pt, dash pattern=on 4pt off 3pt]
            ([xshift=0.55\textwidth+8pt]frame.north west) --
            ([xshift=0.55\textwidth+8pt]frame.south west);
    }
]
\textbf{Q.149} \quad {\footnotesize\textsc{[MSQ $\cdot$ 1 Mark(s) $\cdot$ No Neg]}} \hfill {\footnotesize\textbf{GATE 2024 CS2} \quad \tiny\texttt{2024\_CS2\_CS\_Q26}}
\vspace{2pt}\hrule height 0.3pt\vspace{4pt}

Which of the following file organizations is/are I/O efficient for the scan operation in DBMS?

\par\vspace{6pt}\noindent
\begin{tabular}{@{}p{0.48\linewidth}@{\hspace{0.04\linewidth}}p{0.48\linewidth}@{}}
$\square$\ \textbf{(A)} Sorted & $\square$\ \textbf{(B)} Heap \\
\addlinespace[3pt]
$\square$\ \textbf{(C)} Unclustered tree index & $\square$\ \textbf{(D)} Unclustered hash index
\end{tabular}

\tcblower

\textbf{Official Key:} \textbf{\large A;B} \hfill \textit{\small [Easy]}
\vspace{2pt}\hrule height 0.3pt\vspace{3pt}
\textbf{Topic:} Databases \\
\textbf{Subtopic:} \textit{File organization, indexing (e.g., B and B+ trees)}
\vspace{2pt}\hrule height 0.3pt\vspace{3pt}
\textbf{Guide / Hint:} \\
{\small A full table scan retrieves all pages sequentially. Sorted and Heap files store data records densely packed across consecutive data pages, requiring reading each page exactly once. Unclustered indexes store data out of order, incurring multiple random I/Os per record.}

\end{tcolorbox}
\vspace{6pt}\hrule height 0.35pt\vspace{6pt}
\vspace{8pt}
\subsection*{Computer Networks}
\vspace{-4pt}
\noindent
\begin{tcolorbox}[
    blanker,
    sidebyside,
    sidebyside align=top,
    lefthand width=0.55\textwidth,
    top=4pt, bottom=4pt, left=0pt, right=0pt,
    middle=16pt,
    underlay={
        \draw[dashed, black!80, line width=0.8pt, dash pattern=on 4pt off 3pt]
            ([xshift=0.55\textwidth+8pt]frame.north west) --
            ([xshift=0.55\textwidth+8pt]frame.south west);
    }
]
\textbf{Q.150} \quad {\footnotesize\textsc{[MCQ $\cdot$ 1 Mark(s) $\cdot$ Neg: 0.33]}} \hfill {\footnotesize\textbf{GATE 2018 CS} \quad \tiny\texttt{2018\_CS\_GA\_Q13}}
\vspace{2pt}\hrule height 0.3pt\vspace{4pt}

Match the following:
\begin{center}
\begin{tabular}{ll@{\hspace{2cm}}ll}
\textbf{Field} & & \textbf{Length in bits} & \\
P. & UDP Header's Port Number & I. & 48 \\
Q. & Ethernet MAC Address & II. & 8 \\
R. & IPv6 Next Header & III. & 32 \\
S. & TCP Header's Sequence Number & IV. & 16
\end{tabular}
\end{center}

\par\vspace{6pt}\noindent
\begin{tabular}{@{}p{0.48\linewidth}@{\hspace{0.04\linewidth}}p{0.48\linewidth}@{}}
$\square$\ \textbf{(A)} P-III, Q-I, R-II, S-IV & $\square$\ \textbf{(B)} P-IV, Q-II, R-I, S-III \\
\addlinespace[3pt]
$\square$\ \textbf{(C)} P-IV, Q-I, R-II, S-III & $\square$\ \textbf{(D)} P-I, Q-IV, R-II, S-III
\end{tabular}

\tcblower

\textbf{Official Key:} \textbf{\large C} \hfill \textit{\small [Easy]}
\vspace{2pt}\hrule height 0.3pt\vspace{3pt}
\textbf{Topic:} Computer Networks \\
\textbf{Subtopic:} \textit{Concept of layering: OSI and TCP/IP Protocol Stacks}
\vspace{2pt}\hrule height 0.3pt\vspace{3pt}
\textbf{Guide / Hint:} \\
{\small UDP port numbers are 16 bits long (P-IV). Ethernet MAC addresses are 48 bits long (Q-I). The IPv6 Next Header field is 8 bits long (R-II). TCP sequence numbers are 32 bits long (S-III). Hence the correct matching is P-IV, Q-I, R-II, S-III.}

\end{tcolorbox}
\vspace{6pt}\hrule height 0.35pt\vspace{6pt}
\noindent
\begin{tcolorbox}[
    blanker,
    sidebyside,
    sidebyside align=top,
    lefthand width=0.55\textwidth,
    top=4pt, bottom=4pt, left=0pt, right=0pt,
    middle=16pt,
    underlay={
        \draw[dashed, black!80, line width=0.8pt, dash pattern=on 4pt off 3pt]
            ([xshift=0.55\textwidth+8pt]frame.north west) --
            ([xshift=0.55\textwidth+8pt]frame.south west);
    }
]
\textbf{Q.151} \quad {\footnotesize\textsc{[NAT $\cdot$ 1 Mark(s) $\cdot$ No Neg]}} \hfill {\footnotesize\textbf{GATE 2022 CS} \quad \tiny\texttt{2022\_CS\_CS\_Q35}}
\vspace{2pt}\hrule height 0.3pt\vspace{4pt}

Consider the resolution of the domain name `www.gate.org.in` by a DNS resolver. Assume that no resource records are cached anywhere across the DNS servers and that iterative query mechanism is used in the resolution. The number of DNS query-response pairs involved in completely resolving the domain name is \_\_\_\_\_\_\_\_\_\_\_\_\. (Answer in integer)

\par\vspace{6pt}\noindent\textbf{Your Answer:} \fbox{\rule{0pt}{1.3ex}\hspace{3cm}}

\tcblower

\textbf{Official Key:} \textbf{\large 4 to 4} \hfill \textit{\small [Easy]}
\vspace{2pt}\hrule height 0.3pt\vspace{3pt}
\textbf{Topic:} Computer Networks \\
\textbf{Subtopic:} \textit{Application layer protocols: DNS, SMTP, HTTP, FTP, Email}
\vspace{2pt}\hrule height 0.3pt\vspace{3pt}
\textbf{Guide / Hint:} \\
{\small In an un-cached iterative resolution for `www.gate.org.in`: the resolver queries Root DNS server (1), `.in` TLD server (2), `.org.in` SLD server (3), and `gate.org.in` Authoritative server (4). Each query yields a response, totaling 4 query-response pairs.}

\end{tcolorbox}
\vspace{6pt}\hrule height 0.35pt\vspace{6pt}
\noindent
\begin{tcolorbox}[
    blanker,
    sidebyside,
    sidebyside align=top,
    lefthand width=0.55\textwidth,
    top=4pt, bottom=4pt, left=0pt, right=0pt,
    middle=16pt,
    underlay={
        \draw[dashed, black!80, line width=0.8pt, dash pattern=on 4pt off 3pt]
            ([xshift=0.55\textwidth+8pt]frame.north west) --
            ([xshift=0.55\textwidth+8pt]frame.south west);
    }
]
\textbf{Q.152} \quad {\footnotesize\textsc{[NAT $\cdot$ 2 Mark(s) $\cdot$ No Neg]}} \hfill {\footnotesize\textbf{GATE 2022 CS} \quad \tiny\texttt{2022\_CS\_CS\_Q59}}
\vspace{2pt}\hrule height 0.3pt\vspace{4pt}

Consider a 100 Mbps link between an earth station (sender) and a satellite (receiver) at an altitude of 2100 km. The signal propagates at a speed of $3 \times 10^8\text{ m/s}$. The time taken (in milliseconds, rounded off to two decimal places) for the receiver to completely receive a packet of 1000 bytes transmitted by the sender is \_\_\_\_\_\_\_\_\_\. (Answer in numeric format)

\par\vspace{6pt}\noindent\textbf{Your Answer:} \fbox{\rule{0pt}{1.3ex}\hspace{3cm}}

\tcblower

\textbf{Official Key:} \textbf{\large 7.07 to 7.09} \hfill \textit{\small [Easy]}
\vspace{2pt}\hrule height 0.3pt\vspace{3pt}
\textbf{Topic:} Computer Networks \\
\textbf{Subtopic:} \textit{Data link layer: framing, error detection, Medium Access Control, Ethernet bridging}
\vspace{2pt}\hrule height 0.3pt\vspace{3pt}
\textbf{Guide / Hint:} \\
{\small Propagation time is $T_{\text{prop}} = \frac{2100 \times 10^3\text{ m}}{3 \times 10^8\text{ m/s}} = 7.00\text{ ms}$. Transmission time is $T_{\text{trans}} = \frac{1000 \times 8\text{ bits}}{100 \times 10^6\text{ bps}} = 0.08\text{ ms}$. Total time to completely receive the packet is $T_{\text{prop}} + T_{\text{trans}} = 7.00 + 0.08 = 7.08\text{ ms}$.}

\end{tcolorbox}
\vspace{6pt}\hrule height 0.35pt\vspace{6pt}
\noindent
\begin{tcolorbox}[
    blanker,
    sidebyside,
    sidebyside align=top,
    lefthand width=0.55\textwidth,
    top=4pt, bottom=4pt, left=0pt, right=0pt,
    middle=16pt,
    underlay={
        \draw[dashed, black!80, line width=0.8pt, dash pattern=on 4pt off 3pt]
            ([xshift=0.55\textwidth+8pt]frame.north west) --
            ([xshift=0.55\textwidth+8pt]frame.south west);
    }
]
\textbf{Q.153} \quad {\footnotesize\textsc{[MSQ $\cdot$ 1 Mark(s) $\cdot$ No Neg]}} \hfill {\footnotesize\textbf{GATE 2023 CS} \quad \tiny\texttt{2023\_CS\_CS\_Q25}}
\vspace{2pt}\hrule height 0.3pt\vspace{4pt}

Which of the following statements is/are INCORRECT about the OSPF (Open Shortest Path First) routing protocol used in the Internet?

\par\vspace{6pt}\noindent
\begin{tabular}{@{}p{0.98\linewidth}@{}}
$\square$\ \textbf{(A)} OSPF implements Bellman-Ford algorithm to find shortest paths. \\
\addlinespace[3pt]
$\square$\ \textbf{(B)} OSPF uses Dijkstra's shortest path algorithm to implement least-cost path routing. \\
\addlinespace[3pt]
$\square$\ \textbf{(C)} OSPF is used as an inter-domain routing protocol. \\
\addlinespace[3pt]
$\square$\ \textbf{(D)} OSPF implements hierarchical routing.
\end{tabular}

\tcblower

\textbf{Official Key:} \textbf{\large A, C} \hfill \textit{\small [Easy]}
\vspace{2pt}\hrule height 0.3pt\vspace{3pt}
\textbf{Topic:} Computer Networks \\
\textbf{Subtopic:} \textit{Routing protocols: shortest path, flooding, distance vector and link state routing}
\vspace{2pt}\hrule height 0.3pt\vspace{3pt}
\textbf{Guide / Hint:} \\
{\small OSPF is an intra-domain (Interior Gateway) routing protocol that uses Dijkstra's link-state algorithm (not Bellman-Ford, which is used by distance-vector protocols like RIP). Therefore, statement A (implements Bellman-Ford) and statement C (used as inter-domain) are INCORRECT, which are the required answers.}

\end{tcolorbox}
\vspace{6pt}\hrule height 0.35pt\vspace{6pt}
\noindent
\begin{tcolorbox}[
    blanker,
    sidebyside,
    sidebyside align=top,
    lefthand width=0.55\textwidth,
    top=4pt, bottom=4pt, left=0pt, right=0pt,
    middle=16pt,
    underlay={
        \draw[dashed, black!80, line width=0.8pt, dash pattern=on 4pt off 3pt]
            ([xshift=0.55\textwidth+8pt]frame.north west) --
            ([xshift=0.55\textwidth+8pt]frame.south west);
    }
]
\textbf{Q.154} \quad {\footnotesize\textsc{[MCQ $\cdot$ 1 Mark(s) $\cdot$ Neg: 0.33]}} \hfill {\footnotesize\textbf{GATE 2024 CS1} \quad \tiny\texttt{2024\_CS1\_CS\_Q16}}
\vspace{2pt}\hrule height 0.3pt\vspace{4pt}

A user starts browsing a webpage hosted at a remote server. The browser opens a single TCP connection to fetch the entire webpage from the server. The webpage consists of a top-level index page with multiple embedded image objects. Assume that all caches (e.g., DNS cache, browser cache) are all initially empty. The following packets leave the user's computer in some order:
\begin{enumerate}[label=(\roman*),itemsep=0pt,topsep=2pt]
\item HTTP GET request for the index page
\item DNS request to resolve the web server's name to its IP address
\item HTTP GET request for an image object
\item TCP SYN to open a connection to the web server
\end{enumerate}
Which one of the following is the CORRECT chronological order (earliest in time to latest) of the packets leaving the computer?

\par\vspace{6pt}\noindent
\begin{tabular}{@{}p{0.48\linewidth}@{\hspace{0.04\linewidth}}p{0.48\linewidth}@{}}
$\square$\ \textbf{(A)} (iv), (ii), (iii), (i) & $\square$\ \textbf{(B)} (ii), (iv), (iii), (i) \\
\addlinespace[3pt]
$\square$\ \textbf{(C)} (ii), (iv), (i), (iii) & $\square$\ \textbf{(D)} (iv), (ii), (i), (iii)
\end{tabular}

\tcblower

\textbf{Official Key:} \textbf{\large C} \hfill \textit{\small [Easy]}
\vspace{2pt}\hrule height 0.3pt\vspace{3pt}
\textbf{Topic:} Computer Networks \\
\textbf{Subtopic:} \textit{Concept of layering: OSI and TCP/IP Protocol Stacks}
\vspace{2pt}\hrule height 0.3pt\vspace{3pt}
\textbf{Guide / Hint:} \\
{\small Initially all caches are empty. The client must first send a DNS query (ii) to resolve the server hostname to an IP. Next, it establishes a TCP connection by sending TCP SYN (iv). Once connected, it sends an HTTP GET request for the index HTML page (i). After parsing the HTML and identifying embedded images, it sends HTTP GET requests for the images (iii).}

\end{tcolorbox}
\vspace{6pt}\hrule height 0.35pt\vspace{6pt}
\noindent
\begin{tcolorbox}[
    blanker,
    sidebyside,
    sidebyside align=top,
    lefthand width=0.55\textwidth,
    top=4pt, bottom=4pt, left=0pt, right=0pt,
    middle=16pt,
    underlay={
        \draw[dashed, black!80, line width=0.8pt, dash pattern=on 4pt off 3pt]
            ([xshift=0.55\textwidth+8pt]frame.north west) --
            ([xshift=0.55\textwidth+8pt]frame.south west);
    }
]
\textbf{Q.155} \quad {\footnotesize\textsc{[MSQ $\cdot$ 1 Mark(s) $\cdot$ No Neg]}} \hfill {\footnotesize\textbf{GATE 2024 CS1} \quad \tiny\texttt{2024\_CS1\_CS\_Q29}}
\vspace{2pt}\hrule height 0.3pt\vspace{4pt}

TCP client P successfully establishes a connection to TCP server Q. Let $N_P$ denote the sequence number in the SYN sent from P to Q. Let $N_Q$ denote the acknowledgement number in the SYN ACK from Q to P. Which of the following statements is/are CORRECT?

\par\vspace{6pt}\noindent
\begin{tabular}{@{}p{0.98\linewidth}@{}}
$\square$\ \textbf{(A)} The sequence number $N_P$ is chosen randomly by P \\
\addlinespace[3pt]
$\square$\ \textbf{(B)} The sequence number $N_P$ is always 0 for a new connection \\
\addlinespace[3pt]
$\square$\ \textbf{(C)} The acknowledgement number $N_Q$ is equal to $N_P$ \\
\addlinespace[3pt]
$\square$\ \textbf{(D)} The acknowledgement number $N_Q$ is equal to $N_P + 1$
\end{tabular}

\tcblower

\textbf{Official Key:} \textbf{\large A;D} \hfill \textit{\small [Easy]}
\vspace{2pt}\hrule height 0.3pt\vspace{3pt}
\textbf{Topic:} Computer Networks \\
\textbf{Subtopic:} \textit{Transport layer: flow control and congestion control, UDP, TCP, sockets}
\vspace{2pt}\hrule height 0.3pt\vspace{3pt}
\textbf{Guide / Hint:} \\
{\small During TCP connection establishment, the initial sequence number $N_P$ is chosen pseudo-randomly to avoid cross-connection packet collisions and mitigate spoofing attacks (A is TRUE). Because a SYN packet consumes exactly 1 sequence number in the sequence space, server $Q$ acknowledges it with $N_Q = N_P + 1$ (D is TRUE).}

\end{tcolorbox}
\vspace{6pt}\hrule height 0.35pt\vspace{6pt}
\noindent
\begin{tcolorbox}[
    blanker,
    sidebyside,
    sidebyside align=top,
    lefthand width=0.55\textwidth,
    top=4pt, bottom=4pt, left=0pt, right=0pt,
    middle=16pt,
    underlay={
        \draw[dashed, black!80, line width=0.8pt, dash pattern=on 4pt off 3pt]
            ([xshift=0.55\textwidth+8pt]frame.north west) --
            ([xshift=0.55\textwidth+8pt]frame.south west);
    }
]
\textbf{Q.156} \quad {\footnotesize\textsc{[MSQ $\cdot$ 1 Mark(s) $\cdot$ No Neg]}} \hfill {\footnotesize\textbf{GATE 2024 CS1} \quad \tiny\texttt{2024\_CS1\_CS\_Q31}}
\vspace{2pt}\hrule height 0.3pt\vspace{4pt}

Which of the following fields is/are modified in the IP header of a packet going out of a network address translation (NAT) device from an internal network to an external network?

\par\vspace{6pt}\noindent
\begin{tabular}{@{}p{0.48\linewidth}@{\hspace{0.04\linewidth}}p{0.48\linewidth}@{}}
$\square$\ \textbf{(A)} Source IP & $\square$\ \textbf{(B)} Destination IP \\
\addlinespace[3pt]
$\square$\ \textbf{(C)} Header Checksum & $\square$\ \textbf{(D)} Total Length
\end{tabular}

\tcblower

\textbf{Official Key:} \textbf{\large A;C} \hfill \textit{\small [Easy]}
\vspace{2pt}\hrule height 0.3pt\vspace{3pt}
\textbf{Topic:} Computer Networks \\
\textbf{Subtopic:} \textit{Network Address Translation (NAT)}
\vspace{2pt}\hrule height 0.3pt\vspace{3pt}
\textbf{Guide / Hint:} \\
{\small When a packet leaves a private subnet through an outbound NAT gateway, the router rewrites the private Source IP address to its public external IP address (A is modified). Because the IP header fields have changed, the IPv4 Header Checksum must be recalculated and replaced (C is modified). Destination IP and Total Length are unmodified.}

\end{tcolorbox}
\vspace{6pt}\hrule height 0.35pt\vspace{6pt}
\noindent
\begin{tcolorbox}[
    blanker,
    sidebyside,
    sidebyside align=top,
    lefthand width=0.55\textwidth,
    top=4pt, bottom=4pt, left=0pt, right=0pt,
    middle=16pt,
    underlay={
        \draw[dashed, black!80, line width=0.8pt, dash pattern=on 4pt off 3pt]
            ([xshift=0.55\textwidth+8pt]frame.north west) --
            ([xshift=0.55\textwidth+8pt]frame.south west);
    }
]
\textbf{Q.157} \quad {\footnotesize\textsc{[MSQ $\cdot$ 1 Mark(s) $\cdot$ No Neg]}} \hfill {\footnotesize\textbf{GATE 2024 CS2} \quad \tiny\texttt{2024\_CS2\_CS\_Q23}}
\vspace{2pt}\hrule height 0.3pt\vspace{4pt}

Node X has a TCP connection open to node Y. The packets from X to Y go through an intermediate IP router R. Ethernet switch S is the first switch on the network path between X and R. Consider a packet sent from X to Y over this connection.
Which of the following statements is/are TRUE about the destination addresses?

\par\vspace{6pt}\noindent
\begin{tabular}{@{}p{0.98\linewidth}@{}}
$\square$\ \textbf{(A)} The destination IP address is the IP address of R \\
\addlinespace[3pt]
$\square$\ \textbf{(B)} The destination IP address is the IP address of Y \\
\addlinespace[3pt]
$\square$\ \textbf{(C)} The destination MAC address is the MAC address of S \\
\addlinespace[3pt]
$\square$\ \textbf{(D)} The destination MAC address is the MAC address of Y
\end{tabular}

\tcblower

\textbf{Official Key:} \textbf{\large B} \hfill \textit{\small [Easy]}
\vspace{2pt}\hrule height 0.3pt\vspace{3pt}
\textbf{Topic:} Computer Networks \\
\textbf{Subtopic:} \textit{Concept of layering: OSI and TCP/IP Protocol Stacks}
\vspace{2pt}\hrule height 0.3pt\vspace{3pt}
\textbf{Guide / Hint:} \\
{\small In an end-to-end IP network, the destination IP address remains invariant throughout transit and is always the IP address of final host Y. The destination MAC address changes hop-by-hop and is set to the MAC address of the next-hop router R (not switch S, which operates transparently at Layer 2).}

\end{tcolorbox}
\vspace{6pt}\hrule height 0.35pt\vspace{6pt}
\noindent
\begin{tcolorbox}[
    blanker,
    sidebyside,
    sidebyside align=top,
    lefthand width=0.55\textwidth,
    top=4pt, bottom=4pt, left=0pt, right=0pt,
    middle=16pt,
    underlay={
        \draw[dashed, black!80, line width=0.8pt, dash pattern=on 4pt off 3pt]
            ([xshift=0.55\textwidth+8pt]frame.north west) --
            ([xshift=0.55\textwidth+8pt]frame.south west);
    }
]
\textbf{Q.158} \quad {\footnotesize\textsc{[MSQ $\cdot$ 1 Mark(s) $\cdot$ No Neg]}} \hfill {\footnotesize\textbf{GATE 2024 CS2} \quad \tiny\texttt{2024\_CS2\_CS\_Q28}}
\vspace{2pt}\hrule height 0.3pt\vspace{4pt}

Which of the following statements about IPv4 fragmentation is/are TRUE?

\par\vspace{6pt}\noindent
\begin{tabular}{@{}p{0.98\linewidth}@{}}
$\square$\ \textbf{(A)} The fragmentation of an IP datagram is performed only at the source of the datagram \\
\addlinespace[3pt]
$\square$\ \textbf{(B)} The fragmentation of an IP datagram is performed at any IP router which finds that the size of the datagram is greater than the MTU of the outgoing link \\
\addlinespace[3pt]
$\square$\ \textbf{(C)} The reassembly of fragments is performed only at the destination of the datagram \\
\addlinespace[3pt]
$\square$\ \textbf{(D)} The reassembly of fragments is performed at all intermediate routers along the path from the source to destination
\end{tabular}

\tcblower

\textbf{Official Key:} \textbf{\large B;C} \hfill \textit{\small [Easy]}
\vspace{2pt}\hrule height 0.3pt\vspace{3pt}
\textbf{Topic:} Computer Networks \\
\textbf{Subtopic:} \textit{Fragmentation and IP addressing, IPv4, CIDR notation}
\vspace{2pt}\hrule height 0.3pt\vspace{3pt}
\textbf{Guide / Hint:} \\
{\small In IPv4, any intermediate router encountering a packet larger than the outgoing link MTU performs fragmentation (B). In contrast, fragment reassembly is performed strictly at the final destination host (C) to prevent re-fragmentation overhead and handle out-of-order routing.}

\end{tcolorbox}
\vspace{6pt}\hrule height 0.35pt\vspace{6pt}
\noindent
\begin{tcolorbox}[
    blanker,
    sidebyside,
    sidebyside align=top,
    lefthand width=0.55\textwidth,
    top=4pt, bottom=4pt, left=0pt, right=0pt,
    middle=16pt,
    underlay={
        \draw[dashed, black!80, line width=0.8pt, dash pattern=on 4pt off 3pt]
            ([xshift=0.55\textwidth+8pt]frame.north west) --
            ([xshift=0.55\textwidth+8pt]frame.south west);
    }
]
\textbf{Q.159} \quad {\footnotesize\textsc{[MSQ $\cdot$ 1 Mark(s) $\cdot$ No Neg]}} \hfill {\footnotesize\textbf{GATE 2024 CS2} \quad \tiny\texttt{2024\_CS2\_CS\_Q32}}
\vspace{2pt}\hrule height 0.3pt\vspace{4pt}

Which of the following fields of an IP header is/are always modified by any router before it forwards the IP packet?

\par\vspace{6pt}\noindent
\begin{tabular}{@{}p{0.48\linewidth}@{\hspace{0.04\linewidth}}p{0.48\linewidth}@{}}
$\square$\ \textbf{(A)} Source IP Address & $\square$\ \textbf{(B)} Protocol \\
\addlinespace[3pt]
$\square$\ \textbf{(C)} Time to Live (TTL) & $\square$\ \textbf{(D)} Header Checksum
\end{tabular}

\tcblower

\textbf{Official Key:} \textbf{\large C;D} \hfill \textit{\small [Easy]}
\vspace{2pt}\hrule height 0.3pt\vspace{3pt}
\textbf{Topic:} Computer Networks \\
\textbf{Subtopic:} \textit{Fragmentation and IP addressing, IPv4, CIDR notation}
\vspace{2pt}\hrule height 0.3pt\vspace{3pt}
\textbf{Guide / Hint:} \\
{\small Whenever an IP packet is processed by an intermediate router, its Time to Live (TTL) field is decremented by at least 1. Because the header content changes, the router must recompute and update the Header Checksum field. Source and Protocol remain unchanged.}

\end{tcolorbox}
\vspace{6pt}\hrule height 0.35pt\vspace{6pt}
\noindent
\begin{tcolorbox}[
    blanker,
    sidebyside,
    sidebyside align=top,
    lefthand width=0.55\textwidth,
    top=4pt, bottom=4pt, left=0pt, right=0pt,
    middle=16pt,
    underlay={
        \draw[dashed, black!80, line width=0.8pt, dash pattern=on 4pt off 3pt]
            ([xshift=0.55\textwidth+8pt]frame.north west) --
            ([xshift=0.55\textwidth+8pt]frame.south west);
    }
]
\textbf{Q.160} \quad {\footnotesize\textsc{[MCQ $\cdot$ 2 Mark(s) $\cdot$ Neg: 0.66]}} \hfill {\footnotesize\textbf{GATE 2024 CS2} \quad \tiny\texttt{2024\_CS2\_CS\_Q38}}
\vspace{2pt}\hrule height 0.3pt\vspace{4pt}

Which one of the following CIDR prefixes exactly represents the range of IP addresses 10.12.2.0 to 10.12.3.255?

\par\vspace{6pt}\noindent
\begin{tabular}{@{}p{0.48\linewidth}@{\hspace{0.04\linewidth}}p{0.48\linewidth}@{}}
$\square$\ \textbf{(A)} 10.12.2.0/23 & $\square$\ \textbf{(B)} 10.12.2.0/24 \\
\addlinespace[3pt]
$\square$\ \textbf{(C)} 10.12.0.0/22 & $\square$\ \textbf{(D)} 10.12.2.0/22
\end{tabular}

\tcblower

\textbf{Official Key:} \textbf{\large A} \hfill \textit{\small [Easy]}
\vspace{2pt}\hrule height 0.3pt\vspace{3pt}
\textbf{Topic:} Computer Networks \\
\textbf{Subtopic:} \textit{Fragmentation and IP addressing, IPv4, CIDR notation}
\vspace{2pt}\hrule height 0.3pt\vspace{3pt}
\textbf{Guide / Hint:} \\
{\small The address range spans from $10.12.2.0$ ($00000010_2$) to $10.12.3.255$ ($00000011_2$). The first 23 bits ($8 + 8 + 7 = 23$) are identical ($10.12.2$). Total addresses in range $= 2 \times 256 = 512 = 2^{32 - 23}$, which exactly matches CIDR prefix `10.12.2.0/23`.}

\end{tcolorbox}
\vspace{6pt}\hrule height 0.35pt\vspace{6pt}
\noindent
\begin{tcolorbox}[
    blanker,
    sidebyside,
    sidebyside align=top,
    lefthand width=0.55\textwidth,
    top=4pt, bottom=4pt, left=0pt, right=0pt,
    middle=16pt,
    underlay={
        \draw[dashed, black!80, line width=0.8pt, dash pattern=on 4pt off 3pt]
            ([xshift=0.55\textwidth+8pt]frame.north west) --
            ([xshift=0.55\textwidth+8pt]frame.south west);
    }
]
\textbf{Q.161} \quad {\footnotesize\textsc{[NAT $\cdot$ 2 Mark(s) $\cdot$ No Neg]}} \hfill {\footnotesize\textbf{GATE 2024 CS2} \quad \tiny\texttt{2024\_CS2\_CS\_Q55}}
\vspace{2pt}\hrule height 0.3pt\vspace{4pt}

Consider an Ethernet segment with a transmission speed of $10^8$ bits/sec and a maximum segment length of 500 meters. If the speed of propagation of the signal in the medium is $2 \times 10^8$ meters/sec, then the minimum frame size (in bits) required for collision detection is \_\_\_\_\_\_\. (Answer in integer)

\par\vspace{6pt}\noindent\textbf{Your Answer:} \fbox{\rule{0pt}{1.3ex}\hspace{3cm}}

\tcblower

\textbf{Official Key:} \textbf{\large 500 to 500} \hfill \textit{\small [Easy]}
\vspace{2pt}\hrule height 0.3pt\vspace{3pt}
\textbf{Topic:} Computer Networks \\
\textbf{Subtopic:} \textit{Data link layer: framing, error detection, Medium Access Control, Ethernet bridging}
\vspace{2pt}\hrule height 0.3pt\vspace{3pt}
\textbf{Guide / Hint:} \\
{\small Propagation delay $T_p = \frac{\text{Distance}}{\text{Velocity}} = \frac{500}{2 \times 10^8} = 2.5 \times 10^{-6}\text{ s} = 2.5\,\mu\text{s}$. For collision detection in CSMA/CD, transmission time must satisfy $T_t \ge 2 T_p = 5\,\mu\text{s}$. Minimum frame size $L = 2 T_p \times B = (5 \times 10^{-6}\text{ s}) \times 10^8\text{ bps} = 500$ bits.}

\end{tcolorbox}
\vspace{6pt}\hrule height 0.35pt\vspace{6pt}
\noindent
\begin{tcolorbox}[
    blanker,
    sidebyside,
    sidebyside align=top,
    lefthand width=0.55\textwidth,
    top=4pt, bottom=4pt, left=0pt, right=0pt,
    middle=16pt,
    underlay={
        \draw[dashed, black!80, line width=0.8pt, dash pattern=on 4pt off 3pt]
            ([xshift=0.55\textwidth+8pt]frame.north west) --
            ([xshift=0.55\textwidth+8pt]frame.south west);
    }
]
\textbf{Q.162} \quad {\footnotesize\textsc{[MCQ $\cdot$ 1 Mark(s) $\cdot$ Neg: 0.33]}} \hfill {\footnotesize\textbf{GATE 2025 CS1} \quad \tiny\texttt{2025\_CS1\_CS\_Q16}}
\vspace{2pt}\hrule height 0.3pt\vspace{4pt}

Identify the ONE CORRECT matching between the OSI layers and their corresponding functionalities as shown.
\par\medskip
\begin{tabular}{ll}
\textbf{OSI Layers} & \textbf{Functionalities} \\
(a) Network layer & (I) Packet routing \\
(b) Transport layer & (II) Framing and error handling \\
(c) Datalink layer & (III) Host to host communication
\end{tabular}

\par\vspace{6pt}\noindent
\begin{tabular}{@{}p{0.48\linewidth}@{\hspace{0.04\linewidth}}p{0.48\linewidth}@{}}
$\square$\ \textbf{(A)} (a)-(I), (b)-(II), (c)-(III) & $\square$\ \textbf{(B)} (a)-(I), (b)-(III), (c)-(II) \\
\addlinespace[3pt]
$\square$\ \textbf{(C)} (a)-(II), (b)-(I), (c)-(III) & $\square$\ \textbf{(D)} (a)-(III), (b)-(II), (c)-(I)
\end{tabular}

\tcblower

\textbf{Official Key:} \textbf{\large B} \hfill \textit{\small [Easy]}
\vspace{2pt}\hrule height 0.3pt\vspace{3pt}
\textbf{Topic:} Computer Networks \\
\textbf{Subtopic:} \textit{Concept of layering: OSI and TCP/IP Protocol Stacks}
\vspace{2pt}\hrule height 0.3pt\vspace{3pt}
\textbf{Guide / Hint:} \\
{\small The Network layer is responsible for packet routing across networks ((a)-(I)). The Transport layer provides end-to-end communication services between hosts ((b)-(III)). The Data Link layer handles framing, flow control, and physical error handling ((c)-(II)).}

\end{tcolorbox}
\vspace{6pt}\hrule height 0.35pt\vspace{6pt}
\noindent
\begin{tcolorbox}[
    blanker,
    sidebyside,
    sidebyside align=top,
    lefthand width=0.55\textwidth,
    top=4pt, bottom=4pt, left=0pt, right=0pt,
    middle=16pt,
    underlay={
        \draw[dashed, black!80, line width=0.8pt, dash pattern=on 4pt off 3pt]
            ([xshift=0.55\textwidth+8pt]frame.north west) --
            ([xshift=0.55\textwidth+8pt]frame.south west);
    }
]
\textbf{Q.163} \quad {\footnotesize\textsc{[MSQ $\cdot$ 1 Mark(s) $\cdot$ No Neg]}} \hfill {\footnotesize\textbf{GATE 2025 CS1} \quad \tiny\texttt{2025\_CS1\_CS\_Q22}}
\vspace{2pt}\hrule height 0.3pt\vspace{4pt}

Consider the 3-way handshaking protocol for TCP connection establishment. Let the three packets exchanged during the connection establishment be denoted as P1, P2, and P3, in order. Which of the following option(s) is/are TRUE with respect to TCP header flags that are set in the packets?

\par\vspace{6pt}\noindent
\begin{tabular}{@{}p{0.48\linewidth}@{\hspace{0.04\linewidth}}p{0.48\linewidth}@{}}
$\square$\ \textbf{(A)} P3: SYN = 1, ACK = 1 & $\square$\ \textbf{(B)} P2: SYN = 1, ACK = 1 \\
\addlinespace[3pt]
$\square$\ \textbf{(C)} P2: SYN = 0, ACK = 1 & $\square$\ \textbf{(D)} P1: SYN = 1
\end{tabular}

\tcblower

\textbf{Official Key:} \textbf{\large B;D} \hfill \textit{\small [Easy]}
\vspace{2pt}\hrule height 0.3pt\vspace{3pt}
\textbf{Topic:} Computer Networks \\
\textbf{Subtopic:} \textit{Transport layer: flow control and congestion control, UDP, TCP, sockets}
\vspace{2pt}\hrule height 0.3pt\vspace{3pt}
\textbf{Guide / Hint:} \\
{\small In TCP 3-way handshaking: Packet 1 (Client to Server) sends a SYN segment with flag $[\text{SYN}=1, \text{ACK}=0]$; Packet 2 (Server to Client) responds with $[\text{SYN}=1, \text{ACK}=1]$; Packet 3 (Client to Server) completes setup with $[\text{SYN}=0, \text{ACK}=1]$. Hence, P1 has SYN = 1 and P2 has SYN = 1, ACK = 1.}

\end{tcolorbox}
\vspace{6pt}\hrule height 0.35pt\vspace{6pt}
\noindent
\begin{tcolorbox}[
    blanker,
    sidebyside,
    sidebyside align=top,
    lefthand width=0.55\textwidth,
    top=4pt, bottom=4pt, left=0pt, right=0pt,
    middle=16pt,
    underlay={
        \draw[dashed, black!80, line width=0.8pt, dash pattern=on 4pt off 3pt]
            ([xshift=0.55\textwidth+8pt]frame.north west) --
            ([xshift=0.55\textwidth+8pt]frame.south west);
    }
]
\textbf{Q.164} \quad {\footnotesize\textsc{[NAT $\cdot$ 2 Mark(s) $\cdot$ No Neg]}} \hfill {\footnotesize\textbf{GATE 2025 CS1} \quad \tiny\texttt{2025\_CS1\_CS\_Q56}}
\vspace{2pt}\hrule height 0.3pt\vspace{4pt}

Suppose a 5-bit message is transmitted from a source to a destination through a noisy channel. The probability that a bit of the message gets flipped during transmission is 0.01. Flipping of each bit is independent of one another. The probability that the message is delivered error-free to the destination is \_\_\_\_\_\_\. (rounded off to three decimal places)

\par\vspace{6pt}\noindent\textbf{Your Answer:} \fbox{\rule{0pt}{1.3ex}\hspace{3cm}}

\tcblower

\textbf{Official Key:} \textbf{\large 0.949 to 0.952} \hfill \textit{\small [Easy]}
\vspace{2pt}\hrule height 0.3pt\vspace{3pt}
\textbf{Topic:} Computer Networks \\
\textbf{Subtopic:} \textit{Data link layer: framing, error detection, Medium Access Control, Ethernet bridging}
\vspace{2pt}\hrule height 0.3pt\vspace{3pt}
\textbf{Guide / Hint:} \\
{\small For each bit, the probability of correct transmission without flipping is $1 - 0.01 = 0.99$. Since all 5 bit transmissions are mutually independent, the probability that all 5 bits are transmitted error-free is $(0.99)^5 = 0.95099 \approx 0.951$.}

\end{tcolorbox}
\vspace{6pt}\hrule height 0.35pt\vspace{6pt}
\noindent
\begin{tcolorbox}[
    blanker,
    sidebyside,
    sidebyside align=top,
    lefthand width=0.55\textwidth,
    top=4pt, bottom=4pt, left=0pt, right=0pt,
    middle=16pt,
    underlay={
        \draw[dashed, black!80, line width=0.8pt, dash pattern=on 4pt off 3pt]
            ([xshift=0.55\textwidth+8pt]frame.north west) --
            ([xshift=0.55\textwidth+8pt]frame.south west);
    }
]
\textbf{Q.165} \quad {\footnotesize\textsc{[MCQ $\cdot$ 1 Mark(s) $\cdot$ Neg: 0.33]}} \hfill {\footnotesize\textbf{GATE 2025 CS2} \quad \tiny\texttt{2025\_CS2\_CS\_Q16}}
\vspace{2pt}\hrule height 0.3pt\vspace{4pt}

Consider the following statements:
(i) Address Resolution Protocol (ARP) provides a mapping from an IP address to the corresponding hardware (link-layer) address.
(ii) A single TCP segment from a sender S to a receiver R cannot carry both data from S to R and acknowledgement for a segment from R to S.
Which ONE of the following is CORRECT?

\par\vspace{6pt}\noindent
\begin{tabular}{@{}p{0.48\linewidth}@{\hspace{0.04\linewidth}}p{0.48\linewidth}@{}}
$\square$\ \textbf{(A)} Both (i) and (ii) are TRUE & $\square$\ \textbf{(B)} (i) is TRUE and (ii) is FALSE \\
\addlinespace[3pt]
$\square$\ \textbf{(C)} (i) is FALSE and (ii) is TRUE & $\square$\ \textbf{(D)} Both (i) and (ii) are FALSE
\end{tabular}

\tcblower

\textbf{Official Key:} \textbf{\large B} \hfill \textit{\small [Easy]}
\vspace{2pt}\hrule height 0.3pt\vspace{3pt}
\textbf{Topic:} Computer Networks \\
\textbf{Subtopic:} \textit{Basics of IP support protocols: ARP and TCP headers}
\vspace{2pt}\hrule height 0.3pt\vspace{3pt}
\textbf{Guide / Hint:} \\
{\small ARP resolves an IPv4 address to a physical MAC (link-layer) address, so (i) is true. A TCP segment supports piggybacking, where data and acknowledgement flags/numbers are sent in the same segment, so (ii) is false.}

\end{tcolorbox}
\vspace{6pt}\hrule height 0.35pt\vspace{6pt}
\noindent
\begin{tcolorbox}[
    blanker,
    sidebyside,
    sidebyside align=top,
    lefthand width=0.55\textwidth,
    top=4pt, bottom=4pt, left=0pt, right=0pt,
    middle=16pt,
    underlay={
        \draw[dashed, black!80, line width=0.8pt, dash pattern=on 4pt off 3pt]
            ([xshift=0.55\textwidth+8pt]frame.north west) --
            ([xshift=0.55\textwidth+8pt]frame.south west);
    }
]
\textbf{Q.166} \quad {\footnotesize\textsc{[MCQ $\cdot$ 1 Mark(s) $\cdot$ Neg: 0.33]}} \hfill {\footnotesize\textbf{GATE 2025 CS2} \quad \tiny\texttt{2025\_CS2\_CS\_Q17}}
\vspace{2pt}\hrule height 0.3pt\vspace{4pt}

Consider the routing protocols given in List I and the names given in List II:
\par\medskip
\begin{center}\small
\begin{tabular}{ll}
\toprule
List I & List II \\
\midrule
(i) Distance vector routing & (a) Bellman-Ford \\
(ii) Link state routing & (b) Dijkstra \\
\bottomrule
\end{tabular}\end{center}
\par\medskip
For matching of items in List I with those in List II, which ONE of the following options is CORRECT?

\par\vspace{6pt}\noindent
\begin{tabular}{@{}p{0.48\linewidth}@{\hspace{0.04\linewidth}}p{0.48\linewidth}@{}}
$\square$\ \textbf{(A)} (i) – (a) and (ii) – (b) & $\square$\ \textbf{(B)} (i) – (a) and (ii) – (a) \\
\addlinespace[3pt]
$\square$\ \textbf{(C)} (i) – (b) and (ii) – (a) & $\square$\ \textbf{(D)} (i) – (b) and (ii) – (b)
\end{tabular}

\tcblower

\textbf{Official Key:} \textbf{\large A} \hfill \textit{\small [Easy]}
\vspace{2pt}\hrule height 0.3pt\vspace{3pt}
\textbf{Topic:} Computer Networks \\
\textbf{Subtopic:} \textit{Routing protocols: distance vector and link state}
\vspace{2pt}\hrule height 0.3pt\vspace{3pt}
\textbf{Guide / Hint:} \\
{\small Distance Vector routing is based on the distributed Bellman-Ford algorithm, whereas Link State routing uses Dijkstra's shortest path algorithm.}

\end{tcolorbox}
\vspace{6pt}\hrule height 0.35pt\vspace{6pt}
\end{document}